\chapter{Modular forms and the analytic toolkit}\label{ch:L14}
\begingroup\setlength{\emergencystretch}{3em}

\providecommand{\lfourteenH}{\mathbb{H}}
\providecommand{\lfourteenSL}{\mathrm{SL}_2(\mathbb{Z})}
\providecommand{\lfourteenGt}{\Gamma_\theta}
\providecommand{\lfourteenE}{\mathcal{E}_{46}}
\providecommand{\lfourteenGS}{\mathcal{G}}
\providecommand{\lfourteenev}{\mathrm{ev}}
\providecommand{\lfourteenIm}{\operatorname{Im}}
\providecommand{\lfourteenRe}{\operatorname{Re}}

This chapter documents the release of the LeanVOA library named in the front matter; it was first written against an earlier state of the library, and checks that were
run at an earlier state of the library and not re-run for this edition are marked as such. Every Lean quotation below is taken from the files of the release and is machine-checked by \texttt{scripts/textbook/sync.py} to
occur verbatim (up to whitespace and comments) in the source. The chapter covers the files of
\leanpath{VOA/\allowbreak Modularity/} assigned to it in \texttt{Notes/Textbook/assignment.txt}: the Eisenstein series and the
level-one structure theorem, the bridge between formal $q$-series and holomorphic functions on the upper half-plane, the
Serre derivative, Lambert series, the growth (``subexponential'') estimates and the majorant method for formal modular
differential equations, the cusp $q\to0$, the slash action restricted to subgroups and the theta group $\lfourteenGt$, the
read-off of the modular $S$-matrix, Jacobi characters, the analytic $j$-invariant, and the abstract ``Moonshine corollary''
of Zhu's theorem. About thirty-five files and ten thousand lines of Lean are involved.

\paragraph{How this chapter relates to its neighbours.} Chapter~\ref{ch:05} already develops, for the Moonshine
construction, the classical theory of $\lfourteenSL$, the slash operator, the Eisenstein series $E_4,E_6$, the
discriminant $\Delta$, the functions $j$ and $J$, theta series and the formal series \lean{E4Series},
\lean{E6Series}, \lean{deltaSeries}, \lean{jShift}, \lean{JShift} of \leanpath{VOA/\allowbreak Modularity/\allowbreak JFunction.lean} and
\lean{lambertCoeff} of \leanpath{VOA/\allowbreak Modularity/\allowbreak LambertSeries.lean}. We do not repeat that material: we
recall it in Section~\ref{l14:sec:recall} with pointers and then cover what is \emph{general} or \emph{different}. The
reader should know three other chapters. Chapter~\ref{ch:01} is the book's Lean toolkit (structures, \lean{∑ᶠ},
\lean{tsum}, coercions, \lean{Classical.choose}); we point to it instead of re-explaining Lean idioms.
Chapter~\ref{ch:L12} defines modules, their formal characters \lean{gradedDim} and one-point functions \lean{onePoint}
(Definition~\ref{l12:def:char}), conformal weights (Definition~\ref{l12:def:lowest}), the $T$-periodicity of characters
(Theorem~\ref{l12:thm:Tperiod}), Zhu's algebra $A(V)$ (Section~\ref{l12:sec:zhu}) and $C_2$-cofiniteness
(Definition~\ref{l12:def:c2}). Chapter~\ref{ch:L15} proves Zhu's modularity theorem --- for the one-point functions of a
$C_2$-cofinite, rational $V$ (rational in the sense of the library's \lean{AdmissiblyRational}), through a more general form
for $C_2$-cofinite $V$ whose Zhu algebra is \emph{semisimple} --- and is the main
\emph{consumer} of this chapter: almost every analytic statement here was written to serve one step of Zhu's argument, and we say which.
Where a statement of this chapter needs an object that is formally owned by Chapter~\ref{ch:L15} (the Zhu-normalised
Eisenstein coefficients \lean{zhuG}, the square-bracket modes \lean{vSqZ}, the trace functions \lean{traceFun}) we quote its
definition here, with its file, so that the chapter can be read on its own.

\emph{Release revision.} The chapter was revised for the first public release, applying the release audit \texttt{Notes/\allowbreak Textbook/\allowbreak referee/\allowbreak
RELEASE\_AUDIT\_ZHU.md} (2026-10-27). No quotation of this chapter changed; the edits concern how the consumer theorem of
Chapter~\ref{ch:L15} is described.

\emph{Second release revision (2026-10-28).} At this revision the description of
that consumer theorem was corrected again, following finding F2 of \texttt{Notes/\allowbreak Textbook/\allowbreak
referee/\allowbreak RELEASE\_AUDIT\_NEW\_DECLS.md}. The first release revision said that the library ``does not derive it from
rationality''. That was wrong. The library proves ``rational $\Rightarrow A(V)$ semisimple''
(\lean{isSemisimpleRing_}\allowbreak\lean{A_of_}\allowbreak\lean{admissiblyRational}, Chapter~\ref{ch:L15},
Theorem~\ref{l15:thm:zhuss}), and Zhu's theorem is stated for $C_2$-cofinite rational $V$
(Theorem~\ref{l15:thm:rational}) and proved for every lattice VOA $V_{\sqrt{2m}\Z}$. No quotation of this chapter changed.

\paragraph{The findings of this chapter, in brief.} Every headline theorem of this chapter is proved in Lean and
\texttt{\#print axioms}-clean (only \lean{propext}, \lean{Classical.choice}, \lean{Quot.sound}); none uses \lean{VOA.pbw} (the PBW theorem of Chapter~\ref{ch:L02}; formerly a project axiom, now a theorem). The analytic layer is faithful and, in several places, genuinely new relative to Mathlib (the
equivariance of the Serre derivative, holomorphy and term-by-term differentiation of evaluated $q$-series, the majorant
theorem). The fidelity checks did, however, turn up the following, each explained in a red box where it occurs.
\begin{itemize}
\item The structure theorem \lean{levelOne_coe_mem_E46} (``$M_* = \C[E_4,E_6]$'') proves only that every modular form lies
  in the \emph{algebra generated} by $E_4$ and $E_6$. Algebraic independence and homogeneity are not part of it. They are
  proved, in full, elsewhere in the library (the vendored \lean{ModularForm.levelOne_eq_polynomial} of
  Chapter~\ref{ch:02}); the two developments are not connected, and a predicate documented as ``missing''
  (\lean{E46HasWeightedPolynomialReps}) is in fact a consequence of the vendored theorem. A separate Lean check
  (Theorem~\ref{l14:thm:certreps}) proved it first; the argument is now a library theorem,
  \lean{e46HasWeightedPolynomialReps}.
\item The Ramanujan/Serre closure \lean{E46SerreClosed} ($\partial_k M_k\subseteq M_{k+2}$) is a library theorem,
  \lean{e46SerreClosed} (\leanpath{VOA/\allowbreak Modularity/\allowbreak EisensteinSerreClosure.lean}), proved through
  \lean{serreForm}: the Serre derivative of a level-one form is a level-one form. It is consumed by Chapter~\ref{ch:L16}.
\item The $\Gamma_\theta$ ``conjectures'' of \leanpath{GammaThetaStatement.lean} carry a free proposition slot
  \lean{hFam} through \lean{CompleteRationalFamily}, which is the identity on propositions, and the character families in
  \lean{SuperCharacterSystem} and \lean{SuperRealisesLawOn} are not tied to the super VOA at all. They are statements, not
  theorems, and are honest about it; but as written they are templates whose content depends entirely on what is plugged in
  (checked in Lean, Proposition~\ref{l14:prop:certtheta}).
\item The ``Hauptmodul step'' \lean{char_eq_jInvariant_add_const} assumes a weight-zero modular form equal to $\chi - j$;
  this hypothesis carries the whole content. Nothing in this chapter is instantiated at $\Vnat$: the library's theorem
  $\chi_{\Vnat} = J$ (Chapter~\ref{ch:05}) is a computation and does not pass through these files.
\item The chapter's main consumer, Zhu's modularity theorem, is stated in Chapter~\ref{ch:L15} under $C_2$-cofiniteness and
  rationality (\lean{AdmissiblyRational}), and in a more general form under $C_2$-cofiniteness and semisimplicity of $A(V)$, for
  one-point functions; the convergence of the $q$-series it evaluates is a \emph{separate} theorem
  (\lean{onePointSubexp_of_c2} together with this chapter's \lean{qSummableAt_of_subexponential}), not part of its conclusion.
\item Several docstrings are stale or wrong: one says that $p(n)\le 2^n$ gives subexponential growth (it does not); several
  call residuals open that are proved in other files; the Moonshine-corollary files say $\Vnat$ does not exist.
\end{itemize}

%======================================================================
\section{Conventions, and what is recalled from Chapter 5}\label{l14:sec:recall}
%======================================================================

\subsection{Three kinds of ``$q$-series''}\label{l14:sec:threekinds}

The single most important thing to keep track of in this chapter is \emph{which kind of object} a statement is about. The
library uses three, and each has its own job: \textbf{(1)} carries characters and one-point functions of modules (exponents
$h+n$ with $h\in\C$ arbitrary), and is the language of Zhu's recursion and of the formal differential equations of
Section~\ref{l14:sec:ode}; \textbf{(2)} carries the coefficient rings (Eisenstein series, $E_4$, $E_6$, $\Delta$, $j$) and the
coefficients of those equations; \textbf{(3)} is where modularity is a statement at all, since $\lfourteenSL$ acts on $\tau$, not on
coefficients. Objects of types (1) and (2) are formal: equalities are coefficientwise and no convergence is involved. Objects
of type (3) are honest functions, but whenever one is \emph{produced} from a type-(1) object by an infinite sum, Lean uses
\lean{tsum}, which returns the junk value $0$ on a divergent series (Chapter~\ref{ch:01}, Section~\ref{ch01:sec:junk}); so every
theorem that needs the honest value of such a sum carries a summability or growth hypothesis, and we point this out each
time.
\begin{enumerate}
\item \textbf{Formal $q$-series with complex exponents}, the type \lean{FormalQSeries}, which is literally \lean{ℂ → ℂ}
  (Chapter~\ref{ch:L12}, Definition~\ref{l12:def:char}). A function $\chi:\C\to\C$ is read as the formal sum
  $\sum_{e\in\C}\chi(e)\,q^e$. No convergence is implied, and the sum may have uncountable support; but every series that
  actually occurs is supported on a single progression $h+\N$, because characters of modules with a conformal weight
  are. The formal character of a module $M$ with central charge $c$ is
  \begin{equation*}
    \chi_M = \sum_{\alpha}(\dim M_\alpha)\,q^{\alpha - c/24},\qquad \text{i.e. } \mathtt{gradedDim}\,N\,(e) = \dim M_{e+c/24}.
  \end{equation*}
\item \textbf{Formal power series} \lean{PowerSeries ℂ} $=\C[[X]]$ (Mathlib), with natural-number exponents. This is the
  setting of Chapter~\ref{ch:05}'s $E_4$, $E_6$, $\Delta$, $j$, and of all coefficient rings of differential equations below.
  Notation: \lean{PowerSeries.coeff n F} is $[X^n]F$, \lean{PowerSeries.mk a} is $\sum_na_nX^n$; see Chapter~\ref{ch:05},
  Section~\ref{ch05:sec:formal}.
\item \textbf{Holomorphic functions} on the upper half-plane $\lfourteenH = \{\tau\in\C:\lfourteenIm\tau>0\}$. Mathlib's type is
  \lean{UpperHalfPlane} (notation \lean{ℍ}), whose elements are complex numbers bundled with a proof that the imaginary part
  is positive. For calculus the library prefers the \emph{open subset} \lean{imGT 0} $=\{z\in\C : 0<\lfourteenIm z\}$ of $\C$
  itself and functions $\C\to\C$, so that Mathlib's one-variable \lean{HasDerivAt}, \lean{deriv} and
  \lean{DifferentiableOn} apply directly; the dictionary between the two pictures is \lean{qEvalC_coe},
  \lean{mobiusC_coe}, \lean{slashC} versus Mathlib's \lean{∣[k]} (all below).
\end{enumerate}
The passage between (1)/(2) and (3) is \emph{evaluation}: $\chi\mapsto\sum_n\chi(h+n)q^{h+n}$ at $q = e^{2\pi i\tau}$. It is
only meaningful when the series converges, and a large part of this chapter proves, or reduces to named data, the
estimates that make it converge. The reverse passage, from a holomorphic $q$-periodic function to its coefficients, is
Mathlib's \lean{qExpansion}, and its injectivity (uniqueness of $q$-expansions) is Mathlib's
\lean{qExpansion_coeff_unique}.

\begin{warning}
\textbf{The \lean{tsum} convention.} An infinite sum in Lean, \lean{∑' n, f n} (\lean{tsum}), is defined for every $f$: if
$f$ is not summable it is $0$ by convention (Chapter~\ref{ch:01}). So a definition such as \lean{qEvalAt} below is
\emph{always} a complex number, and an equation between two \lean{tsum}s can be ``true'' because both sides are the junk
value $0$. Every theorem below that manipulates such sums either carries a summability hypothesis or proves summability
first; we check this in each fidelity box. Similarly \lean{∑ᶠ} (\lean{finsum}) is $0$ on infinite support.
\end{warning}

\subsection{The objects recalled from Chapter 5}

From Chapter~\ref{ch:05} we use: the upper half-plane and the action of $\lfourteenSL$ by $\gamma\tau = (a\tau+b)/(c\tau+d)$;
the generators $S = \left(\begin{smallmatrix}0&-1\\1&0\end{smallmatrix}\right)$, $T =
\left(\begin{smallmatrix}1&1\\0&1\end{smallmatrix}\right)$ (Mathlib \lean{ModularGroup.S}, \lean{ModularGroup.T}); the
slash operator $(f|_k\gamma)(\tau) = (c\tau+d)^{-k}f(\gamma\tau)$ and modular forms of level one
(Definition~\ref{ch05:def:modform}; Mathlib's type \lean{ModularForm 𝒮ℒ k}, where \lean{𝒮ℒ} is $\lfourteenSL$ viewed as a
subgroup of $\mathrm{GL}_2(\mathbb R)$); Mathlib's normalised Eisenstein series \lean{ModularForm.E hk} of weight
$k\ge3$, with \lean{ModularForm.E₄}, \lean{ModularForm.E₆}; the discriminant \lean{ModularForm.discriminant} and the
identity $\Delta = (E_4^3-E_6^2)/1728$ (Remark~\ref{ch05:rem:deltadef}); the Lambert coefficient
$\lambda_s(n) = \sum_{ab=n}a^s = \sigma_s(n)$ (\lean{lambertCoeff}, \lean{lambertCoeff_eq_sigma},
Definition~\ref{ch05:def:lambert}, Theorem~\ref{ch05:thm:lambert}); and the formal $E_4$, $E_6$, $\Delta$, $qj$, $qJ$
(Definitions~\ref{ch05:def:E46}--\ref{ch05:def:jshift}).

We also use from Mathlib, and explain where it first appears: the quasimodular Eisenstein series
$E_2 = 1 - 24\sum_{n\ge1}\sigma_1(n)q^n$ (\lean{EisensteinSeries.E2}) and its unnormalised version
$G_2 = 2\zeta(2)E_2 = \frac{\pi^2}{3}E_2$ (\lean{EisensteinSeries.G2}), the Bernoulli numbers \lean{bernoulli} with
$B_2 = \tfrac16$, $B_4 = -\tfrac1{30}$, $B_6 = \tfrac1{42}$, and Mathlib's limit at the cusp
\lean{valueAtInfty} and filter \lean{atImInfty} (``$\lfourteenIm\tau\to\infty$'').

\begin{physics}
Throughout, the physical dictionary is that of Chapter~\ref{ch:01}: a VOA is the chiral algebra of a two-dimensional
conformal field theory, a module is a sector, and the character
$\chi_M(\tau) = \tr_M q^{L_0-c/24}$, $q = e^{2\pi i\tau}$, is the \emph{torus partition function} of that sector: the
trace over states of the evolution operator $e^{2\pi i\tau(L_0-c/24)}$ around a torus of modular parameter $\tau$. The
shift $-c/24$ is the Casimir energy of the vacuum on a circle. Two parameters $\tau$ related by $\lfourteenSL$ describe the
same torus, so a consistent theory has partition functions that transform into each other under $\lfourteenSL$; this
\emph{modular invariance} is the physical origin of everything in this chapter. The generator $T:\tau\mapsto\tau+1$ is a
Dehn twist, and acts on a sector by the phase $e^{2\pi i(h-c/24)}$; the generator $S:\tau\mapsto-1/\tau$
\emph{exchanges the space and time circles of the torus}, and its matrix $\rho(S)$ is the deep piece of data (it determines
the fusion rules through Verlinde's formula).
\end{physics}

%======================================================================
\section{The slash action in the complex-plane picture}\label{l14:sec:slash}
%======================================================================

The slash operator $f\mapsto f|_k\gamma$ is the standard device of modular-form theory: a function $f$ is modular of weight $k$
exactly when $f|_k\gamma = f$ for all $\gamma$, and $|_k$ is a right action of $\lfourteenSL$ (Chapter~\ref{ch:05},
Section~\ref{ch05:sec:modular}). Nothing about it is new here except the domain.
Mathlib's slash operator \lean{∣[k]} acts on functions \lean{ℍ → ℂ}. To differentiate, the library re-implements it on
functions $\C\to\C$, restricted to statements on \lean{imGT 0}. We record the definitions.

\begin{definition}[Half-planes; automorphy factor; M\"obius action; slash, in $\C$]\label{l14:def:slashC}
This is the standard slash operator; we implement it on functions $\C\to\C$ rather than $\lfourteenH\to\C$ so that Mathlib's
one-variable derivative lemmas apply verbatim. For $c\in\mathbb R$ put $\lfourteenH_c = \{z\in\C : c<\lfourteenIm z\}$, so
$\lfourteenH_0 = \lfourteenH$. For
$\gamma = \left(\begin{smallmatrix}a&b\\c&d\end{smallmatrix}\right)\in\lfourteenSL$ and $z\in\C$ put
$j(\gamma,z) = cz+d$ and $\gamma z = (az+b)/(cz+d)$. For $k\in\Z$ and $f:\C\to\C$ put
$(f|_k\gamma)(z) = f(\gamma z)\,j(\gamma,z)^{-k}$.
\end{definition}
\noindent File \leanpath{VOA/\allowbreak Modularity/\allowbreak QEvalHolomorphic.lean}.
\begin{leancode}
def imGT (c : ℝ) : Set ℂ := {z : ℂ | c < z.im}
\end{leancode}
\noindent File \leanpath{VOA/\allowbreak Modularity/\allowbreak SerreDerivative.lean}.
\begin{leancode}
noncomputable def denomC (γ : SL(2, ℤ)) (z : ℂ) : ℂ := (γ 1 0 : ℂ) * z + (γ 1 1 : ℂ)

noncomputable def mobiusC (γ : SL(2, ℤ)) (z : ℂ) : ℂ :=
  ((γ 0 0 : ℂ) * z + (γ 0 1 : ℂ)) / denomC γ z

noncomputable def slashC (k : ℤ) (γ : SL(2, ℤ)) (f : ℂ → ℂ) (z : ℂ) : ℂ :=
  f (mobiusC γ z) * (denomC γ z) ^ (-k)
\end{leancode}
\begin{dictionary}
\lean{γ i j} is the $(i,j)$ matrix entry with indices in \lean{Fin 2} $=\{0,1\}$, so \lean{γ 1 0} $=c$, \lean{γ 1 1} $=d$,
\lean{γ 0 0} $=a$, \lean{γ 0 1} $=b$; the casts \lean{(· : ℂ)} send integers to complex numbers. \lean{SL(2, ℤ)} is
Mathlib's notation for $\lfourteenSL$. The exponent \lean{(-k)} is an integer power (\lean{zpow}), so the formula is the
textbook $(cz+d)^{-k}f(\gamma z)$. Off $\lfourteenH$ the functions are junk (for instance \lean{mobiusC} divides by zero at
$z=-d/c$, where Lean's division returns $0$); every theorem below is stated for $z\in$ \lean{imGT 0}, where
\lean{denomC_ne_zero} and \lean{mobiusC_mem_imGT} (both proved) guarantee $cz+d\neq0$ and $\gamma z\in\lfourteenH$. The
agreement with Mathlib is \lean{mobiusC_coe}: \lean{mobiusC γ (τ : ℂ) = ((γ • τ : UpperHalfPlane) : ℂ)}. These are
definitions; there is nothing to be vacuous.
\end{dictionary}
\begin{physics}
$j(\gamma,\tau)^k$ is the factor by which a field of conformal weight $k$ (more precisely, a one-point function of a state
of weight $k$) rescales when the torus is reparametrised; a weight-$k$ object is a $k$-form $f(\tau)(d\tau)^{k/2}$ on the
moduli space, and $d(\gamma\tau) = (c\tau+d)^{-2}d\tau$. Partition functions are weight $0$; one-point functions of weight-$k$
states are weight $k$.
\end{physics}
\begin{status}
Definitions. Supporting lemmas \lean{denomC_ne_zero}, \lean{mobiusC_mem_imGT}, \lean{hasDerivAt_mobiusC}
($\frac{d}{dz}\gamma z = (cz+d)^{-2}$, using $ad-bc=1$) and \lean{slashC_mul}
($f|_{k}\gamma\cdot g|_l\gamma = (fg)|_{k+l}\gamma$) are proved without hypotheses.
\end{status}

\begin{theorem}[The derivative does not commute with slashing]\label{l14:thm:slashderiv}
Let $k\in\Z$, $\gamma\in\lfourteenSL$, $z\in\lfourteenH$, and let $f$ be complex-differentiable at $\gamma z$. Then
\begin{equation*}
  (f')|_{k+2}\gamma\,(z) \;=\; \frac{d}{dz}\bigl(f|_k\gamma\bigr)(z) \;+\; \frac{k\,c}{cz+d}\,(f|_k\gamma)(z).
\end{equation*}
\end{theorem}
\noindent File \leanpath{VOA/\allowbreak Modularity/\allowbreak SerreDerivative.lean}.
\begin{leancode}
theorem slashC_deriv_eq (k : ℤ) (γ : SL(2, ℤ)) (f : ℂ → ℂ) {z : ℂ} (hz : z ∈ imGT 0)
    (hf : DifferentiableAt ℂ f (mobiusC γ z)) :
    slashC (k + 2) γ (_root_.deriv f) z
      = _root_.deriv (slashC k γ f) z
        + (k : ℂ) * (γ 1 0 : ℂ) * (denomC γ z)⁻¹ * slashC k γ f z
\end{leancode}
\begin{dictionary}
\lean{_root_.deriv f} is Mathlib's complex derivative $f'$ (the prefix \lean{_root_} only disambiguates the name). The
hypothesis \lean{hf} is needed: \lean{deriv} of a non-differentiable function is $0$ by convention, and the identity would
then fail. Only differentiability at the single point $\gamma z$ is required. Same as the textbook identity.
\end{dictionary}
\noindent\emph{Proof.} Chain rule: $\frac{d}{dz}[f(\gamma z)(cz+d)^{-k}] = f'(\gamma z)(cz+d)^{-2}(cz+d)^{-k} -
kc\,f(\gamma z)(cz+d)^{-k-1}$, using $\frac{d}{dz}\gamma z = (cz+d)^{-2}$. Rearranging gives the claim.
\begin{physics}
The anomalous term $kc/(cz+d)$ is why the naive $\tau$-derivative of a modular form of weight $k$ is \emph{not} modular of
weight $k+2$. In CFT it is the reason that inserting the stress tensor $T(z)$ (the field of $\omega$, whose torus zero mode $L_{[0]}$ is what $\partial_\tau$ inserts
into the trace) produces a non-covariant term; it is cured by the quasi-modular $E_2$, Theorem~\ref{l14:thm:G2slash}.
\end{physics}
\begin{status}
Unconditional (beyond the stated differentiability). \texttt{\#print axioms}: the standard three.
\end{status}

\begin{definition}[$G_2$ in the $\C$ picture]\label{l14:def:G2C}
$G_2^{\C}(z) = G_2(z)$ for $z\in\lfourteenH$, where $G_2 = \sum_{c}\sum_{d}{}'(cz+d)^{-2}$ is Mathlib's (Eisenstein-summed)
weight-two Eisenstein series, $G_2 = 2\zeta(2)E_2 = \frac{\pi^2}{3}(1-24\sum\sigma_1(n)q^n)$.
\end{definition}
\noindent File \leanpath{VOA/\allowbreak Modularity/\allowbreak SerreDerivative.lean}.
\begin{leancode}
noncomputable def G2C (z : ℂ) : ℂ := EisensteinSeries.G2 (UpperHalfPlane.ofComplex z)
\end{leancode}
\begin{dictionary}
\lean{UpperHalfPlane.ofComplex} is Mathlib's retraction $\C\to\lfourteenH$, the identity on $\lfourteenH$ and an
arbitrary point elsewhere; so \lean{G2C} is $G_2$ on $\lfourteenH$ and junk outside (\lean{G2C_coe} records the agreement
on $\lfourteenH$).
\end{dictionary}
\begin{physics}
$G_2$ (equivalently $E_2$) is the torus one-point function of the stress tensor, up to normalisation; see the physics box of Theorem~\ref{l14:thm:G2slash}.
\end{physics}

\begin{theorem}[Quasimodularity of $G_2$; Mathlib's \lean{G2_slash_action}]\label{l14:thm:G2slash}
For $\gamma\in\lfourteenSL$ and $z\in\lfourteenH$: $(G_2|_2\gamma)(z) = G_2(z) - \dfrac{2\pi i\,c}{cz+d}$.
\end{theorem}
\noindent File \leanpath{VOA/\allowbreak Modularity/\allowbreak SerreDerivative.lean}.
\begin{leancode}
theorem slashC_G2C (γ : SL(2, ℤ)) {z : ℂ} (hz : z ∈ imGT 0) :
    slashC 2 γ G2C z
      = G2C z - 2 * (Real.pi : ℂ) * Complex.I * (γ 1 0 : ℂ) * (denomC γ z)⁻¹
\end{leancode}
\begin{dictionary}
A transcription of Mathlib's \lean{EisensteinSeries.G2_slash_action}, $G_2|_2\gamma = G_2 - D_2(\gamma)$ with
$D_2(\gamma)(z) = 2\pi i c/(cz+d)$, into the $\C$ picture. Not reproved: Mathlib supplies it. Exactly the classical law
(Zagier, \emph{Elliptic modular forms}, eq.~(17)).
\end{dictionary}
\noindent\emph{Proof.} Mathlib. Classically: $G_2$ is defined by an iterated, conditionally convergent sum; exchanging the
order of summation under $\tau\mapsto-1/\tau$ produces the defect $2\pi i/\tau$, computed by Hecke's trick
$G_2 = \lim_{s\to0}\sum(m\tau+n)^{-2}|m\tau+n|^{-s}$.
\begin{physics}
$E_2$ is, up to normalisation, the torus one-point function of the stress tensor in the vacuum of a theory with no
states but the vacuum descendants of the identity: $\langle T\rangle_{\text{torus}}\propto\partial_\tau\log\eta$, and
$\eta(\tau)^{24}=\Delta$ has weight $12$, so $\partial_\tau\log\eta = \frac{2\pi i}{24}E_2$ fails to be modular exactly by
the anomaly of Theorem~\ref{l14:thm:slashderiv}. The defect $2\pi ic/(cz+d)$ is the conformal anomaly of the stress tensor
on the torus.
\end{physics}
\begin{status}
Unconditional. Standard axioms.
\end{status}

%======================================================================
\section{Eisenstein series in Zhu's normalisation}\label{l14:sec:eis}
%======================================================================

Zhu's recursion (Chapter~\ref{ch:L15}) produces, from traces of modes, the coefficients of a specific normalisation of the
Eisenstein series, which we call $\mathcal G_{2j}$; it is \emph{neither} $E_{2j}$ (constant term $1$) nor the lattice sum
$G_{2j} = 2\zeta(2j)E_{2j}$. Getting this normalisation wrong is, as the library itself warns, the most likely failure mode
of the whole modularity development, so we are careful.

\begin{definition}[Zhu's Eisenstein coefficients]\label{l14:def:zhuG}
For $j\in\N$ let $g_j:\N\to\C$ be
\begin{equation*}
  g_j(0) = -\frac{B_{2j}}{(2j)!},\qquad g_j(m) = \frac{2}{(2j-1)!}\,\sigma_{2j-1}(m)\quad(m\ge1),
\end{equation*}
and $\mathcal G_{2j} = \sum_{m\ge0}g_j(m)q^m$.
\end{definition}
\noindent File \leanpath{VOA/\allowbreak Modularity/\allowbreak ZhuNormalization.lean} (a Chapter~\ref{ch:L15} file, quoted here).
\begin{leancode}
noncomputable def zhuG (j : ℕ) : ℕ → ℂ :=
  fun m =>
    if m = 0 then -(((bernoulli (2 * j) : ℚ) : ℂ) / (Nat.factorial (2 * j) : ℂ))
    else (2 / (Nat.factorial (2 * j - 1) : ℂ)) * lambertCoeff (2 * j - 1) m
\end{leancode}
\begin{dictionary}
\lean{bernoulli n} is Mathlib's $B_n$ (a rational number, cast to $\C$), with the convention $B_1 = -\tfrac12$ (irrelevant
here, only even indices occur). \lean{lambertCoeff (2j−1) m} is $\sigma_{2j-1}(m)$ by \lean{lambertCoeff_eq_sigma}
(Chapter~\ref{ch:05}, Theorem~\ref{ch05:thm:lambert}). Natural-number subtraction \lean{2 * j - 1} is truncated at $0$; it
is only used for $j\ge1$. At $j=0$ the definition is meaningless junk ($g_0(0) = -B_0 = -1$) and is never used.
\textbf{Normalisation check.} With $E_{2j} = 1 - \frac{4j}{B_{2j}}\sum\sigma_{2j-1}(m)q^m$ one has
$g_j(0)E_{2j} = -\frac{B_{2j}}{(2j)!} + \frac{4j}{(2j)!}\sum\sigma_{2j-1}q^m$, and $\frac{4j}{(2j)!} =
\frac{2}{(2j-1)!}$; so $\mathcal G_{2j} = g_j(0)\,E_{2j}$ (proved analytically, Theorem~\ref{l14:thm:zhuGCE}). Equivalently,
$\mathcal G_{2j} = G_{2j}/(2\pi i)^{2j}$ with $G_{2j}$ the lattice sum: Zhu's $G_{2k}(\tau)$ in J.~AMS 9 (1996), §3. The library
records (in \leanpath{ZhuNormalization.lean}) that the alternatives ``constant term $1$'' and ``constant term $0$'' are
refuted by the Heisenberg computation.
\end{dictionary}
\begin{physics}
In Zhu's derivation the Eisenstein series arise as $\sum_{m\neq0} m^{2j-1}q^m/(1-q^m)$-type sums from moving a mode
$a_{(n)}$ once around the torus: each passage around the time circle costs $q^{n}$, and summing the geometric series over
the winding produces Lambert series. The constant term $-B_{2j}/(2j)!$ comes from the zero-winding (``plane'') contribution
and is the Taylor coefficient of $z^{2j}$ in the Weierstrass-type function $\wp$ on the cylinder. A heuristic gloss on a
definition.
\end{physics}

The library also contains a second, older formal family in the normalisation $E_k$:

\noindent File \leanpath{VOA/\allowbreak Modularity/\allowbreak EisensteinFormal.lean}.
\begin{leancode}
noncomputable def eisensteinCoeff (k : ℕ) (m : ℕ) : ℂ :=
  if m = 0 then 1
  else -(2 * (k : ℂ) / ((bernoulli k : ℚ) : ℂ))
        * ((ArithmeticFunction.sigma (k - 1) m : ℕ) : ℂ)

noncomputable def zhuEisenstein (j : ℕ) : ℕ → ℂ := eisensteinCoeff (2 * j)

def ZhuRecursionEisenstein (N : OrdVModule W M) (a v : 𝒱) (ka : ℕ) : Prop :=
  ZhuRecursionAt N zhuEisenstein a v ka
\end{leancode}
\begin{dictionary}
\lean{eisensteinCoeff k m} is the $m$-th coefficient of $E_k = 1-\frac{2k}{B_k}\sum\sigma_{k-1}(m)q^m$, exactly Mathlib's
formula \lean{EisensteinSeries.E_qExpansion_coeff}; and \lean{zhuEisenstein j} $=E_{2j}$. \lean{ZhuRecursionEisenstein}
instantiates Zhu's recursion (\lean{ZhuRecursionAt}, Chapter~\ref{ch:L15}) at this family. It is a \lean{Prop}, never proved.
\end{dictionary}
\begin{warning}
\textbf{Superseded normalisation; stale docstring (documentation only).} The recursion that the library proves (Chapter~\ref{ch:L15}) uses
\lean{zhuG}, not \lean{zhuEisenstein}; since $g_j(0)\neq1$, \lean{ZhuRecursionEisenstein} is the recursion with the
\emph{wrong} normalisation and would be false. It has no consumer. The module docstring of \leanpath{EisensteinFormal.lean}
still says that the link to Mathlib's $E_k$ is ``deferred'' and that the normalisation is unsettled; both were settled
(\leanpath{EisensteinAnalytic.lean}, \leanpath{QExpansionBridge.lean}, \leanpath{ZhuNormalization.lean}).
\end{warning}

\begin{theorem}[$B_{2j}\neq0$]\label{l14:thm:bernoulli}
For every $j\ge1$, $B_{2j}\neq0$.
\end{theorem}
\noindent File \leanpath{VOA/\allowbreak Modularity/\allowbreak EisensteinAnalytic.lean}.
\begin{leancode}
theorem bernoulli_two_mul_ne_zero {j : ℕ} (hj : j ≠ 0) :
    ((bernoulli (2 * j) : ℚ) : ℂ) ≠ 0
\end{leancode}
\noindent\emph{Proof.} Euler: $\zeta(2j) = (-1)^{j+1}\frac{(2\pi)^{2j}B_{2j}}{2(2j)!}$ (Mathlib \lean{riemannZeta_two_mul_nat}),
and $\zeta(s)\neq0$ for $\lfourteenRe s>1$ (Euler product; Mathlib \lean{riemannZeta_ne_zero_of_one_lt_re}).
\begin{physics}
A purely arithmetic lemma; it ensures $g_j(0)\neq0$, i.e.\ that Zhu's Eisenstein functions have a nonzero value at the cusp,
which is needed when the cusp is used as a reference point (Section~\ref{l14:sec:cusp}).
\end{physics}
\begin{status}
Unconditional. Standard axioms. It discharges a hypothesis that \leanpath{ZhuNormalization.lean} once carried.
\end{status}

\begin{definition}[Zhu's Eisenstein function]\label{l14:def:zhuGC}
For $j\in\N$ and $z\in\C$, $\mathcal G_{2j}(z) = \sum_{n\ge0}g_j(n)\,e^{2\pi inz}$; and $q_{\C}(z) = e^{2\pi iz}$.
\end{definition}
\noindent File \leanpath{VOA/\allowbreak Modularity/\allowbreak EisensteinAnalytic.lean}.
\begin{leancode}
noncomputable def qC (z : ℂ) : ℂ := Complex.exp (2 * (Real.pi : ℂ) * Complex.I * z)

noncomputable def zhuGC (j : ℕ) (z : ℂ) : ℂ :=
  ∑' n : ℕ, zhuG j n * Complex.exp (2 * (Real.pi : ℂ) * Complex.I * (n : ℂ) * z)
\end{leancode}
\begin{dictionary}
A \lean{tsum}: junk where the series diverges. On $\lfourteenH$ it converges absolutely
(\lean{summable_norm_zhuG_qexp}, from $|g_j(n)|\le|g_j(0)| + Cn^{2j}$ and $|q|<1$), so on $\lfourteenH$ it is the honest sum.
\end{dictionary}
\begin{physics}
$\mathcal G_{2j}(\tau)$ is the Eisenstein coefficient of the genus-one Ward identity evaluated on an actual torus of modulus $\tau$, rather than as a formal series.
\end{physics}

\begin{theorem}[Identification with Mathlib's Eisenstein series]\label{l14:thm:zhuGCE}
For $j\ge2$ and $z\in\lfourteenH$, $\mathcal G_{2j}(z) = g_j(0)\,E_{2j}(z)$, where $E_{2j}$ is Mathlib's normalised Eisenstein
modular form of weight $2j$. Consequently, for all $\gamma\in\lfourteenSL$, $\mathcal G_{2j}|_{2j}\gamma = \mathcal G_{2j}$ on
$\lfourteenH$. (Note $g_j(0) = -B_{2j}/(2j)!\neq1$: Zhu's $\mathcal G_{2j}$ is \emph{not} $E_{2j}$, but a nonzero constant multiple of it.)
\end{theorem}
\noindent File \leanpath{VOA/\allowbreak Modularity/\allowbreak EisensteinAnalytic.lean}.
\begin{leancode}
theorem zhuGC_eq_mul_E {j : ℕ} (hj : 2 ≤ j) {z : ℂ} (hz : z ∈ imGT 0) :
    zhuGC j z = zhuG j 0 * ModularForm.E (show 3 ≤ 2 * j by omega) (⟨z, hz⟩ : UpperHalfPlane)

theorem slashC_zhuGC {j : ℕ} (hj : 2 ≤ j) (γ : SL(2, ℤ)) {z : ℂ} (hz : z ∈ imGT 0) :
    slashC (2 * j : ℤ) γ (zhuGC j) z = zhuGC j z
\end{leancode}
\begin{dictionary}
\lean{ModularForm.E h} is Mathlib's normalised Eisenstein series of weight $k$, needing a proof \lean{h : 3 ≤ k}, here
supplied inline by \lean{omega}. \lean{⟨z, hz⟩ : UpperHalfPlane} packages $z$ with the proof that $\lfourteenIm z>0$. The
hypothesis $j\ge2$ is genuine: at $j=1$ there is no weight-two modular form, and the correct statement is
Theorem~\ref{l14:thm:zhuGCone}. So these are exactly the textbook statements, with Mathlib's $E_{2j}$ used, not reproved.
\end{dictionary}
\noindent\emph{Proof.} Mathlib's \lean{q_expansion_bernoulli} gives $E_{2j} = 1 - \frac{4j}{B_{2j}}\sum_{n\ge1}\sigma_{2j-1}(n)q^n$
on $\lfourteenH$; multiply by $g_j(0)$ and use $\frac{4j}{(2j)!} = \frac2{(2j-1)!}$. Modularity then follows from that of
$E_{2j}$ (Mathlib's \lean{slash_action_eqn}).
\begin{physics}
For $j\ge2$ the coefficients that appear in Zhu's genus-one recursion are honest modular forms. This is what makes the
recursion $\lfourteenSL$-covariant: the torus Ward identities have modular coefficients, so the space of solutions (the
conformal blocks on the torus) is a representation of the mapping class group.
\end{physics}
\begin{status}
Unconditional. Standard axioms (both).
\end{status}

\begin{theorem}[The anomalous weight $j=1$]\label{l14:thm:zhuGCone}
For $z\in\lfourteenH$: $\sum_{n\ge0}g_1(n)q^n = -E_2(z)/12$, and $\mathcal G_2(z) = (2\pi i)^{-2}G_2(z)$.
\end{theorem}
\noindent File \leanpath{VOA/\allowbreak Modularity/\allowbreak EisensteinBridge.lean}.
\begin{leancode}
theorem zhuG_one_eq_neg_E2_div_twelve (z : ℍ) :
    zhuG 1 0 + ∑' n : ℕ+, zhuG 1 n * (qParam z) ^ (n : ℕ) = -(E2 z) / 12
\end{leancode}
\noindent File \leanpath{VOA/\allowbreak Modularity/\allowbreak EisensteinAnalytic.lean}.
\begin{leancode}
theorem zhuGC_one {z : ℂ} (hz : z ∈ imGT 0) :
    zhuGC 1 z = ((2 * (Real.pi : ℂ) * Complex.I)⁻¹) ^ 2 * G2C z
\end{leancode}
\begin{dictionary}
\lean{ℕ+} is the type of positive integers, so the \lean{tsum} runs over $n\ge1$. \lean{qParam z} $=e^{2\pi iz}$
(Definition~\ref{l14:def:qeval}). \lean{E2} is Mathlib's $E_2$. Arithmetic check: $g_1(0) = -B_2/2! = -\tfrac1{12}$ and
$g_1(n) = 2\sigma_1(n)$, so the left side is $-\tfrac1{12}(1 - 24\sum\sigma_1q^n)$. And $(2\pi i)^{-2}G_2 =
-\frac{1}{4\pi^2}\cdot\frac{\pi^2}{3}E_2 = -E_2/12$. The \lean{tsum} converges (Mathlib's \lean{E2_eq_tsum_cexp} is proved
for the convergent series), so there is no junk-value issue. Same as textbook.
\end{dictionary}
\begin{physics}
The weight-two term of Zhu's recursion is the only non-modular one, and it is exactly $E_2$: this is the torus conformal
anomaly again (Theorem~\ref{l14:thm:G2slash}). In Zhu's argument it is absorbed into the Serre derivative
(Section~\ref{l14:sec:serre}): the $L_{[0]}$-insertion on the torus is $q\frac{d}{dq}$ plus an $E_2$ correction.
\end{physics}
\begin{status}
Unconditional. Standard axioms.
\end{status}

\begin{theorem}[Values at the cusp]\label{l14:thm:cusplim}
(i) Every modular form $f$ (for a subgroup with the cusp $\infty$) satisfies $f(\tau)\to f(\infty)$ as
$\lfourteenIm\tau\to\infty$. (ii) $E_k(\infty) = 1$ for even $k\ge3$. (iii) $\mathcal G_{2j}(\tau)\to g_j(0)$ as
$\lfourteenIm\tau\to\infty$, for $j\ge2$.
\end{theorem}
\noindent File \leanpath{VOA/\allowbreak Modularity/\allowbreak EisensteinCuspLimit.lean}.
\begin{leancode}
theorem valueAtInfty_E {k : ℕ} (hk : 3 ≤ k) (hk2 : Even k) :
    valueAtInfty (fun τ : UpperHalfPlane => ModularForm.E hk τ) = 1

theorem tendsto_zhuGC_atImInfty {j : ℕ} (hj : 2 ≤ j) :
    Tendsto (fun τ : UpperHalfPlane => zhuGC j (τ : ℂ)) atImInfty (𝓝 (zhuG j 0))
\end{leancode}
\begin{dictionary}
\lean{Tendsto f atImInfty (𝓝 L)} is ``$f(\tau)\to L$ as $\lfourteenIm\tau\to\infty$'' in Mathlib's filter language
(Chapter~\ref{ch:01}); \lean{valueAtInfty f} is Mathlib's name for that limit (a \lean{limUnder}, junk if the limit does not
exist). Statement (i) is \lean{tendsto_modularForm_atImInfty}, stated for any \lean{ModularFormClass} over a subgroup $\Gamma$
of $\mathrm{GL}_2(\mathbb R)$ with $\infty$ a cusp and determinant $\pm1$ (type-class hypotheses, all satisfied by
$\lfourteenSL$). Same as textbook.
\end{dictionary}
\noindent\emph{Proof.} (i) Mathlib's \lean{exp_decay_sub_atImInfty'}: $f - f(\infty) = O(e^{-c\lfourteenIm\tau})$. (ii) The
constant term of the $q$-expansion of $E_k$ is $1$ and equals $f(\infty)$ (\lean{qExpansion_coeff_zero}). (iii) From (i),
(ii) and Theorem~\ref{l14:thm:zhuGCE}.
\begin{physics}
$\lfourteenIm\tau\to\infty$ is the limit of a very long torus: the time circle is infinitely long, so only the ground state
propagates (``$q\to0$ projects onto the vacuum''). The cusp values of the Eisenstein coefficients therefore describe Zhu's
recursion on the infinite cylinder, i.e.\ the algebra of zero modes on the ground states; Section~\ref{l14:sec:cusp} makes
this precise.
\end{physics}
\begin{status}
Unconditional. Standard axioms.
\end{status}

%======================================================================
\section{The level-one structure theorem}\label{l14:sec:levelone}
%======================================================================

\subsection{The algebra generated by \texorpdfstring{$E_4$}{E4} and \texorpdfstring{$E_6$}{E6}}

\begin{definition}[$\lfourteenE$]\label{l14:def:E46}
$\lfourteenE$ is the $\C$-subalgebra of the algebra of all functions $\lfourteenH\to\C$ (pointwise operations) generated
by $E_4$ and $E_6$; its elements are the functions $P(E_4,E_6)$ for $P\in\C[X,Y]$ \emph{not necessarily homogeneous}.
\end{definition}
\noindent File \leanpath{VOA/\allowbreak Modularity/\allowbreak LevelOneStructure.lean}.
\begin{leancode}
noncomputable def E46 : Subalgebra ℂ (UpperHalfPlane → ℂ) :=
  Algebra.adjoin ℂ {(⇑ModularForm.E₄ : UpperHalfPlane → ℂ), (⇑ModularForm.E₆ : UpperHalfPlane → ℂ)}
\end{leancode}
\begin{dictionary}
\lean{Algebra.adjoin ℂ s} is the smallest $\C$-subalgebra containing the set $s$; \lean{⇑F} is the underlying function of a
bundled modular form. So \lean{E46} is a set of functions, not a graded ring: a sum $E_4 + E_6$ of different weights is an
element. This is a deliberate simplification (``no graded-ring plumbing''), and costs nothing for the use in
Chapter~\ref{ch:L15}, which needs only finite generation. It is \emph{not} Mathlib's graded ring \lean{ModularForm.ring}.
\end{dictionary}
\begin{physics}
$\lfourteenE$ is the ring in which the coefficients of modular linear differential equations live (Mathur--Mukhi--Sen): every
torus amplitude of a rational CFT satisfies an ODE in $\tau$ whose coefficients are polynomials in $E_4$, $E_6$.
\end{physics}

\begin{theorem}[Every level-one modular form is a polynomial in $E_4$, $E_6$]\label{l14:thm:levelone}
For every $k\in\Z$ and every $f\in M_k(\lfourteenSL)$, the function $f$ lies in $\lfourteenE$. In particular $E_k\in\lfourteenE$
for $k\ge3$ and $\mathcal G_{2j}\in\lfourteenE$ for $j\ge2$.
\end{theorem}
\noindent File \leanpath{VOA/\allowbreak Modularity/\allowbreak LevelOneStructure.lean}.
\begin{leancode}
theorem levelOne_coe_mem_E46 : ∀ (k : ℤ) (f : ModularForm 𝒮ℒ k),
    (⇑f : UpperHalfPlane → ℂ) ∈ E46

theorem E_mem_E46 {k : ℕ} (hk : 3 ≤ k) :
    (⇑(ModularForm.E hk) : UpperHalfPlane → ℂ) ∈ E46

theorem zhuGC_mem_E46 {j : ℕ} (hj : 2 ≤ j) :
    (fun z : UpperHalfPlane => zhuGC j (z : ℂ)) ∈ E46
\end{leancode}
\begin{dictionary}
\lean{ModularForm 𝒮ℒ k} is Mathlib's space $M_k$ of holomorphic level-one modular forms of weight $k$, holomorphic at the cusp
(Chapter~\ref{ch:05}, Definition~\ref{ch05:def:modform}). No hypotheses. \textbf{Weaker than the textbook statement.} The
textbook ``$M_* = \C[E_4,E_6]$'' has three parts: (a) generation, (b) homogeneity (a weight-$k$ form is a combination of
$E_4^aE_6^b$ with $4a+6b=k$), (c) algebraic independence of $E_4$, $E_6$. \lean{levelOne_coe_mem_E46} is (a) only. Parts (b)
and (c) are \emph{true and proved in the library}, but in a different file: the vendored
\lean{ModularForm.levelOne_eq_polynomial} in \leanpath{VOA/\allowbreak LatticesAndCodes/\allowbreak ModularForms/\allowbreak LevelOne.lean}
(Chapter~\ref{ch:02}) gives, for even $k\ge0$, a basis of $M_k$ consisting of the monomials $E_4^aE_6^b$, $4a+6b=k$. The
docstring's name ``structure theorem'' should be read as ``generation''. \emph{Could it be vacuous?} No: $\lfourteenE$ is a
proper subalgebra of all functions, and membership is a genuine statement.
\end{dictionary}
\noindent\emph{Proof (Serre's induction).} Negative and odd weights: $M_k = 0$ (Mathlib). $k=0$: constants. $k=2$: $M_2 = 0$
(Mathlib). For even $k\ge4$ pick a monomial $F = E_4^aE_6^b$ of weight $k$; it tends to $1$ at the cusp
(\lean{exists_unit_modularForm}, by induction on $k$ in steps of $4$, from $E_4,E_6\to1$). Put $c = f(\infty)$. Then
$f - cF$ has $q$-expansion with zero constant term, so it is a cusp form; Mathlib's isomorphism
\lean{CuspForm.discriminantEquiv} $:S_k\cong M_{k-12}$ divides it by $\Delta$, and $f = cF + \Delta\cdot g$ with
$g\in M_{k-12}$. By induction $g\in\lfourteenE$; and $\Delta = (E_4^3-E_6^2)/1728\in\lfourteenE$ (Mathlib's
\lean{discriminant_eq_E₄_cube_sub_E₆_sq}). For $\mathcal G_{2j}$ use Theorem~\ref{l14:thm:zhuGCE}.
\begin{physics}
The heuristic content (Chapter~\ref{ch:05}): modular covariant partition functions have very little freedom; any weight-$k$
modular object is fixed by finitely many leading $q$-coefficients. Here it is used for a structural reason: the
coefficients of Zhu's genus-one relations lie in a \emph{finitely generated, Noetherian} ring, so the relation module is
finitely generated and dimension counts are possible (Chapter~\ref{ch:L15}).
\end{physics}
\begin{status}
Unconditional. \texttt{\#print axioms levelOne\_coe\_mem\_E46}, \texttt{zhuGC\_mem\_E46}: the standard three (run for this
chapter). Mathlib's dimension formula and $\Delta$-division are used, not reproved.
\end{status}

\begin{warning}
\textbf{Duplicate development (documentation only).} The library proves level-one generation twice: once here (non-homogeneous, functions)
and once, with homogeneity and independence, in the vendored \leanpath{LatticesAndCodes/\allowbreak ModularForms/\allowbreak LevelOne.lean}
(\lean{span_eisBasisFun}, \lean{levelOne_eq_polynomial}). Neither file mentions the other. The stronger vendored theorem
discharges the predicate below that \leanpath{EisensteinRamanujan.lean} used to record as ``missing''
(Theorem~\ref{l14:thm:certreps}; now the library theorem \lean{e46HasWeightedPolynomialReps}).
\end{warning}

\subsection{Weights inside \texorpdfstring{$\lfourteenE$}{E46}, and the Ramanujan boundary}\label{l14:sec:ramanujan}

Because $\lfourteenE$ is not graded, the library adds a weight predicate. We separate what is \emph{defined and used} (the
weight predicate and homogeneous representatives, Definition~\ref{l14:def:E46Weight}) from three \emph{named propositions}
(Definition~\ref{l14:def:serreclosed}), all three now theorems of the library.

\paragraph{The Serre derivative, in advance.} Serre closure involves the \emph{Serre derivative}, studied in
Section~\ref{l14:sec:serre} (Definition~\ref{l14:def:serre}). For the reader's convenience we give it now: for $k\in\Z$ and
$f:\C\to\C$,
\begin{equation*}
  \partial_kf \;=\; \frac{1}{2\pi i}f' + \frac{k}{(2\pi i)^2}\,G_2\,f \;=\; q\frac{df}{dq} - \frac{k}{12}E_2f\qquad\text{on }\lfourteenH,
\end{equation*}
with $G_2$ the weight-two Eisenstein series of Definition~\ref{l14:def:G2C} ($G_2 = \frac{\pi^2}3E_2$). The point of the
correction term is that $\partial_k$ commutes with slashing (weight $k$ in, weight $k+2$ out), which the plain derivative does
not (Theorem~\ref{l14:thm:slashderiv}); this equivariance is proved in Theorem~\ref{l14:thm:serreslash}. In Lean it is
\lean{serreDerivC k f z}.

\begin{definition}[Weighted elements; homogeneous representatives]\label{l14:def:E46Weight}
For $f\in\lfourteenE$ and $k\in\Z$: $f$ \emph{has weight $k$} if $f$ is (the function of) some $F\in M_k$. $f$ \emph{has a
homogeneous polynomial representative of weight $k$} if $f = \sum_{(a,b)\in t}c_{ab}E_4^aE_6^b$ for a finite set $t$ of
pairs with $4a+6b = k$.
\end{definition}
\noindent File \leanpath{VOA/\allowbreak Modularity/\allowbreak EisensteinRamanujan.lean}.
\begin{leancode}
noncomputable def e46C (f : E46) (z : ℂ) : ℂ :=
  (f : UpperHalfPlane → ℂ) (UpperHalfPlane.ofComplex z)

def E46Weight (k : ℤ) (f : E46) : Prop :=
  ∃ F : ModularForm 𝒮ℒ k, (f : UpperHalfPlane → ℂ) = F

noncomputable def e46E4 : E46 := ⟨⇑ModularForm.E₄, E4_mem_E46⟩

noncomputable def e46E6 : E46 := ⟨⇑ModularForm.E₆, E6_mem_E46⟩

noncomputable def e46Monomial (a b : ℕ) : E46 := e46E4 ^ a * e46E6 ^ b

def E46WeightedPolynomialRep (k : ℤ) (f : E46) : Prop :=
  ∃ (t : Finset (ℕ × ℕ)) (c : ℕ × ℕ → ℂ),
    f = ∑ ab ∈ t, c ab • e46Monomial ab.1 ab.2 ∧
      ∀ ab ∈ t, 4 * (ab.1 : ℤ) + 6 * (ab.2 : ℤ) = k
\end{leancode}
\begin{dictionary}
\lean{e46C f} extends $f$ from $\lfourteenH$ to $\C$ through Mathlib's retraction \lean{ofComplex}; only its values on
\lean{imGT 0} matter, and all equalities below are restricted there. \lean{E46Weight k f} is an existential over modular forms:
it retains holomorphy at the cusp, so it is the right notion of ``coefficient of weight $k$''. These are predicates about a
given $f$; they assert nothing by themselves.
\end{dictionary}
\begin{physics}
A weight-$k$ element is a candidate coefficient of a modular linear differential equation: the coefficient of
$\partial^{n-r}$ in an order-$n$ MLDE must be a weight-$2r$ modular form.
\end{physics}
\begin{status}
Definitions. Proved closure properties: \lean{E46Weight} is closed under $+,-,\cdot,$ scalars and finite sums
(\lean{E46Weight.add}, \dots, \lean{E46Weight.sum}); $E_4^aE_6^b$ has weight $4a+6b$ (\lean{e46Monomial_weight}). Standard axioms.
\end{status}

\begin{definition}[Three named propositions: homogeneous representatives, monomial and full Serre closure]\label{l14:def:serreclosed}
(a) \emph{Homogeneous representatives}: every $f\in\lfourteenE$ of weight $k$ has a homogeneous polynomial representative of
weight $k$. (b) \emph{Monomial Serre closure}: for all $a,b$, $\partial_{4a+6b}(E_4^aE_6^b)$ equals on $\lfourteenH$ an element of
$\lfourteenE$ of weight $4a+6b+2$. (c) \emph{Serre closure}: for every $f$ of weight $k$ there is $g\in\lfourteenE$ of weight $k+2$
with $g = \partial_kf$ on $\lfourteenH$.
\end{definition}
\noindent File \leanpath{VOA/\allowbreak Modularity/\allowbreak EisensteinRamanujan.lean}.
\begin{leancode}
def E46HasWeightedPolynomialReps : Prop :=
  ∀ (k : ℤ) (f : E46), E46Weight k f → E46WeightedPolynomialRep k f

def E46SerreMonomialClosed : Prop :=
  ∀ (a b : ℕ), ∃ g : E46,
    E46Weight (4 * (a : ℤ) + 6 * (b : ℤ) + 2) g ∧
      ∀ z ∈ imGT 0,
        e46C g z = serreDerivC (4 * (a : ℤ) + 6 * (b : ℤ))
          (e46C (e46Monomial a b)) z

def E46SerreClosed : Prop :=
  ∀ (k : ℤ) (f : E46), E46Weight k f →
    ∃ g : E46, E46Weight (k + 2) g ∧
      ∀ z ∈ imGT 0, e46C g z = serreDerivC k (e46C f) z
\end{leancode}
\begin{dictionary}
Each is a \lean{def … : Prop}, i.e.\ a \emph{statement}; all three are proved separately (status box below). (c) is the classical fact $\partial_kM_k\subseteq M_{k+2}$.
\textbf{Equivariance is not closure.} Theorem~\ref{l14:thm:serreslash} (proved) shows that $\partial_kf$ transforms with weight
$k+2$; what it does not give is that $\partial_kf$ is \emph{holomorphic at the cusp} with a $q$-expansion, i.e.\ is again in
$M_{k+2}$, nor that it lies in $\lfourteenE$; that extra content (classically: Ramanujan's identities plus the structure theorem)
is exactly (b)/(c). The library proves the aggregation step
\lean{e46SerreClosed_of_weightedPolynomialReps : E46HasWeightedPolynomialReps → E46SerreMonomialClosed → E46SerreClosed}.
Mathlib (\leanpath{NumberTheory/\allowbreak ModularForms/\allowbreak Derivative.lean}) defines its own \lean{serreDerivative} on
\lean{ℍ → ℂ} with the same normalisation ($DF - \frac{k}{12}E_2F$) but proves no modularity-preservation theorem for it.
\end{dictionary}
\begin{physics}
Serre closure is the statement that the ring of modular forms is closed under the covariant derivative of the torus
moduli space. Physically it is what makes a modular linear differential equation (MLDE) $[\partial^n + \phi_4\partial^{n-2}+
\cdots]\chi = 0$ meaningful: each $\partial$ raises the weight by $2$ and the coefficient $\phi_{2r}$ must be a weight-$2r$
modular form. The monomial case is Ramanujan's identities $\partial_4E_4 = -\frac13E_6$, $\partial_6E_6 = -\frac12E_4^2$.
\end{physics}
\begin{status}
\textbf{Status of the three propositions.}
\begin{itemize}
\item (a) \lean{E46HasWeightedPolynomialReps}: \textbf{proved} in the library (\lean{e46HasWeightedPolynomialReps},
  from the argument of Theorem~\ref{l14:thm:certreps}).
\item (b) \lean{E46SerreMonomialClosed}: \textbf{proved} (\lean{e46SerreMonomialClosed}, the monomial case of (c)).
\item (c) \lean{E46SerreClosed}: \textbf{proved} (\lean{e46SerreClosed}). The proof does not go through Ramanujan's
  identities: \lean{serreForm} packages $\partial_nF$, for $F$ a level-one form of weight $n$, as a level-one form of weight
  $n+2$ (holomorphy on $\lfourteenH$ and boundedness at $i\infty$ from the $q$-expansion, the weight-$(n+2)$ law from
  Theorem~\ref{l14:thm:serreslash}), and \lean{levelOne_coe_mem_E46} places it in $\lfourteenE$. The reduction
  \lean{e46SerreClosed_of_weightedPolynomialReps} gives a second proof (\lean{e46SerreClosed'}). The hypothesis
  \lean{hclose} of the Arakawa--Kawasetsu MLDE of Chapter~\ref{ch:L16} is thereby discharged.
  (\leanpath{VOA/\allowbreak Modularity/\allowbreak EisensteinSerreClosure.lean}; not compiled for this chapter.)
\end{itemize}
\end{status}

\begin{theorem}[Homogeneous representatives exist]\label{l14:thm:certreps}
Every $f\in\lfourteenE$ of weight $k$ is a linear combination of monomials $E_4^aE_6^b$ with $4a+6b=k$; that is,
\lean{E46HasWeightedPolynomialReps} holds.
\end{theorem}
\noindent File \leanpath{VOA/\allowbreak Modularity/\allowbreak EisensteinSerreClosure.lean}.
\begin{plaincode}
theorem e46HasWeightedPolynomialReps : E46HasWeightedPolynomialReps
\end{plaincode}
\noindent\emph{Proof.} If $k<0$, $F=0$. If $k = n\ge0$, the vendored \lean{ModularForm.span_eisBasisFun n} says the
weight-$n$ monomials span $M_n$; write $F$ as a finite combination, push it through the linear coercion \lean{coeLm}, and
read the result as an identity in $\lfourteenE$. (Odd $n$: the index set \lean{wtIdx n} is empty and $F=0$.)
\begin{physics}
Arithmetic bookkeeping: the statement that ``a weight-$k$ coefficient is a homogeneous polynomial''. It reduces Serre closure
to Ramanujan's two identities.
\end{physics}
\begin{status}
First compiled for this chapter as a separate check with \texttt{lake env lean -DmaxSynthPendingDepth=3} at an earlier commit (not re-run for this edition): no errors, no
\lean{sorry}. The library has since adopted this proof as \lean{e46HasWeightedPolynomialReps}
(\leanpath{EisensteinSerreClosure.lean}), and the docstring of \lean{E46HasWeightedPolynomialReps}, which used to call it
missing, now points to it.
\end{status}

\subsection{Points of \texorpdfstring{$\lfourteenE$}{E46}: the cusp and the interior}\label{l14:sec:points}

The \emph{cusp} is the boundary point $\tau = i\infty$ of $\lfourteenH$ (the one point added to $\lfourteenSL\backslash\lfourteenH$ to
compactify it, Chapter~\ref{ch:05}). Approaching it means $\lfourteenIm\tau\to\infty$, i.e.\ $q\to0$. Geometrically the torus
degenerates: its time circle becomes infinitely long, the torus becomes a cylinder, and in a trace only the ground states survive.
Every element of $\lfourteenE$ has a limit there, so the cusp behaves like a point of $\lfourteenE$ on the same footing as the
interior points.

\begin{definition}[Evaluation at the cusp and at interior points]\label{l14:def:evInfty}
$\lfourteenev_\infty:\lfourteenE\to\C$, $f\mapsto\lim_{\lfourteenIm\tau\to\infty}f(\tau)$, and for $z_0\in\lfourteenH$,
$\lfourteenev_{z_0}:\lfourteenE\to\C$, $f\mapsto f(z_0)$. Both are $\C$-algebra homomorphisms.
\end{definition}
\noindent File \leanpath{VOA/\allowbreak Modularity/\allowbreak CuspPoint.lean}.
\begin{leancode}
noncomputable def cuspLimSubalg : Subalgebra ℂ (UpperHalfPlane → ℂ) where
  carrier := {f : UpperHalfPlane → ℂ | ∃ L : ℂ, Tendsto f atImInfty (𝓝 L)}
  mul_mem' := fun ⟨L₁, h₁⟩ ⟨L₂, h₂⟩ => ⟨L₁ * L₂, h₁.mul h₂⟩
  add_mem' := fun ⟨L₁, h₁⟩ ⟨L₂, h₂⟩ => ⟨L₁ + L₂, h₁.add h₂⟩
  zero_mem' := ⟨0, tendsto_const_nhds⟩
  one_mem' := ⟨1, tendsto_const_nhds⟩
  algebraMap_mem' := fun c => ⟨c, tendsto_const_nhds⟩

theorem evInfty_zhuGC {j : ℕ} (hj : 2 ≤ j) :
    evInfty ⟨_, zhuGC_mem_E46 hj⟩ = zhuG j 0
\end{leancode}
\noindent File \leanpath{VOA/\allowbreak Modularity/\allowbreak E46Point.lean}.
\begin{leancode}
noncomputable def evAt (z₀ : UpperHalfPlane) : E46 →ₐ[ℂ] ℂ where
  toFun f := (f : UpperHalfPlane → ℂ) z₀
  map_one' := rfl
  map_mul' _ _ := rfl
  map_zero' := rfl
  map_add' _ _ := rfl
  commutes' _ := rfl

theorem ker_isMaximal (φ : E46 →ₐ[ℂ] ℂ) :
    (RingHom.ker (φ : E46 →+* ℂ)).IsMaximal
\end{leancode}
\begin{dictionary}
\lean{evInfty : E46 →ₐ[ℂ] ℂ} (same file \leanpath{CuspPoint.lean}) is \lean{valueAtInfty}, shown to be an algebra hom because
$\lfourteenE\le$ \lean{cuspLimSubalg} (\lean{E46_le_cuspLimSubalg}: $E_4$, $E_6$ have cusp limits, and limits respect
$+$ and $\cdot$). \lean{→ₐ[ℂ]} is ``$\C$-algebra homomorphism''. \lean{ker_isMaximal}: the kernel of any such hom is a maximal
ideal (residue field $\C$). \lean{evInfty_zhuGC}: $\lfourteenev_\infty(\mathcal G_{2j}) = g_j(0)$. Statements are exactly the
textbook ones. They are infrastructure: no theorem in this chapter compares the two fibres (that is the open ``flatness''
residual of Chapter~\ref{ch:L15}).
\end{dictionary}
\begin{physics}
The two kinds of points are two geometries of the torus: an interior point $z_0$ is a torus of fixed shape, the cusp is the
degenerate torus (a pinched cycle = an infinitely long cylinder). Conformal blocks over the family of tori form a sheaf over
the $(E_4,E_6)$-plane; its fibre at the cusp is computed algebraically from $A(V)$ (Section~\ref{l14:sec:cusp}), and
transporting the dimension to an interior point is the statement that the number of torus blocks does not jump.
\end{physics}
\begin{status}
Unconditional, standard axioms (\lean{evInfty_zhuGC}, \lean{ker_evAt_isMaximal}).
\end{status}

\begin{theorem}[No fibre of {$\C[[q]]$} at $q_0\neq0$]\label{l14:thm:noalghom}
There is no $\C$-algebra homomorphism $\varphi:\C[[q]]\to\C$ with $\varphi(q)\neq0$.
\end{theorem}
\noindent File \leanpath{VOA/\allowbreak Modularity/\allowbreak FormalSpecialization.lean}.
\begin{leancode}
theorem no_algHom_powerSeries_X_ne_zero (φ : PowerSeries ℂ →ₐ[ℂ] ℂ) {q₀ : ℂ} (hq : q₀ ≠ 0) :
    φ PowerSeries.X ≠ q₀
\end{leancode}
\noindent\emph{Proof.} $1 - q_0^{-1}q$ has constant term $1$, hence is a unit of $\C[[q]]$; $\varphi$ would send it to
$1 - q_0^{-1}q_0 = 0$, which is not a unit.
\begin{dictionary}
Exactly the stated fact. It is recorded as a \emph{dead end}: a purely formal argument over $\C[[q]]$ can only reach the
special fibre $q=0$ and the generic fibre $\C((q))$, never a genuine torus $0<|q|<1$. This is why the library moved its
coefficient ring from $\C[[q]]$ to the function algebra $\lfourteenE$, where evaluation at $z_0$ \emph{is} a homomorphism.
\end{dictionary}
\begin{physics}
Formal power series know only infinitesimally long tori. A physical statement about a torus of finite size needs
convergence, i.e.\ analysis. The theorem is the algebraic shadow of that fact.
\end{physics}
\begin{status}
Unconditional. Standard axioms.
\end{status}

%======================================================================
\section{From formal \texorpdfstring{$q$}{q}-series to holomorphic functions}\label{l14:sec:bridge}
%======================================================================

\subsection{Subexponential growth and summability}

\begin{definition}[Subexponential sequence]\label{l14:def:subexp}
A real sequence $(a_n)_{n\in\N}$ is \emph{subexponential} if for every $r>1$ there is $C$ with $a_n\le Cr^n$ for all $n$.
\end{definition}
\noindent File \leanpath{VOA/\allowbreak Modularity/\allowbreak ConvergenceBridge.lean}.
\begin{leancode}
def Subexponential (a : ℕ → ℝ) : Prop :=
  ∀ r : ℝ, 1 < r → ∃ C : ℝ, ∀ n : ℕ, a n ≤ C * r ^ n
\end{leancode}
\begin{dictionary}
Quantifiers: $\forall r>1\ \exists C\ \forall n$; $C$ depends on $r$. This is an upper bound only (no absolute value); it is
always applied to $a_n = \|\chi(h+n)\|\ge0$. Equivalent to $\limsup a_n^{1/n}\le1$ for nonnegative $a$, i.e.\ to the power
series $\sum a_nx^n$ having radius of convergence $\ge1$. So it is exactly the condition for convergence on the whole disc
$|q|<1$, the right condition for a character.
\end{dictionary}
\begin{physics}
The degeneracies of a CFT grow like $\dim V_n\sim\exp(2\pi\sqrt{c_{\mathrm{eff}}n/6})$ (Cardy's formula, derived from the
$S$-transformation: the high-temperature behaviour of $Z(\tau)$ is controlled by the vacuum in the dual channel). That growth
is subexponential: $e^{a\sqrt n}\le C_r r^n$ for every $r>1$ (\lean{subexponential_exp_sqrt}). Subexponentiality is the
weakest growth condition that makes the partition function a well-defined function on all of $\lfourteenH$; Cardy growth is
much stronger.
\end{physics}

\begin{theorem}[Closure properties]\label{l14:thm:subexpclosure}
Subexponential sequences are closed under domination, nonnegative scalar multiples, addition, finite sums, polynomial factors
$n^p$, index shifts in both directions, and Cauchy products (for complex sequences, in absolute value); constants and
$e^{c\sqrt n}$ are subexponential. The subexponential formal power series form a $\C$-subalgebra $\mathcal S$ of $\C[[q]]$,
stable under $\theta = q\frac{d}{dq}$.
\end{theorem}
\noindent Files \leanpath{VOA/\allowbreak Modularity/\allowbreak ConvergenceBridge.lean}, \leanpath{GrowthClosure.lean},
\leanpath{SubexpClosure.lean}, \leanpath{SubexpSeriesRing.lean}, \leanpath{QDerivProductRule.lean}. File
\leanpath{VOA/\allowbreak Modularity/\allowbreak ConvergenceBridge.lean}.
\begin{leancode}
theorem subexponential_exp_sqrt (c : ℝ) :
    Subexponential (fun n : ℕ => Real.exp (c * Real.sqrt n))
\end{leancode}
\noindent File \leanpath{VOA/\allowbreak Modularity/\allowbreak GrowthClosure.lean}.
\begin{leancode}
theorem Subexponential.mul_pow {a : ℕ → ℝ} (ha : Subexponential a) (hnn : ∀ n, 0 ≤ a n) (p : ℕ) :
    Subexponential (fun n : ℕ => (n : ℝ) ^ p * a n)
\end{leancode}
\noindent File \leanpath{VOA/\allowbreak Modularity/\allowbreak SubexpClosure.lean}.
\begin{leancode}
theorem Subexponential.of_shift {a : ℕ → ℝ} (k : ℕ)
    (ha : Subexponential (fun n => a (n + k))) : Subexponential a
\end{leancode}
\noindent File \leanpath{VOA/\allowbreak Modularity/\allowbreak SubexpSeriesRing.lean}.
\begin{leancode}
theorem subexponential_cauchy {a b : ℕ → ℂ}
    (ha : Subexponential fun m : ℕ => ‖a m‖) (hb : Subexponential fun m : ℕ => ‖b m‖) :
    Subexponential fun n : ℕ => ‖∑ m ∈ Finset.range (n + 1), a m * b (n - m)‖

def SubexpPS : Subalgebra ℂ (PowerSeries ℂ) where
  carrier := {F : PowerSeries ℂ | Subexponential fun n : ℕ => ‖PowerSeries.coeff n F‖}
\end{leancode}
\noindent File \leanpath{VOA/\allowbreak Modularity/\allowbreak QDerivProductRule.lean}.
\begin{leancode}
noncomputable def theta (F : PowerSeries ℂ) : PowerSeries ℂ :=
  PowerSeries.mk fun n => (n : ℂ) * PowerSeries.coeff n F

theorem theta_mem_subexpPS {F : PowerSeries ℂ} (hF : F ∈ SubexpPS) : theta F ∈ SubexpPS
\end{leancode}
\begin{dictionary}
\lean{‖·‖} is the complex absolute value. \lean{Finset.range (n+1)} is $\{0,\dots,n\}$, so the Cauchy-product sum is
$\sum_{m\le n}a_mb_{n-m}$. The quoted part of \lean{SubexpPS} is its carrier; the closure fields (\lean{mul_mem'} etc.) are
proved from \lean{subexponential_cauchy}. \lean{Subexponential.mul_pow} needs $a\ge0$ (otherwise the inequality could be
reversed by a negative factor). All statements are the textbook ones. \lean{theta F} is $qF'(q)$:
$[q^n]\theta F = n[q^n]F$; \lean{theta_mul} proves it is a derivation.
\end{dictionary}
\noindent\emph{Proof sketch.} $e^{c\sqrt n - n\log r}\le e^{c^2/(4\log r)}$ by completing the square. $n^p\le C_1(\sqrt r)^n$
and $a_n\le C_2(\sqrt r)^n$ give $n^pa_n\le C_1C_2r^n$. Cauchy: choose $1<s<r$; then $\sum_{m\le n}Ds^m\cdot Cr^{n-m}\le
CD(1-s/r)^{-1}r^n$. Shifts absorb finitely many terms into the constant.
\begin{physics}
These are the bookkeeping rules of state counting: tensoring theories (Cauchy product of characters), multiplying by a
polynomial in the energy (inserting $L_0^p$, which is what $\theta^p$ does), and shifting the vacuum energy all preserve the
property ``the partition function converges for every $\tau\in\lfourteenH$''.
\end{physics}
\begin{status}
All unconditional; standard axioms (\lean{subexponential_exp_sqrt}, \lean{theta_mem_subexpPS} checked for this chapter).
\end{status}

\begin{theorem}[Subexponential coefficients converge on the open disc]\label{l14:thm:summable}
If $a_n\ge0$ is subexponential and $0\le t<1$, then $\sum_na_nt^n$ converges.
\end{theorem}
\noindent File \leanpath{VOA/\allowbreak Modularity/\allowbreak ConvergenceBridge.lean}.
\begin{leancode}
theorem summable_of_subexponential {a : ℕ → ℝ} (hpos : ∀ n, 0 ≤ a n) (hsub : Subexponential a)
    {t : ℝ} (ht0 : 0 ≤ t) (ht : t < 1) : Summable (fun n : ℕ => a n * t ^ n)
\end{leancode}
\noindent\emph{Proof.} Choose $1<r<t^{-1}$; compare with $C\sum(rt)^n$.
\begin{physics}
The partition function at inverse temperature $\beta$ ($t = e^{-2\pi\beta}$) is finite.
\end{physics}
\begin{status}
Unconditional; standard axioms.
\end{status}

\subsection{Evaluation}

\begin{definition}[The $q$-parameter; summability at $\tau$; evaluation]\label{l14:def:qeval}
For $\tau\in\lfourteenH$ put $q = e^{2\pi i\tau}$. For $\chi$ a formal $q$-series and a base point $h\in\C$, $\chi$ is
\emph{summable at $\tau$ along $h$} if $\sum_{n\ge0}\chi(h+n)q^n$ converges, and its \emph{evaluation} is
$\chi_h(\tau) = e^{2\pi i\tau h}\sum_{n\ge0}\chi(h+n)q^n$, the formal $\sum_n\chi(h+n)q^{h+n}$ with $q^h := e^{2\pi i\tau h}$.
In the $\C$ picture, $\chi_h^{\C}(z) = e^{2\pi izh}\sum_n\chi(h+n)e^{2\pi inz}$. (The function $q_{\C}(z) = e^{2\pi iz}$ of
Definition~\ref{l14:def:zhuGC} is the same $q$ in the $\C$ picture.) The base point $h$ is the lowest exponent of the series:
for the character of a module of conformal weight $h_M$ it is the ground-state energy $h = h_M - c/24$, so that the evaluation
is $\tr\,q^{L_0-c/24}$ (Chapter~\ref{ch:L12}, Definition~\ref{l12:def:char}).
\end{definition}
\noindent File \leanpath{VOA/\allowbreak Modularity/\allowbreak ConvergenceBridge.lean}.
\begin{leancode}
noncomputable def qParam (τ : UpperHalfPlane) : ℂ :=
  Complex.exp (2 * (Real.pi : ℂ) * Complex.I * (τ : ℂ))

def QSummableAt (χ : FormalQSeries) (h : ℂ) (τ : UpperHalfPlane) : Prop :=
  Summable (fun n : ℕ => χ (h + n) * (qParam τ) ^ n)

noncomputable def qEvalAt (χ : FormalQSeries) (h : ℂ) (τ : UpperHalfPlane) : ℂ :=
  Complex.exp (2 * (Real.pi : ℂ) * Complex.I * (τ : ℂ) * h) *
    ∑' n : ℕ, χ (h + n) * (qParam τ) ^ n
\end{leancode}
\noindent File \leanpath{VOA/\allowbreak Modularity/\allowbreak QEvalHolomorphic.lean}.
\begin{leancode}
noncomputable def qTerm (χ : FormalQSeries) (h : ℂ) (n : ℕ) (z : ℂ) : ℂ :=
  χ (h + (n : ℂ)) * Complex.exp (2 * (Real.pi : ℂ) * Complex.I * (n : ℂ) * z)

noncomputable def qSeriesC (χ : FormalQSeries) (h : ℂ) (z : ℂ) : ℂ := ∑' n : ℕ, qTerm χ h n z

noncomputable def qEvalC (χ : FormalQSeries) (h : ℂ) (z : ℂ) : ℂ :=
  Complex.exp (2 * (Real.pi : ℂ) * Complex.I * z * h) * qSeriesC χ h z
\end{leancode}
\begin{dictionary}
\textbf{Decoding.} The prefactor \lean{exp(2πi τ h)} is $q^h$ with the principal choice $q^h = e^{2\pi i\tau h}$ (no branch
ambiguity, because it is defined through $\tau$, not through $q$). \lean{qEvalAt} reads $\chi$ only along $h+\N$; values of
$\chi$ off that progression are ignored, so \lean{qEvalAt χ h} depends on the choice of $h$, and the caller must choose $h$
to be the lowest exponent (for a character of a module of conformal weight $h_M$, $h = h_M - c/24$). \lean{qEvalC_coe}
proves \lean{qEvalC χ h (τ : ℂ) = qEvalAt χ h τ}. \textbf{Junk value.} If the series diverges, \lean{qEvalAt} is $0$.
Theorems that need the honest value carry \lean{QSummableAt} or a \lean{Subexponential} hypothesis. \textbf{Not Mathlib's
\lean{qExpansion}.} Mathlib goes from functions to coefficients; this goes from coefficients to functions.
\end{dictionary}
\begin{physics}
$\chi_h(\tau) = \tr_Mq^{L_0-c/24}$ is the torus partition function of the sector $M$; the choice $h = h_M - c/24$ is the
ground-state energy on the cylinder.
\end{physics}

\begin{theorem}[The convergence bridge]\label{l14:thm:qsummable}
If $\|\chi(h+n)\|$ is subexponential then $\chi$ is summable at every $\tau\in\lfourteenH$ along $h$. In particular, if the
graded dimensions of a module $N$ along $h+\N$ are subexponential, its character converges on $\lfourteenH$.
\end{theorem}
\noindent File \leanpath{VOA/\allowbreak Modularity/\allowbreak ConvergenceBridge.lean}.
\begin{leancode}
theorem qSummableAt_of_subexponential (χ : FormalQSeries) (h : ℂ)
    (hsub : Subexponential (fun n : ℕ => ‖χ (h + n)‖)) (τ : UpperHalfPlane) :
    QSummableAt χ h τ

def SubexponentialCoeffs (N : OrdVModule W M) (h : ℂ) : Prop :=
  Subexponential (fun n : ℕ => ‖gradedDim N (h + n)‖)
\end{leancode}
\begin{dictionary}
\lean{OrdVModule W M} is an ordinary $V$-module on $M$ (Chapter~\ref{ch:L12}, Definition~\ref{l12:def:ordmodule});
\lean{gradedDim N (h+n)} $=\dim M_{h+n+c/24}$ (cast to $\C$). \lean{SubexponentialCoeffs} is the named ``growth datum''. The
docstring of \leanpath{ConvergenceBridge.lean} says its supply from $C_2$-cofiniteness is ``not built''; that is stale: it
is proved for the adjoint module (\lean{subexponentialCoeffs_adjoint_of_C2Cofinite}, Chapter~\ref{ch:L15}) and for modules
with a finite lowest-weight generating family (\lean{subexponentialCoeffs_of_C2Cofinite}).
\end{dictionary}
\noindent\emph{Proof.} Fix $\tau\in\lfourteenH$ and put $t = |q| = e^{-2\pi\lfourteenIm\tau}$; then $0\le t<1$
(\lean{norm_qParam_lt_one}). Choose any $r$ with $1<r<t^{-1}$ (if $t = 0$ there is nothing to prove), for instance
$r = (1+t^{-1})/2$. By subexponentiality, applied at this $r$, there is $C$ with $\|\chi(h+n)\|\le Cr^n$ for all $n$. Hence
\begin{equation*}
  \sum_{n\ge0}\|\chi(h+n)q^n\| = \sum_{n\ge0}\|\chi(h+n)\|\,t^n \le C\sum_{n\ge0}(rt)^n = \frac{C}{1-rt}<\infty ,
\end{equation*}
since $0\le rt<1$. This is Theorem~\ref{l14:thm:summable} with $a_n = \|\chi(h+n)\|$. An absolutely convergent series of complex
numbers converges, which is \lean{QSummableAt}.
\begin{physics}
Growth bound on degeneracies $\Rightarrow$ finite partition function at every temperature and every twist.
\end{physics}
\begin{status}
Unconditional; standard axioms.
\end{status}

\begin{warning}
\textbf{A wrong docstring claim (documentation only).} The ``Residual'' section of \leanpath{ConvergenceBridge.lean} says that only
``that $p(n)$ is subexponential'' is needed, ``which follows from $p(n)\le2^n$-type bounds''. It does not: $2^n$ is not
subexponential ($2^n\le C(3/2)^n$ fails for large $n$). The library's actual partition bound (Theorem~\ref{l14:thm:pcount})
is the correct generating-function estimate $p_g(n)t^n\le\exp(gt/(1-t)^2)$ for all $t\in(0,1)$, which does give
subexponential growth. No theorem is affected.
\end{warning}

\begin{theorem}[Holomorphy and term-by-term differentiation]\label{l14:thm:holo}
Let $\chi$ be a formal $q$-series with $\|\chi(h+n)\|$ subexponential. Then $z\mapsto\chi^{\C}_h(z)$ is holomorphic on
$\lfourteenH$, and for $z\in\lfourteenH$
\begin{equation*}
  \frac{d}{dz}\chi^{\C}_h(z) = 2\pi i\,(D\chi)^{\C}_h(z),\qquad (D\chi)(e) := e\,\chi(e).
\end{equation*}
In particular the character of a $C_2$-cofinite VOA (with $\mathbf 1\neq0$) is holomorphic on $\lfourteenH$.
\end{theorem}
\noindent File \leanpath{VOA/\allowbreak Modularity/\allowbreak ZhuBlockConditions.lean} (a Chapter~\ref{ch:L15} file).
\begin{leancode}
noncomputable def qDeriv (χ : FormalQSeries) : FormalQSeries := fun e => e * χ e
\end{leancode}
\noindent File \leanpath{VOA/\allowbreak Modularity/\allowbreak QEvalHolomorphic.lean}.
\begin{leancode}
theorem differentiableOn_qEvalC (χ : FormalQSeries) (h : ℂ)
    (hsub : Subexponential (fun n : ℕ => ‖χ (h + (n : ℂ))‖)) :
    DifferentiableOn ℂ (qEvalC χ h) (imGT 0)

theorem deriv_qEvalC (χ : FormalQSeries) (h : ℂ)
    (hsub : Subexponential (fun n : ℕ => ‖χ (h + (n : ℂ))‖)) {z : ℂ} (hz : z ∈ imGT 0) :
    _root_.deriv (qEvalC χ h) z = 2 * (Real.pi : ℂ) * Complex.I * qEvalC (qDeriv χ) h z

theorem differentiableOn_char_adjoint (W : VertexOperatorAlgebra 𝒱) (hne : W.vacuum ≠ 0)
    (hC2 : C2Cofinite W) :
    DifferentiableOn ℂ
      (qEvalC (gradedDim (adjointModule W)) (0 - W.centralCharge / 24)) (imGT 0)
\end{leancode}
\begin{dictionary}
\lean{qDeriv} is the formal Euler operator $q\frac{d}{dq}$ on series with complex exponents: $q\frac{d}{dq}q^e = eq^e$.
\lean{DifferentiableOn ℂ f U} is complex differentiability at every point of $U$, i.e.\ holomorphy on the open set $U$.
The formula is exactly $\frac{d}{d\tau} = 2\pi i\,q\frac{d}{dq}$ applied termwise. \lean{adjointModule W} is $V$ as a module
over itself (Definition~\ref{l12:def:adjoint}); its conformal weight is $0$, whence the base point $0 - c/24$. The hypothesis
\lean{hne} excludes the zero VOA. The derivative statement uses $\lfourteenH$ only (\lean{hz}); the growth hypothesis is the
only analytic input. Same as textbook.
\end{dictionary}
\noindent\emph{Proof.} On $\lfourteenH_c$ ($c>0$), $|q^n|\le r^n$ with $r = e^{-2\pi c}<1$; subexponentiality at
$s = (1+r^{-1})/2$ (so $s>1$ and $sr<1$) gives $|\chi(h+n)\,2\pi n\,q^n|\le C\,2\pi n(sr)^n$, a summable dominant
independent of $z$. Mathlib's \lean{hasDerivAt_tsum_of_isPreconnected} differentiates the series termwise on the convex
(hence preconnected) open set $\lfourteenH_c$. The prefactor $e^{2\pi izh}$ is entire; product rule. For the adjoint
module, Chapter~\ref{ch:L15} supplies the growth datum from $C_2$-cofiniteness.
\begin{physics}
$\frac{1}{2\pi i}\frac{d}{d\tau}\chi = \tr(L_0 - c/24)q^{L_0-c/24}$: differentiating in the modular parameter inserts the
Hamiltonian, i.e.\ the zero mode of the stress tensor on the cylinder. This is the bridge between Zhu's recursion (stated
with the formal operator $D$) and differential equations in $\tau$.
\end{physics}
\begin{status}
\lean{differentiableOn_qEvalC}, \lean{deriv_qEvalC}: unconditional beyond the growth hypothesis. \lean{differentiableOn_char_adjoint}:
unconditional for $C_2$-cofinite $V$ with $\mathbf 1\neq0$. All \texttt{\#print axioms}-clean (standard three); no
\lean{VOA.pbw}.
\end{status}

\subsection{Algebra of evaluation}

\begin{definition}[Convolution with an $\N$-indexed series]\label{l14:def:qConv}
For $f:\N\to\C$ and a formal $q$-series $g$: $(f\star g)(e) = \sum_{n\in\N}f(n)\,g(e-n)$ (a \lean{finsum}: $0$ if the support
is infinite).
\end{definition}
\noindent File \leanpath{VOA/\allowbreak Modularity/\allowbreak ZhuRecursion.lean} (a Chapter~\ref{ch:L15} file).
\begin{leancode}
noncomputable def qConv (f : ℕ → ℂ) (g : FormalQSeries) : FormalQSeries :=
  fun e => ∑ᶠ n : ℕ, f n * g (e - (n : ℂ))
\end{leancode}
\begin{physics}
On coefficients, multiplying a torus amplitude by a $q$-series (an Eisenstein coefficient of a Ward identity) convolves the state counts level by level.
\end{physics}

\begin{theorem}[Convolution becomes multiplication]\label{l14:thm:qconv}
If $\chi(h+i) = 0$ for all integers $i<0$, and both $\sum|f(n)q^n|$ and $\sum|\chi(h+n)q^n|$ converge at $z$, then
$(f\star\chi)^{\C}_h(z) = \bigl(\sum_nf(n)e^{2\pi inz}\bigr)\cdot\chi^{\C}_h(z)$.
\end{theorem}
\noindent File \leanpath{VOA/\allowbreak Modularity/\allowbreak QEvalConv.lean}.
\begin{leancode}
theorem qEvalC_qConv (f : ℕ → ℂ) (χ : FormalQSeries) (h : ℂ) {z : ℂ}
    (hsupp : ∀ i : ℤ, i < 0 → χ (h + (i : ℂ)) = 0)
    (hf : Summable fun n : ℕ => ‖f n * Complex.exp (2 * (Real.pi : ℂ) * Complex.I * (n : ℂ) * z)‖)
    (hχ : Summable fun n : ℕ =>
      ‖χ (h + (n : ℂ)) * Complex.exp (2 * (Real.pi : ℂ) * Complex.I * (n : ℂ) * z)‖) :
    qEvalC (qConv f χ) h z
      = (∑' n : ℕ, f n * Complex.exp (2 * (Real.pi : ℂ) * Complex.I * (n : ℂ) * z))
          * qEvalC χ h z
\end{leancode}
\begin{dictionary}
\lean{hsupp} makes $(f\star\chi)(h+m) = \sum_{n\le m}f(n)\chi(h+m-n)$ a genuine finite sum (\lean{qConv_apply_of_support});
without it the \lean{finsum} could have infinite support and be junk. The two absolute-summability hypotheses are what
Mertens' theorem on Cauchy products needs (Mathlib \lean{tsum_mul_tsum_eq_tsum_sum_range_of_summable_norm}). Same as textbook.
\end{dictionary}
\noindent\emph{Proof.} Cauchy product of two absolutely convergent series.
\begin{physics}
Multiplying a torus amplitude by an Eisenstein series (as Zhu's recursion does) is, on coefficients, a convolution of state
counts.
\end{physics}
\begin{status}
Unconditional beyond hypotheses; standard axioms.
\end{status}

\begin{theorem}[Linearity, rebasing, leading coefficient]\label{l14:thm:qevallin}
(i) Evaluation is additive and commutes with finite sums where the series converge, and is homogeneous unconditionally.
(ii) $(D\chi)^{\C}_h = (2\pi i)^{-1}\frac{d}{dz}\chi^{\C}_h$ under the growth hypothesis.
(iii) If $\chi(d+s) = 0$ for $s<k$ and $\chi$ is summable at $\tau$ along $d$, then $\chi_d(\tau) = \chi_{d+k}(\tau)$.
(iv) If $\|\chi(e_0+n)\|$ is subexponential and $\chi_{e_0}(it) = 0$ for all $t>0$, then $\chi(e_0) = 0$.
\end{theorem}
\noindent File \leanpath{VOA/\allowbreak Modularity/\allowbreak QEvalLinear.lean}.
\begin{leancode}
theorem qEvalC_finsum {ι : Type*} (g : ι → FormalQSeries) (h : ℂ) {z : ℂ}
    (hfin : (Function.support g).Finite)
    (hsum : ∀ i, Summable (fun n : ℕ => qTerm (g i) h n z)) :
    qEvalC (∑ᶠ i, g i) h z = ∑ᶠ i, qEvalC (g i) h z

theorem qEvalC_qDeriv (χ : FormalQSeries) (h : ℂ)
    (hsub : Subexponential (fun n : ℕ => ‖χ (h + (n : ℂ))‖)) {z : ℂ} (hz : z ∈ imGT 0) :
    qEvalC (qDeriv χ) h z
      = (2 * (Real.pi : ℂ) * Complex.I)⁻¹ * _root_.deriv (qEvalC χ h) z
\end{leancode}
\noindent File \leanpath{VOA/\allowbreak Modularity/\allowbreak QEvalShift.lean}.
\begin{leancode}
theorem qEvalAt_shift (χ : FormalQSeries) (d : ℂ) (k : ℕ) (τ : UpperHalfPlane)
    (hsum : QSummableAt χ d τ) (hvan : ∀ s : ℕ, s < k → χ (d + (s : ℂ)) = 0) :
    qEvalAt χ d τ = qEvalAt χ (d + (k : ℂ)) τ
\end{leancode}
\noindent File \leanpath{VOA/\allowbreak Modularity/\allowbreak LeadingCoeff.lean}.
\begin{leancode}
theorem leading_coeff_eq_zero_of_qEvalAt_imAxis_eq_zero (χ : FormalQSeries) (e₀ : ℂ)
    (hsub : Subexponential (fun n : ℕ => ‖χ (e₀ + (n : ℂ))‖))
    (hzero : ∀ (t : ℝ) (ht : 0 < t), qEvalAt χ e₀ (imAxis t ht) = 0) :
    χ e₀ = 0
\end{leancode}
\noindent File \leanpath{VOA/\allowbreak Modularity/\allowbreak TraceNonvanish.lean} (a Chapter~\ref{ch:L15} file).
\begin{leancode}
noncomputable def imAxis (t : ℝ) (ht : 0 < t) : UpperHalfPlane :=
  ⟨Complex.I * (t : ℂ), by simp [ht]⟩
\end{leancode}
\begin{dictionary}
(i) \lean{qEvalC_finsum}: \lean{(Function.support g).Finite} says only finitely many $g_i$ are nonzero, so both
\lean{finsum}s are honest finite sums; summability of each term series is required because \lean{tsum} is additive only on
summable families. \lean{qEvalAt_finsetSum} (\leanpath{QEvalLinearity.lean}) is the same for \lean{qEvalAt} and a
\lean{Finset}. (iii) \lean{qEvalAt_shift}: the summability hypothesis is needed to split off finitely many terms. Because the
library's progressions are rigid, rebasing is what allows comparison of characters of modules whose weights differ by
integers. (iv) \lean{imAxis t} is $\tau = it$, where $q = e^{-2\pi t}\in(0,1)$. The conclusion is only about the
\emph{leading} coefficient; iterating it (with (iii)) gives uniqueness of all coefficients. The hypothesis \lean{hzero}
quantifies over the whole imaginary axis, which is strictly more than needed (a sequence $t\to\infty$ suffices), so the
theorem is slightly weaker than it could be but adequate.
\end{dictionary}
\noindent\emph{Proof of (iv).} On the imaginary axis $q = \rho\in(0,1)$ is real; $0 = \chi(e_0) + \sum_{n\ge1}\chi(e_0+n)\rho^n$
and the tail is $\le\rho\,C$ for $\rho\le\frac12$, with $C = \sum_n|\chi(e_0+n+1)|2^{-n}<\infty$. Let $\rho\to0$. No identity
theorem is needed.
\begin{physics}
(iv) is the elementary form of ``the low-temperature limit of the partition function sees the ground state'': as $t\to\infty$
only the lowest level survives. It is used to prove linear independence of characters of inequivalent sectors (Zhu's
Theorem~5.3.3, Chapter~\ref{ch:L15}).
\end{physics}
\begin{status}
All four unconditional beyond hypotheses; standard axioms.
\end{status}

\subsection{The \texorpdfstring{$T$}{T}-transformation, analytically}

\begin{definition}[Modular data; $T$-entries]\label{l14:def:moddata}
For a finite family $(M_i)_{i\in\iota}$ of irreducible $V$-modules with conformal weights $h_i$, a \emph{modular datum} is a
group homomorphism $\rho:\lfourteenSL\to\mathrm{Mat}_\iota(\C)$ (multiplicative) together with a matrix $\rho_S$ such that
$\rho(S) = \rho_S$ and $\rho(T) = \mathrm{diag}\bigl(e^{2\pi i(h_i-c/24)}\bigr)$.
\end{definition}
\noindent File \leanpath{VOA/\allowbreak Modularity/\allowbreak ZhuTheorem.lean} (a Chapter~\ref{ch:L15} file).
\begin{leancode}
noncomputable def modularTEntry (W : VertexOperatorAlgebra 𝒱) (h : ℂ) : ℂ :=
  Complex.exp (2 * (Real.pi : ℂ) * Complex.I * (h - W.centralCharge / 24))

noncomputable def modularRepT {ι : Type*} [DecidableEq ι] (W : VertexOperatorAlgebra 𝒱)
    (h : ι → ℂ) : ι → ι → ℂ :=
  fun i j => if i = j then modularTEntry W (h i) else 0

structure ModularData (W : VertexOperatorAlgebra 𝒱) (ι : Type*) [Fintype ι] [DecidableEq ι] where
  Mfam : ι → IrrModule W
  h : ι → ℂ
  hasWeight : ∀ i, HasConformalWeight (Mfam i).isMod (h i)
  rep : Matrix.SpecialLinearGroup (Fin 2) ℤ →* Matrix ι ι ℂ
  repS : ι → ι → ℂ
  rep_S : rep ModularGroup.S = Matrix.of repS
  rep_T : rep ModularGroup.T = Matrix.of (modularRepT W h)
\end{leancode}
\begin{dictionary}
\lean{IrrModule W} bundles an irreducible ordinary module (carrier, structure, irreducibility; \leanpath{VOA/\allowbreak Modules/\allowbreak Rationality.lean});
\lean{HasConformalWeight} is Definition~\ref{l12:def:lowest} ($M_h\neq0$ and all weights in $h+\N$). \lean{→*} is a monoid
homomorphism, i.e.\ a representation of $\lfourteenSL$ on $\C^\iota$. \lean{ModularData} contains \emph{no} link between
$\rho$ and the characters; that link is the proposition \lean{ZhuModularityAnalytic} (Chapter~\ref{ch:L15}). Note that
$\rho(T)$ is pinned but $\rho_S$ is free data. A \lean{ModularData} need not exist for an arbitrary weight vector: $\rho$
must be a homomorphism, so $\rho(T)$ must be compatible with the relations of $\lfourteenSL$ for some $\rho(S)$. The interesting objects are those produced by Zhu's theorem.
\end{dictionary}
\begin{physics}
A modular datum is the data of a rational CFT's modular tensor category at the level of the torus: the $T$-matrix (spins and central charge) and the $S$-matrix (exchange of space and time), forming a representation of the mapping class group $\lfourteenSL$ of the torus.
\end{physics}

\begin{theorem}[Zhu's law on the translation subgroup, analytically]\label{l14:thm:Tanalytic}
For every modular datum $D$, every $i$, every $k\in\N$ and $\tau'=\tau+k$:
$\chi_{M_i}(\tau') = \sum_j\rho(T^k)_{ij}\chi_{M_j}(\tau)$, where $\chi_{M_i}$ is the evaluated character based at
$h_i - c/24$. No convergence hypothesis is needed.
\end{theorem}
\noindent File \leanpath{VOA/\allowbreak Modularity/\allowbreak ConvergenceBridge.lean}.
\begin{leancode}
theorem zhu_modularity_T_pow_analytic {ι : Type*} [Fintype ι] [DecidableEq ι]
    {W : VertexOperatorAlgebra 𝒱} (D : ModularData W ι) (i : ι) (k : ℕ)
    {τ τ' : UpperHalfPlane} (hτ : (τ' : ℂ) = (τ : ℂ) + k) :
    qEvalAt (gradedDim (D.Mfam i).isMod) (D.h i - W.centralCharge / 24) τ'
      = ∑ j : ι, D.rep (ModularGroup.T ^ k) i j
          * qEvalAt (gradedDim (D.Mfam j).isMod) (D.h i - W.centralCharge / 24) τ
\end{leancode}
\begin{dictionary}
The right side evaluates \emph{every} $\chi_{M_j}$ at the base point of $M_i$; this is harmless because $\rho(T^k)$ is
diagonal and only $j=i$ survives. \textbf{Why no summability is needed:} $\tau\mapsto\tau+k$ multiplies each term
$\chi(h+n)q^{h+n}$ by the same phase $e^{2\pi ikh}$, so the identity holds termwise, and \lean{tsum} commutes with constant
multiples even for non-summable families (both sides junk-$0$ together). Hence this is a true statement also when the
character diverges; it is fully meaningful when it converges. The companion \lean{zhu_modularity_T_analytic} (same file, $k=1$)
is stated with an arbitrary base point $h$ for all $j$, which is also true by the same termwise argument.
\end{dictionary}
\noindent\emph{Proof.} \lean{qEvalAt_translate} identifies $\tau\mapsto\tau+1$ with the formal $q$-shift $q^e\mapsto
e^{2\pi ie}q^e$; the character's $T$-periodicity (Theorem~\ref{l12:thm:Tperiod}) says this is multiplication by
$e^{2\pi i(h_i-c/24)}$; and $\rho(T)^k$ is diagonal (\lean{Matrix.diagonal_pow}).
\begin{physics}
The Dehn twist $T$ acts on each sector by its spin phase $e^{2\pi i(h-c/24)}$; this is the trivially true half of modular
invariance, valid for any theory with a well-defined $L_0$.
\end{physics}
\begin{status}
Unconditional. Standard axioms.
\end{status}

%======================================================================
\section{Formal \texorpdfstring{$q$}{q}-series as a module over \texorpdfstring{$\C[[q]]$}{C[[q]]}}\label{l14:sec:qconvmodule}
%======================================================================

The convolution $\star$ of Definition~\ref{l14:def:qConv} is the only multiplication on \lean{FormalQSeries}. For it to make
``the $R$-span of the block values'' a meaningful object, it must be an \emph{action} of a ring.

\begin{definition}[Series based at $h$; the {$\C[[q]]$}-action]\label{l14:def:QSeriesAt}
$\mathcal Q_h$ is the subspace of formal $q$-series supported on $h+\N$ (i.e.\ $\chi(e) = 0$ unless $e = h+n$, $n\in\N$). For
$F = \sum F_nq^n\in\C[[q]]$ and $\chi\in\mathcal Q_h$, $F\cdot\chi := (n\mapsto F_n)\star\chi$.
\end{definition}
\noindent File \leanpath{VOA/\allowbreak Modularity/\allowbreak QConvModule.lean}.
\begin{leancode}
def QSeriesAt (h : ℂ) : Submodule ℂ FormalQSeries where
  carrier := {χ : FormalQSeries | ∀ e : ℂ, (∀ n : ℕ, e ≠ h + (n : ℂ)) → χ e = 0}

noncomputable def psSMul (F : PowerSeries ℂ) (χ : FormalQSeries) : FormalQSeries :=
  qConv (fun n => PowerSeries.coeff n F) χ

theorem psSMul_mul (F G : PowerSeries ℂ) {h : ℂ} {χ : FormalQSeries} (hχ : χ ∈ QSeriesAt h) :
    psSMul (F * G) χ = psSMul F (psSMul G χ)

noncomputable instance instModulePowerSeriesQSeriesAt (h : ℂ) :
    Module (PowerSeries ℂ) ↥(QSeriesAt h) where
\end{leancode}
\begin{dictionary}
Only the carrier of \lean{QSeriesAt} is quoted; closure under $+$, $0$, scalars is proved in the omitted fields.
\lean{↥(QSeriesAt h)} is the subspace regarded as a type. The \lean{instance} line is quoted without its six proof fields
(\lean{one_smul}, \lean{mul_smul}, \dots), which are discharged by \lean{psSMul_one}, \lean{psSMul_mul}, etc.
\textbf{Why the support condition is on all of $\C$.} The library notes (and we confirm) that requiring vanishing only at
$h+i$, $i<0$, is \emph{not} enough: it constrains $\chi$ only on the coset $h+\Z$, and at $e\notin h+\Z$ the \lean{finsum}
$\sum_n F_n\chi(e-n)$ may have infinite support (junk value $0$), destroying associativity. With the support on $h+\N$ every
convolution is a finite sum $\sum_{n\le N}F_n\chi(h+N-n)$. Faithful.
\end{dictionary}
\noindent\emph{Proof.} Unitality is immediate; associativity is the re-indexing of $\sum_{n+k\le N}$ by rows versus
antidiagonals (\lean{sum_triangle_swap}), i.e.\ associativity of the Cauchy product.
\begin{physics}
$\mathcal Q_h$ is the space of formal ``partition functions'' of sectors with ground-state energy $h$, and $\C[[q]]$ acts by
multiplication with $q$-series (e.g.\ Eisenstein series), exactly as in the coefficient form of a Ward identity.
\end{physics}
\begin{status}
Unconditional. \texttt{\#print axioms psSMul\_mul}: standard three.
\end{status}

\begin{theorem}[The Euler operator is a derivation over the action]\label{l14:thm:productrule}
For $\chi\in\mathcal Q_h$ and $F\in\C[[q]]$: $D(F\cdot\chi) = F\cdot D\chi + (\theta F)\cdot\chi$.
\end{theorem}
\noindent File \leanpath{VOA/\allowbreak Modularity/\allowbreak QDerivProductRule.lean}.
\begin{leancode}
theorem qDeriv_psSMul (F : PowerSeries ℂ) {h : ℂ} {χ : FormalQSeries} (hχ : χ ∈ QSeriesAt h) :
    qDeriv (psSMul F χ) = psSMul F (qDeriv χ) + psSMul (theta F) χ
\end{leancode}
\noindent\emph{Proof.} At $h+N$: $(h+N)F_n\chi(h+N-n) = F_n(h+N-n)\chi(h+N-n) + nF_n\chi(h+N-n)$.
\begin{dictionary}
Exactly the Leibniz rule for $q\frac d{dq}$. The docstring's point: the coefficient ring of Zhu's formal ODE need not be
closed under $\theta$ (which would require Ramanujan's identities, e.g.\ $\theta E_2 = (E_2^2-E_4)/12$), because the
downstream consumer only needs \emph{subexponential} coefficients, and $\mathcal S$ is $\theta$-stable
(Theorem~\ref{l14:thm:subexpclosure}). This is a genuine simplification, correctly argued.
\end{dictionary}
\begin{physics}
Differentiating a product of an Eisenstein coefficient and an amplitude: the $E_2$-type terms that appear are the
non-covariant pieces that the Serre derivative reorganises; for the growth estimate they are harmless.
\end{physics}
\begin{status}
Unconditional; standard axioms.
\end{status}

%======================================================================
\section{The \texorpdfstring{$q$}{q}-expansion bridge and the Noetherian coefficient ring}\label{l14:sec:qexp}
%======================================================================

\begin{definition}[Functions with analytic cusp function; the $q$-expansion map]\label{l14:def:cuspAnalytic}
A $1$-periodic $f:\lfourteenH\to\C$ has a \emph{cusp function} $\tilde f$ with $f(\tau) = \tilde f(e^{2\pi i\tau})$ (Mathlib
\lean{cuspFunction 1 f}, defined on the punctured disc and extended to $0$ by its limit). Let $\mathcal C$ be the set of $f$
whose cusp function is analytic at $0$; its \emph{$q$-expansion} is the Taylor series of $\tilde f$ at $0$ (Mathlib
\lean{qExpansion 1 f}).
\end{definition}
\noindent File \leanpath{VOA/\allowbreak Modularity/\allowbreak QExpansionBridge.lean}.
\begin{leancode}
noncomputable def CuspAnalytic : Subalgebra ℂ (UpperHalfPlane → ℂ) where
  carrier := {f : UpperHalfPlane → ℂ | AnalyticAt ℂ (cuspFunction 1 f) 0}

noncomputable def qExpansionAlgHom : CuspAnalytic →ₐ[ℂ] PowerSeries ℂ where
  toFun f := qExpansion 1 (f : UpperHalfPlane → ℂ)
  map_one' := qExpansion_one 1
  map_mul' f g := qExpansion_mul f.2 g.2
  map_zero' := qExpansion_zero 1
  map_add' f g := qExpansion_add f.2 g.2
\end{leancode}
\begin{dictionary}
\lean{AnalyticAt ℂ g 0}: $g$ is given by a convergent power series near $0$. Mathlib's \lean{qExpansion} is defined for any
function; it is a ring homomorphism only where the cusp functions are analytic (\lean{qExpansion_mul} needs both), which is why
the domain is restricted to $\mathcal C$. The quoted carrier omits closure proofs; \lean{qExpansionAlgHom} is quoted without
its \lean{commutes'} field. This is the statement ``$q$-expansion turns pointwise products into Cauchy products''.
\end{dictionary}
\begin{physics}
The partition function as a function of the torus versus its list of state counts: multiplication of partition functions
(stacking decoupled theories) corresponds to convolution of degeneracies.
\end{physics}

\begin{theorem}[Formal and analytic Eisenstein coefficients agree]\label{l14:thm:qexpzhuG}
For $j\ge2$: the $q$-expansion of $\tau\mapsto\mathcal G_{2j}(\tau)$ is $\sum_mg_j(m)q^m$; the $q$-expansion coefficients of
$E_{2j}$ are $g_j(m)/g_j(0)$; and $\sum g_j(m)q^m$ lies in the subalgebra $\C[\widehat{E}_4,\widehat{E}_6]\subseteq\C[[q]]$
generated by the $q$-expansions of $E_4,E_6$.
\end{theorem}
\noindent File \leanpath{VOA/\allowbreak Modularity/\allowbreak QExpansionBridge.lean}.
\begin{leancode}
theorem qExpansion_zhuGC {j : ℕ} (hj : 2 ≤ j) :
    qExpansion 1 (fun z : UpperHalfPlane => zhuGC j (z : ℂ)) = PowerSeries.mk (zhuG j)

noncomputable def E46Series : Subalgebra ℂ (PowerSeries ℂ) :=
  Algebra.adjoin ℂ ({qExpansion 1 (⇑ModularForm.E₄ : UpperHalfPlane → ℂ),
    qExpansion 1 (⇑ModularForm.E₆ : UpperHalfPlane → ℂ)} : Set (PowerSeries ℂ))

theorem mk_zhuG_mem_E46Series {j : ℕ} (hj : 2 ≤ j) :
    PowerSeries.mk (zhuG j) ∈ E46Series
\end{leancode}
\begin{dictionary}
This is the honest formal $\leftrightarrow$ analytic dictionary for the Eisenstein coefficients: the \emph{formal} sequence
\lean{zhuG j} (a definition) is proved to be the Fourier coefficient sequence of the \emph{holomorphic} function
\lean{zhuGC j}, by Mathlib's uniqueness of $q$-expansions (\lean{ModularFormClass.qExpansion_coeff_unique}). Then
generation (Theorem~\ref{l14:thm:levelone}) is transported to $\C[[q]]$ along \lean{qExpansionAlgHom}. Faithful.
\end{dictionary}
\noindent\emph{Proof.} $\mathcal G_{2j}$ has the convergent expansion $\sum g_j(m)q^m$ on $\lfourteenH$
(\lean{hasSum_zhuG}); a function with a convergent $q$-expansion has that expansion as its \lean{qExpansion}. For membership:
$\mathcal G_{2j}\in\lfourteenE\subseteq\mathcal C$ and \lean{qExpansionAlgHom} maps $\lfourteenE$ into \lean{E46Series}.
\begin{physics}
The Eisenstein coefficients of Zhu's recursion are, as formal series, polynomials in the $q$-expansions of $E_4$, $E_6$:
the formal recursion lives over a finitely generated ring.
\end{physics}
\begin{status}
Unconditional; standard axioms.
\end{status}

\begin{definition}[Zhu's quasimodular coefficient ring]\label{l14:def:GSeries}
$\lfourteenGS = \C[\mathcal G_2,\mathcal G_4,\mathcal G_6]\subseteq\C[[q]]$, the subalgebra generated by the three formal
series $\sum_mg_j(m)q^m$, $j=1,2,3$.
\end{definition}
\noindent File \leanpath{VOA/\allowbreak Modularity/\allowbreak QExpansionBridge.lean}.
\begin{leancode}
noncomputable def GSeries : Subalgebra ℂ (PowerSeries ℂ) :=
  Algebra.adjoin ℂ ({PowerSeries.mk (zhuG 1), PowerSeries.mk (zhuG 2),
    PowerSeries.mk (zhuG 3)} : Set (PowerSeries ℂ))

theorem mk_zhuG_mem_GSeries {j : ℕ} (hj : 1 ≤ j) :
    PowerSeries.mk (zhuG j) ∈ GSeries

theorem isNoetherianRing_GSeries : IsNoetherianRing GSeries
\end{leancode}
\noindent File \leanpath{VOA/\allowbreak Modularity/\allowbreak SubexpSeriesRing.lean}.
\begin{leancode}
theorem subexponential_coeff_of_mem_GSeries {F : PowerSeries ℂ} (hF : F ∈ GSeries) :
    Subexponential fun n : ℕ => ‖PowerSeries.coeff n F‖
\end{leancode}
\begin{dictionary}
\lean{IsNoetherianRing}: every ideal finitely generated (Hilbert's basis theorem, via \lean{Algebra.FiniteType}). The
\emph{formal} $\mathcal G_2$ is included as a generator because $\mathcal G_2$ is not a polynomial in $E_4$, $E_6$. All
$\mathcal G_{2j}$, $j\ge1$, lie in $\lfourteenGS$, every element of $\lfourteenGS$ has subexponential coefficients
($\lfourteenGS\le\mathcal S$ because $|g_j(n)|\le C+Cn^{2j}$). Faithful; these are exactly the properties used in
Chapter~\ref{ch:L15}.
\end{dictionary}
\begin{physics}
$\C[E_2,E_4,E_6]$ is the ring of quasimodular forms, the natural home of torus amplitudes with stress-tensor insertions
(Kaneko--Zagier). Its Noetherianity is what ensures that Zhu's infinitely many relations reduce to finitely many, i.e.\ that
torus characters satisfy a differential equation of finite order.
\end{physics}
\begin{status}
Unconditional; standard axioms (\lean{mk_zhuG_mem_GSeries}, \lean{isNoetherianRing_GSeries},
\lean{subexponential_coeff_of_mem_GSeries}).
\end{status}

%======================================================================
\section{The Serre derivative}\label{l14:sec:serre}
%======================================================================

\begin{definition}[Serre derivative]\label{l14:def:serre}
For $k\in\Z$ and $f:\C\to\C$: $\partial_kf = \frac{1}{2\pi i}f' + \frac{k}{(2\pi i)^2}G_2f$. On $\lfourteenH$, with
$G_2 = \frac{\pi^2}3E_2$, this is $\partial_kf = q\frac{df}{dq} - \frac{k}{12}E_2f$.
\end{definition}
\noindent File \leanpath{VOA/\allowbreak Modularity/\allowbreak SerreDerivative.lean}.
\begin{leancode}
noncomputable def serreDerivC (k : ℤ) (f : ℂ → ℂ) (z : ℂ) : ℂ :=
  (2 * (Real.pi : ℂ) * Complex.I)⁻¹ * _root_.deriv f z
    + (k : ℂ) * ((2 * (Real.pi : ℂ) * Complex.I)⁻¹) ^ 2 * (G2C z * f z)
\end{leancode}
\begin{dictionary}
$(2\pi i)^{-1}\frac{d}{dz} = q\frac{d}{dq}$. $(2\pi i)^{-2} = -\frac1{4\pi^2}$, so $\frac{k}{(2\pi i)^2}G_2 = -\frac{k}{4\pi^2}\cdot
\frac{\pi^2}{3}E_2 = -\frac k{12}E_2$: the classical $\vartheta_k = D - \frac k{12}E_2$ (Zagier; Kaneko--Zagier), and the same
as Mathlib's \lean{serreDerivative} (which is defined on \lean{ℍ → ℂ}). The docstring's claim that the coefficient
$k(2\pi i)^{-2}$ is \emph{forced} by equivariance is correct: by Theorems~\ref{l14:thm:slashderiv} and~\ref{l14:thm:G2slash}
the anomalies cancel for exactly this coefficient.
\end{dictionary}
\begin{physics}
The Serre derivative is the covariant derivative on the moduli space of tori; $E_2$ is the connection. It is the operation of
inserting the zero mode $L_{[0]}$ of the \emph{torus} stress tensor (Zhu's square-bracket $\tilde\omega = \omega -
\frac c{24}\mathbf 1$): this is the content of Zhu's $\omega$-block condition, where the $E_2$ term is precisely the $j=1$
Eisenstein coefficient (Theorem~\ref{l14:thm:zhuGCone}). In Mathur--Mukhi--Sen's classification of rational CFTs, characters
satisfy MLDEs $\bigl[\partial^n + \sum_r\phi_{2r}\partial^{n-r}\bigr]\chi = 0$ with $\partial^n = \partial_{2n-2}\circ\cdots
\circ\partial_0$ and $\phi_{2r}\in M_{2r}$; the second-order case with $\phi_4 = \mu E_4$ produces the famous list
$c = -\frac{22}5,1,2,\frac{14}5,4,\frac{26}5,6,7,\frac{38}5,8$ (Lee--Yang and the Deligne series of WZW models; quoted from
the literature, not formalised).
\end{physics}

\begin{theorem}[Equivariance of the Serre derivative]\label{l14:thm:serreslash}
For $k\in\Z$, $\gamma\in\lfourteenSL$, $z\in\lfourteenH$, and $f$ differentiable at $\gamma z$:
$(\partial_kf)|_{k+2}\gamma\,(z) = \partial_k(f|_k\gamma)(z)$. (No modularity of $f$ is required. This is equivariance only:
it does \emph{not} show that $\partial_k$ maps $M_k$ into $M_{k+2}$, which also needs holomorphy at the cusp; that is
Serre closure, Definition~\ref{l14:def:serreclosed}, proved separately in \leanpath{EisensteinSerreClosure.lean}.)
\end{theorem}
\noindent File \leanpath{VOA/\allowbreak Modularity/\allowbreak SerreDerivative.lean}.
\begin{leancode}
theorem serreDerivC_slashC (k : ℤ) (γ : SL(2, ℤ)) (f : ℂ → ℂ) {z : ℂ} (hz : z ∈ imGT 0)
    (hf : DifferentiableAt ℂ f (mobiusC γ z)) :
    slashC (k + 2) γ (serreDerivC k f) z = serreDerivC k (slashC k γ f) z
\end{leancode}
\begin{dictionary}
Pointwise in $\gamma$ and $z$; differentiability at one point. No modularity of $f$ is assumed: the statement is an identity of
operators, and it implies at once that $\partial_k$ maps weight-$k$ slash-invariant functions to weight-$(k+2)$ ones.
\textbf{Not in Mathlib}, which defines \lean{serreDerivative} but proves no equivariance; this is the library's own theorem.
It does \emph{not} prove holomorphy at the cusp of $\partial_kf$, which is why Serre closure (\lean{E46SerreClosed}) is
a separate statement, proved separately in \leanpath{EisensteinSerreClosure.lean} (\lean{e46SerreClosed}). Same as the textbook (Zagier, Prop.~5 for the case of modular $f$).
\end{dictionary}
\noindent\emph{Proof.} Split $\partial_kf$ into its two terms. The derivative term picks up the anomaly
$\frac{k}{2\pi i}\frac{c}{cz+d}(f|_k\gamma)$ (Theorem~\ref{l14:thm:slashderiv}); the $G_2$ term picks up
$\frac{k}{(2\pi i)^2}(-\frac{2\pi ic}{cz+d})(f|_k\gamma)$ (Theorem~\ref{l14:thm:G2slash}); they cancel.
\begin{physics}
This one computation is ``the only genuine computation in Zhu's $S$-law'' (library docstring), and the physical statement is
that the trace of the torus stress-tensor zero mode is a modular-covariant insertion: the conformal anomaly of $\partial_\tau$
and the anomaly of $E_2$ cancel exactly. It is what makes MLDEs $\lfourteenSL$-covariant and hence what forces the solution
space (the characters of a RCFT) to carry a representation of $\lfourteenSL$.
\end{physics}
\begin{status}
Unconditional; \texttt{\#print axioms serreDerivC\_slashC}: standard three. Every declaration in this section is pointwise in
$\gamma$, so it restricts verbatim to any subgroup (Section~\ref{l14:sec:theta}).
\end{status}

\begin{proposition}[Leibniz rule]\label{l14:prop:serremul}
$\partial_{k+l}(fg) = (\partial_kf)g + f(\partial_lg)$ at points where $f$, $g$ are differentiable
(\lean{serreDerivC_mul}, \leanpath{EisensteinRamanujan.lean}); $\partial_k$ is additive and $\C$-linear
(\lean{serreDerivC_add}, \lean{serreDerivC_smul}).
\end{proposition}
\noindent\emph{Proof.} The product rule; the $G_2$ terms add up because $(k+l) = k + l$. Unconditional, standard axioms.
\begin{physics}
Inserting the covariant torus derivative into a product of amplitudes obeys the Leibniz rule, with weights adding; this is what makes the coefficient algebra of MLDEs a differential ring.
\end{physics}

%======================================================================
\section{Growth of degeneracies from ordered spanning sets}\label{l14:sec:growth}
%======================================================================

\paragraph{Why growth estimates.} A character or one-point function is first constructed as a formal series. To speak of its
modular transformation one must evaluate it at every $\tau\in\lfourteenH$, and by Theorems~\ref{l14:thm:summable}
and~\ref{l14:thm:qsummable} this is possible exactly when the coefficients (the degeneracies $\dim M_{h+n}$, or the traces
$\tr_{M_{h+n}}o(v)$) are subexponential: for every $r>1$, bounded by $Cr^n$. Holomorphy and term-by-term differentiation
(Theorem~\ref{l14:thm:holo}) then come for free. Nothing in the VOA axioms gives such a bound directly; the library has two
routes to it, and this section and the next develop them. \emph{Route 1} (this section): a finite ordered spanning set,
available under $C_2$-cofiniteness, bounds $\dim M_{h+n}$ by the number of states of finitely many free bosons, which grows like
$e^{c\sqrt n}$, a subexponential rate. \emph{Route 2} (Section~\ref{l14:sec:ode}): if the one-point functions satisfy a formal
modular differential equation with subexponential coefficients, a majorant induction transfers the bound from the
coefficients to the solution. Route 2 is needed for the traces $\tr\,o(v)$ with $v\neq\mathbf 1$, which are not dimensions.

\begin{definition}[Coloured partition count]\label{l14:def:pcount}
For $g,n\in\N$, $p_g(\le n)$ is the number of functions $\mu:\{1..g\}\times\{1..n\}\to\{0..n\}$ with
$\sum_{(i,m)}\mu(i,m)\,m\le n$: the number of $g$-coloured partitions of integers $\le n$.
\end{definition}
\noindent File \leanpath{VOA/\allowbreak Modularity/\allowbreak GrowthBound.lean}.
\begin{leancode}
def partMultsLe (g n : ℕ) : Finset (Fin g × Fin n → Fin (n + 1)) :=
  Finset.univ.filter (fun μ => ∑ p, (μ p : ℕ) * (p.2 + 1) ≤ n)

def pcountLe (g n : ℕ) : ℕ := (partMultsLe g n).card
\end{leancode}
\begin{dictionary}
\lean{Fin n} $=\{0,\dots,n-1\}$; the part size is \lean{p.2 + 1} $\in\{1,\dots,n\}$, the multiplicity \lean{μ p}
$\in\{0,\dots,n\}$ (enough, since a multiplicity exceeding $n$ is impossible when the total is $\le n$). So the cardinality is
exactly the number of $g$-coloured partitions of size $\le n$.
\end{dictionary}
\begin{physics}
$p_g(\le n)$ counts the states of $g$ free chiral bosons up to level $n$; it is the reference count against which degeneracies of any finitely strongly generated theory are bounded.
\end{physics}

\begin{theorem}[The generating-function bound]\label{l14:thm:pcount}
For $0<t<1$: $p_g(\le n)\,t^n\le\exp\!\bigl(g\,t/(1-t)^2\bigr)$. Hence $n\mapsto p_g(\le n)$ is subexponential.
\end{theorem}
\noindent File \leanpath{VOA/\allowbreak Modularity/\allowbreak GrowthBound.lean}.
\begin{leancode}
theorem pcountLe_mul_pow_le (g n : ℕ) {t : ℝ} (ht0 : 0 < t) (ht1 : t < 1) :
    (pcountLe g n : ℝ) * t ^ n ≤ Real.exp (g * (t / (1 - t) ^ 2))

theorem pcountLe_subexponential (g : ℕ) : Subexponential (fun n => (pcountLe g n : ℝ))
\end{leancode}
\noindent\emph{Proof.} $p_g(\le n)t^n\le\sum_\mu t^{|\mu|}\le\prod_{k=1}^n(1-t^k)^{-g}\le\exp\bigl(g\sum_k\frac{t^k}{1-t^k}\bigr)$
$\le\exp(gt/(1-t)^2)$, using $(1-x)^{-1}\le e^{x/(1-x)}$. A uniform bound $a_nt^n\le K(t)$ for every $t\in(0,1)$ is equivalent
to subexponentiality (\lean{Subexponential.of_mul_pow_le}, with $t = r^{-1}$).
\begin{physics}
$\prod(1-q^k)^{-g}$ is the partition function of $g$ free chiral bosons (up to $q^{-g/24}$), and this is the elementary
(non-modular) half of Cardy's estimate. The true asymptotic $e^{\pi\sqrt{2gn/3}}$ needs the $S$-transformation of $\eta$;
here only subexponentiality is needed and the crude bound suffices.
\end{physics}
\begin{status}
Unconditional; standard axioms.
\end{status}

\begin{definition}[PBW and ordered spanning data]\label{l14:def:pbw}
For a module $N$ with base weight $h$: \emph{PBW spanning data} are finitely many generator weights $w_1,\dots,w_g\ge1$,
finitely many bottom weights $b_1,\dots,b_B\in\N$, and vectors $x_{\mu,j}$, indexed by multiplicity functions
$\mu:\{1..g\}\times\N\to\N$ of finite support and $j\le B$, with $x_{\mu,j}\in M_{h + b_j + \sum\mu(i,m)(w_i+m)}$, which span
$M$. \emph{Ordered spanning data} supply the $x_{\mu,j}$ as ordered monomials
$\prod_{(i,m)}(s_i)_{(-m-1)}^{\mu(i,m)}u_j$ in homogeneous generators $s_i\in V_{w_i}$ applied to bottom vectors $u_j$.
\end{definition}
\noindent File \leanpath{VOA/\allowbreak Modularity/\allowbreak GrowthBound.lean}.
\begin{leancode}
structure PBWSpanningData (N : OrdVModule W M) (h : ℂ) where
  g : ℕ
  wt : Fin g → ℕ
  wt_pos : ∀ i, 1 ≤ wt i
  b : ℕ
  botWt : Fin b → ℕ
  mono : (Fin g × ℕ →₀ ℕ) → Fin b → M
  mono_mem : ∀ μ j, mono μ j ∈ N.weight (h + (botWt j : ℂ) + (multWDeg wt μ : ℂ))
  spans : (⊤ : Submodule ℂ M) = Submodule.span ℂ {x | ∃ μ j, x = mono μ j}

def BuhlSpanning (W : VertexOperatorAlgebra 𝒱) : Prop :=
  ∀ (M : Type u) [AddCommGroup M] [Module ℂ M] (N : OrdVModule W M) (h : ℂ),
    HasConformalWeight N h → Nonempty (OrderedSpanningData N h)
\end{leancode}
\begin{dictionary}
\lean{Fin g × ℕ →₀ ℕ} is a finitely supported function $\mu$; \lean{multWDeg wt μ} $=\sum\mu(i,m)(w_i+m)$. \lean{N.weight α}
is the weight space $M_\alpha$. \lean{OrderedSpanningData} (same file) has the same shape with \lean{mono μ j} replaced by
\lean{orderedMonomial N gen μ (bot j)}, factors applied in lexicographic order of $(i,m)$; \lean{OrderedSpanningData.toPBW}
converts. \textbf{Weaker than Buhl / Gaberdiel--Neitzke.} No strict decrease of mode indices is required: repetitions are
allowed. This suffices for growth (coloured partitions allow repeats) but not for $C_m$-cofiniteness, as the library itself
warns. \textbf{Stronger than Buhl's theorem as a hypothesis.} \lean{BuhlSpanning W} asks for ordered spanning data for
\emph{every} module with a conformal weight, with a \emph{finite} bottom family; Buhl's theorem concerns modules generated by
one vector. Chapter~\ref{ch:L15} proves the adjoint case unconditionally and the module case under an explicit finite
lowest-weight generation hypothesis (\leanpath{BuhlSpanningC2.lean}); the ``Residual'' section of
\leanpath{GrowthBound.lean} is partly stale on this point (documentation only).
\end{dictionary}
\begin{physics}
An ordered (PBW) spanning set realises the module as a quotient of a Fock space of $g$ generalised oscillators built on finitely
many primaries; state counting is then bounded by that of $g$ free bosons.
\end{physics}

\begin{theorem}[PBW spanning data give the growth datum]\label{l14:thm:pbwgrowth}
Given PBW spanning data for $N$ at $h$: $\dim M_{h+n}\le B\cdot p_g(\le n)$ for all $n$, hence the graded dimensions are
subexponential along $h - c/24$. Under \lean{BuhlSpanning}, every module with a conformal weight has subexponential graded
dimensions.
\end{theorem}
\noindent File \leanpath{VOA/\allowbreak Modularity/\allowbreak GrowthBound.lean}.
\begin{leancode}
theorem finrank_weight_le_of_pbw (N : OrdVModule W M) (h : ℂ) (d : PBWSpanningData N h) (n : ℕ) :
    Module.finrank ℂ (N.weight (h + n)) ≤ d.b * pcountLe d.g n

theorem subexponentialCoeffs_of_pbw (N : OrdVModule W M) (h : ℂ) (d : PBWSpanningData N h) :
    SubexponentialCoeffs N (h - W.centralCharge / 24)

theorem subexponentialCoeffs_of_buhl {W : VertexOperatorAlgebra 𝒱} (hB : BuhlSpanning W)
    {M : Type u} [AddCommGroup M] [Module ℂ M] (N : OrdVModule W M) (h : ℂ)
    (hw : HasConformalWeight N h) : SubexponentialCoeffs N (h - W.centralCharge / 24)
\end{leancode}
\begin{dictionary}
The base point shifts from the $L_0$-weight $h$ to $h - c/24$ because \lean{gradedDim} indexes by $L_0 - c/24$. Since
\lean{Module.finrank} of an infinite-dimensional space is $0$ (Mathlib), the dimension bound could be vacuous only if weight
spaces were infinite-dimensional; they are finite-dimensional for every \lean{OrdVModule} (field \lean{finrank}). Faithful.
\end{dictionary}
\noindent\emph{Proof.} A spanning vector of weight exactly $h+n$ has $b_j + \sum\mu(i,m)(w_i+m) = n$, hence
$\sum\mu(i,m)(m+1)\le n$ (as $w_i\ge1$); restrict $\mu$ to the box $\{1..g\}\times\{0..n-1\}$, injective on such $\mu$. Then
the weight space is spanned by at most $B\cdot p_g(\le n)$ vectors. Theorem~\ref{l14:thm:pcount}.
\begin{physics}
Cardy-type growth of degeneracies, in the crude form ``at most free-boson growth'', from a finite set of strong generators.
\end{physics}
\begin{status}
Unconditional (beyond the data); standard axioms (\lean{subexponentialCoeffs_of_buhl} checked).
\end{status}

%======================================================================
\section{Formal modular differential equations and the indicial equation}\label{l14:sec:ode}
%======================================================================

Classically, Zhu proves convergence of one-point functions by showing that they satisfy a linear ODE in $q$ with a regular
singular point at $q=0$ and then invoking the Frobenius--Fuchs theory. The library replaces the whole theory by its
\emph{coefficient-estimate half}, which is elementary and is all that is needed for subexponential growth. Mathlib has no
regular-singular-point theory.

\begin{theorem}[The majorant theorem]\label{l14:thm:majorant}
Let $\chi_n, a_n, b_n\ge0$ with $a$, $b$ subexponential, let $p_n\to+\infty$, and suppose that for some $N_0$ and all $n\ge N_0$
\begin{equation*}
  p_n\,\chi_n \;\le\; \sum_{m<n}a_{n-m}\,\chi_m + b_n.
\end{equation*}
Then $\chi$ is subexponential.
\end{theorem}
\noindent File \leanpath{VOA/\allowbreak Modularity/\allowbreak Majorant.lean}.
\begin{leancode}
theorem subexponential_of_majorant_recursion {χ a b p : ℕ → ℝ}
    (hχ0 : ∀ n, 0 ≤ χ n) (ha0 : ∀ n, 0 ≤ a n) (hb0 : ∀ n, 0 ≤ b n)
    (ha : Subexponential a) (hb : Subexponential b)
    (hp : Tendsto p atTop atTop) (N₀ : ℕ)
    (hrec : ∀ n, N₀ ≤ n → p n * χ n ≤ (∑ m ∈ Finset.range n, a (n - m) * χ m) + b n) :
    Subexponential χ
\end{leancode}
\begin{dictionary}
\lean{Tendsto p atTop atTop}: $p_n\to+\infty$. All hypotheses are visible and are the intended ones. \textbf{Why $p_n\to\infty$
and not merely $p_n\ge c>0$} (the library's warning, which we verify): to prove $\chi_n\le Kr^n$ for a given $r$ one bounds the
right side by $(D\theta+1)Kr^n$ with $\theta = \frac{\sigma/r}{1-\sigma/r}$ and $D = D(\sigma)$ from $a_n\le D\sigma^n$
($1<\sigma<r$); one needs $p_n\ge D\theta+1$, and since $D(\sigma)$ may blow up as $\sigma\to1$ (equivalently as $r\to1$), a
fixed lower bound on $p$ does not serve every $r$. With $p_n\to\infty$ it suffices to start the induction late enough. The
finitely many initial values $n<N_0$ (where the recursion may fail, e.g.\ at resonances) are absorbed into $K$.
\end{dictionary}
\noindent\emph{Proof.} Fix $r>1$, $\sigma = (1+r)/2$, $t = \sigma/r$. Choose $N\ge N_0$ with $p_n\ge D\theta+1$ for $n\ge N$, and
$K\ge E$ (the constant of $b$ at $r$) with $\chi_n\le Kr^n$ for $n<N$. Strong induction: for $n\ge N$,
$p_n\chi_n\le\sum_{m<n}D\sigma^{n-m}Kr^m + Er^n\le(D\theta+1)Kr^n\le p_nKr^n$.
\begin{physics}
The majorant method is the elementary core of the statement ``a solution of a Fuchsian ODE in $q$ around a regular singular
point converges in the largest disc free of other singularities''. For characters the ODE is the MLDE of the theory, and the
conclusion is that the partition function converges on $|q|<1$: the theory has a finite free energy at every temperature.
\end{physics}
\begin{status}
Unconditional; standard axioms. The docstring's ``not yet connected to a VOA'' is stale: it is connected through
Theorem~\ref{l14:thm:formalODE} (documentation only).
\end{status}

\begin{definition}[Indicial polynomial, kernel and divisor]\label{l14:def:indicial}
For a monic order-$s$ equation $D^sF + \sum_{i<s}g_i(q)D^iF = 0$ with $g_i = \sum_kg_{i,k}q^k$, and an exponent $\lambda$:
$P(\lambda+n) = (\lambda+n)^s + \sum_{i<s}g_{i,0}(\lambda+n)^i$ (indicial polynomial along $\lambda+\N$),
$a_k = \sum_{i<s}|g_{i,k}|$ (kernel), and $p_n = |P(\lambda+n)|/(\max(|\lambda|,1)+n)^{s-1}$ (divisor).
\end{definition}
\noindent File \leanpath{VOA/\allowbreak Modularity/\allowbreak IndicialRecursion.lean}.
\begin{leancode}
noncomputable def indicialP (lam : ℂ) (s : ℕ) (g : ℕ → ℕ → ℂ) (n : ℕ) : ℂ :=
  (lam + (n : ℂ)) ^ s + ∑ i ∈ Finset.range s, g i 0 * (lam + (n : ℂ)) ^ i

noncomputable def indicialKernel (s : ℕ) (g : ℕ → ℕ → ℂ) (k : ℕ) : ℝ :=
  ∑ i ∈ Finset.range s, ‖g i k‖

noncomputable def indicialDivisor (lam : ℂ) (s : ℕ) (g : ℕ → ℕ → ℂ) (n : ℕ) : ℝ :=
  ‖indicialP lam s g n‖ / (max ‖lam‖ 1 + (n : ℝ)) ^ (s - 1)
\end{leancode}
\begin{dictionary}
\lean{g i k} is $g_{i,k}$, the $k$-th coefficient of the $i$-th ODE coefficient. \lean{indicialP} is the classical indicial
polynomial evaluated at $\lambda+n$; its roots are the Frobenius exponents. The divisor normalisation by $(B+n)^{s-1}$ is the
device that turns the $m$-dependent weights $(\lambda+m)^i$ into a single factor (see below).
\end{dictionary}
\begin{physics}
The roots of the indicial polynomial at $q=0$ are the possible leading exponents $h_i - c/24$ of the characters: in
Mathur--Mukhi--Sen, the indicial equation of the MLDE determines the central charge and the conformal weights of the RCFT
from the order $n$ and the number $\ell$ of zeros of the Wronskian, via $\sum_i(h_i - c/24) = \frac{n(n-1)}{12} - \frac\ell6$.
\end{physics}

\begin{theorem}[Indicial recursion $\Rightarrow$ subexponential; the divisor grows]\label{l14:thm:indicial}
(i) If $s\ge1$, all $\|g_{i,\cdot}\|$ are subexponential, $p_n\to\infty$, and
$P(\lambda+n)\chi_n = -\sum_{m<n}\sum_{i<s}g_{i,n-m}(\lambda+m)^i\chi_m$ for all $n$, then $|\chi_n|$ is subexponential.
(ii) For $s\ge1$, $p_n\to\infty$ always (monicity).
\end{theorem}
\noindent Files \leanpath{VOA/\allowbreak Modularity/\allowbreak IndicialRecursion.lean} and \leanpath{IndicialGrowth.lean}. File
\leanpath{VOA/\allowbreak Modularity/\allowbreak IndicialRecursion.lean}.
\begin{leancode}
theorem subexponential_of_indicial_recursion {χ : ℕ → ℂ} (lam : ℂ) {s : ℕ} (hs : 1 ≤ s)
    (g : ℕ → ℕ → ℂ) (hg : ∀ i, Subexponential (fun k => ‖g i k‖))
    (hp : Tendsto (indicialDivisor lam s g) atTop atTop)
    (hrec : ∀ n : ℕ, indicialP lam s g n * χ n
      = - ∑ m ∈ Finset.range n, ∑ i ∈ Finset.range s,
          g i (n - m) * (lam + (m : ℂ)) ^ i * χ m) :
    Subexponential (fun n => ‖χ n‖)
\end{leancode}
\noindent File \leanpath{VOA/\allowbreak Modularity/\allowbreak IndicialGrowth.lean}.
\begin{leancode}
theorem tendsto_indicialDivisor_atTop (lam : ℂ) {s : ℕ} (hs : 1 ≤ s) (g : ℕ → ℕ → ℂ) :
    Tendsto (indicialDivisor lam s g) atTop atTop
\end{leancode}
\begin{dictionary}
In (i) the hypothesis \lean{hp} is redundant given (ii); it is kept separate for modularity of the proof. \textbf{Resonances
are harmless.} If $P(\lambda+n) = 0$ for some $n$ (a second Frobenius exponent at integer distance, the logarithmic
case), the recursion at that $n$ says $0 = \cdots$ and gives no information on $\chi_n$; the majorant inequality
$p_n|\chi_n|\le\cdots$ then holds trivially because $p_n = 0$, and since $p_n\to\infty$ such $n$ are finitely many. So the theorem
is correct in the resonant case, where classical Frobenius theory needs logarithms. The docstring of
\leanpath{IndicialRecursion.lean} cites the growth lemma as \lean{tendsto_indicial_quotient}; its actual name is
\lean{tendsto_indicialDivisor_atTop} (documentation only, trivial).
\end{dictionary}
\noindent\emph{Proof.} (i) For $m<n$, $i<s$: $|\lambda+m|^i\le(B+n)^{s-1}$ with $B = \max(|\lambda|,1)$; divide the norm of the
recursion by $(B+n)^{s-1}$ and apply Theorem~\ref{l14:thm:majorant} with $b=0$. (ii) With $u = |\lambda+n|$ and
$C = \sum|g_{i,0}|$: $|P|\ge u^{s-1}(u-C)$ for $u\ge1$; for $n\ge2B+4C+4$, $u/(B+n)\ge\frac14$ and $u\ge n/2$, so
$p_n\ge4^{-(s-1)}(n/4)$.
\begin{physics}
The exponents of a Fuchsian equation are finitely many; beyond them each coefficient of the solution is determined from the
lower ones with a ``large denominator'' $P(\lambda+n)\sim n^s$, which is what keeps the degeneracies from growing
exponentially.
\end{physics}
\begin{status}
Both unconditional; standard axioms.
\end{status}

\begin{definition}[One-point functions on the torus; support and growth data]\label{l14:def:onepoint}
For a module $N$ and $v\in V$, $Z_N(v)$ is the formal one-point function $\sum_\alpha\tr_{M_\alpha}o(v)\,q^{\alpha-c/24}$
(Chapter~\ref{ch:L12}, Definition~\ref{l12:def:char}). The \emph{support datum} at base $h$ says $Z_N(v)(h+i) = 0$ for integers
$i<0$, for all $v$; the \emph{growth datum} says $n\mapsto|Z_N(v)(h+n)|$ is subexponential for all $v$.
\end{definition}
\noindent File \leanpath{VOA/\allowbreak Modularity/\allowbreak OnePointFunctions.lean} (a Chapter~\ref{ch:L15} file).
\begin{leancode}
noncomputable def torusOnePoint (N : OrdVModule W M) (v : 𝒱) : FormalQSeries := onePoint N v
\end{leancode}
\noindent File \leanpath{VOA/\allowbreak Modularity/\allowbreak AnalyticBlock.lean} (a Chapter~\ref{ch:L15} file).
\begin{leancode}
def OnePointSupported (N : OrdVModule W M) (h : ℂ) : Prop :=
  ∀ (v : 𝒱) (i : ℤ), i < 0 → torusOnePoint N v (h + (i : ℂ)) = 0

def OnePointSubexp (N : OrdVModule W M) (h : ℂ) : Prop :=
  ∀ v : 𝒱, Subexponential (fun n : ℕ => ‖torusOnePoint N v (h + (n : ℂ))‖)
\end{leancode}
\begin{dictionary}
\lean{onePoint N v e} $=\tr_{M_{e+c/24}}o(v)$ with $o(v)$ the zero mode (Chapter~\ref{ch:L12}). At $v=\mathbf 1$, $o(\mathbf
1) = \mathrm{id}$ and \lean{OnePointSubexp} specialises to \lean{SubexponentialCoeffs}. The support datum follows from a
conformal weight (\lean{onePointSupported_of_hasConformalWeight}).
\end{dictionary}
\begin{physics}
$Z_N(v)$ is the torus one-point function of the field $v$ in the sector $N$; the support datum says the sector has no states below its ground-state energy, the growth datum that its twisted partition functions converge.
\end{physics}

\begin{theorem}[A formal ODE gives the growth datum]\label{l14:thm:formalODE}
Let $N$ have the support datum at $h$. Suppose that for every $v$ there are $s\ge1$ and sequences $g_i$ ($i<s$) with
$|g_{i,\cdot}|$ subexponential such that, as formal $q$-series,
\begin{equation*}
  D^sZ_N(v) + \sum_{i<s}g_i\star D^iZ_N(v) = 0 .
\end{equation*}
Then $N$ has the growth datum at $h$. The same conclusion holds if the coefficients $g_i$ are elements of $\lfourteenGS =
\C[\mathcal G_2,\mathcal G_4,\mathcal G_6]$.
\end{theorem}
\noindent File \leanpath{VOA/\allowbreak Modularity/\allowbreak FormalODE.lean}.
\begin{leancode}
theorem onePointSubexp_of_formalODE (N : OrdVModule W M) (hbase : ℂ)
    (hsupp : OnePointSupported N hbase)
    (hODE : ∀ v : 𝒱, ∃ (s : ℕ) (g : ℕ → ℕ → ℂ), 1 ≤ s ∧
      (∀ i, Subexponential (fun k => ‖g i k‖)) ∧
      ∀ e : ℂ, (qDeriv^[s] (torusOnePoint N v)) e
        + ∑ i ∈ Finset.range s, qConv (g i) (qDeriv^[i] (torusOnePoint N v)) e = 0) :
    OnePointSubexp N hbase
\end{leancode}
\noindent File \leanpath{VOA/\allowbreak Modularity/\allowbreak SubexpSeriesRing.lean}.
\begin{leancode}
theorem onePointSubexp_of_GSeriesODE (N : OrdVModule W M) (hbase : ℂ)
    (hsupp : OnePointSupported N hbase)
    (hODE : ∀ v : 𝒱, ∃ (s : ℕ) (G : ℕ → PowerSeries ℂ), 1 ≤ s ∧
      (∀ i, G i ∈ GSeries) ∧
      ∀ e : ℂ, (qDeriv^[s] (torusOnePoint N v)) e
        + ∑ i ∈ Finset.range s,
            qConv (fun k => PowerSeries.coeff k (G i)) (qDeriv^[i] (torusOnePoint N v)) e = 0) :
    OnePointSubexp N hbase
\end{leancode}
\begin{dictionary}
\lean{qDeriv^[i]} is the $i$-th iterate $D^i$; since $D$ is diagonal, $D^i\chi(e) = e^i\chi(e)$ (\lean{qDeriv_iterate_apply}).
The equation is stated at \emph{every} $e\in\C$, monic of order $s\ge1$, with $s$ and $g$ allowed to depend on $v$ (as in
Zhu). The ODE is a \textbf{hypothesis}: producing it for a $C_2$-cofinite VOA is Zhu's Lemma 4.4.4, done (or not) in
Chapter~\ref{ch:L15}. These theorems are exactly ``ODE $\Rightarrow$ convergence''. They are formal-power-series statements:
no holomorphy is assumed or concluded; holomorphy then follows from Theorem~\ref{l14:thm:holo}.
\end{dictionary}
\noindent\emph{Proof.} Evaluating the equation at $\lambda+n$ and splitting off the $m=n$ term gives the indicial recursion
(\lean{indicialP_recursion_of_formalODE}); apply Theorem~\ref{l14:thm:indicial}. For $\lfourteenGS$ use
\lean{subexponential_coeff_of_mem_GSeries} (Section~\ref{l14:sec:qexp}).
\begin{physics}
``Characters satisfying a modular linear differential equation converge'', in its formal form: the MMS equation for the
one-point functions forces Cardy-compatible (subexponential) growth.
\end{physics}
\begin{status}
Unconditional given the ODE hypothesis; standard axioms (both).
\end{status}

\begin{theorem}[First-order systems]\label{l14:thm:system}
Let $z_1,\dots,z_g$ be sequences and $a_{kl}$ subexponential sequences with
$(h+N)z_k(N) = \sum_l\sum_{j\le N}a_{kl}(j)z_l(N-j)$ for all $k$, $N$. Then every $|z_k|$ is subexponential.
\end{theorem}
\noindent File \leanpath{VOA/\allowbreak Modularity/\allowbreak SystemGrowth.lean}.
\begin{leancode}
theorem subexponential_of_system {g : ℕ} {hbase : ℂ}
    (z : Fin g → ℕ → ℂ) (a : Fin g → Fin g → ℕ → ℂ)
    (ha : ∀ k l, Subexponential (fun j => ‖a k l j‖))
    (hrec : ∀ (k : Fin g) (N : ℕ),
      (hbase + (N : ℂ)) * z k N
        = ∑ l, ∑ j ∈ Finset.range (N + 1), a k l j * z l (N - j)) :
    ∀ k, Subexponential (fun N => ‖z k N‖)
\end{leancode}
\begin{dictionary}
This is the coefficient form of $q\frac{d}{dq}\vec z = A(q)\vec z$ (a first-order system) for a vector of series based at
\lean{hbase}. Note that the $j=0$ term of the right side involves $z_l(N)$ itself, which is why the proof moves it to the left
with $K = \sum|a_{kl}(0)|$. No eigenvalue or matrix-inverse condition is needed. Faithful.
\end{dictionary}
\noindent\emph{Proof.} Sum the norms over $k$: $Z_N = \sum_k|z_k(N)|$ satisfies $(|h+N| - K)Z_N\le\sum_{m<N}B_{N-m}Z_m$; apply
Theorem~\ref{l14:thm:majorant} with $p_N = |h+N| - K\to\infty$.
\begin{physics}
A vector-valued modular form (the vector of characters of all sectors) satisfying a first-order modular system
$D\vec\chi = A\vec\chi$; the system form is how block conditions in Zhu's argument are naturally organised.
\end{physics}
\begin{status}
Unconditional; standard axioms.
\end{status}

%======================================================================
\section{The cylinder-to-torus coordinate: formal identities}\label{l14:sec:coord}
%======================================================================

Zhu's theorem rests on the change of coordinate $w = e^z - 1$ from the plane to the cylinder (and thence the torus). The library
does all of its consequences as identities in $\C[[w]]$, without Laurent series or analysis. We collect the statements.

\subsection{The square-bracket modes}

\begin{definition}[Square-bracket coefficients and modes]\label{l14:def:sqbracket}
In $\C[[w]]$ let $\log(1+w) = \sum_{i\ge1}(-1)^{i+1}w^i/i$ and $u = \log(1+w)/w = \sum_i(-1)^iw^i/(i+1)$, a unit. For
$k\in\N$, $m,i\in\Z$ put
\begin{equation*}
  c(k,m,i) = [w^{i-m}]\bigl((1+w)^{k-1}u^m\bigr)\quad(i\ge m),\qquad c(k,m,i) = 0\quad(i<m).
\end{equation*}
For $a\in V$ and $m\in\Z$, $a[m]_k = \sum_{i\in\Z}c(k,m,i)\,a_{(i)}$ (a finite sum on each vector), and the torus Virasoro modes are
$L_{[n]} = \omega[n+1]_2 - \frac{c}{24}\mathbf 1[n+1]_0$.
\end{definition}
\noindent File \leanpath{VOA/\allowbreak Modularity/\allowbreak OnePointFunctions.lean} (a Chapter~\ref{ch:L15} file).
\begin{leancode}
noncomputable def logOnePlusXSeries : PowerSeries ℂ :=
  PowerSeries.mk (fun i => if i = 0 then 0 else (-1 : ℂ) ^ (i + 1) / (i : ℂ))

noncomputable def onePlusXUnit : (PowerSeries ℂ)ˣ := onePlusX_isUnit.unit

noncomputable def uSeries : PowerSeries ℂ := PowerSeries.mk (fun i => (-1 : ℂ) ^ i / (i + 1 : ℂ))

noncomputable def sqBracketCoeffZ (k : ℕ) (m i : ℤ) : ℂ :=
  if 0 ≤ i - m then
    PowerSeries.coeff (R := ℂ) (i - m).toNat
      ((↑(onePlusXUnit ^ ((k : ℤ) - 1)) : PowerSeries ℂ) * (↑(uUnit ^ m) : PowerSeries ℂ))
  else 0

noncomputable def vSqZ (v : 𝒱) (k : ℕ) (m : ℤ) : 𝒱 →ₗ[ℂ] 𝒱 where
  toFun x := ∑ᶠ i : ℤ, sqBracketCoeffZ k m i • mode (W.Y v) i x

noncomputable def Lsq (n : ℤ) : Module.End ℂ 𝒱 :=
  vSqZ W W.conformalVector 2 (n + 1)
    - (W.centralCharge / 24) • vSqZ W W.vacuum 0 (n + 1)
\end{leancode}
\begin{dictionary}
\lean{mode (W.Y v) i} is $v_{(i)}$ (Chapter~\ref{ch:01}). \lean{onePlusXUnit} is $1+w$ as a unit, so that $(1+w)^{k-1}$
makes sense for $k=0$ ($=(1+w)^{-1}$); \lean{uUnit} is $u$ as a unit, so $u^m$ makes sense for negative $m$. The \lean{finsum}
is finite on each $x$ (\lean{sqModeZ_support_finite}: $c(k,m,i) = 0$ for $i<m$, and $v_{(i)}x = 0$ for $i\gg0$). In Lean
\lean{vSqZ W v k m} takes the explicit VOA argument \lean{W} (a section variable). \textbf{Agreement with Zhu.} Zhu defines
$Y[a,z] = Y(a,e^z-1)e^{z\,\mathrm{wt}\,a}$ and $a[m] = \Res_zz^mY[a,z]$ for $a\in V_k$. Substituting $w = e^z-1$:
$z = \log(1+w)$, $dz = dw/(1+w)$, $e^{kz} = (1+w)^k$, so $a[m] = \Res_wY(a,w)(1+w)^{k-1}\log(1+w)^m = \sum_ia_{(i)}[w^i]
\bigl((1+w)^{k-1}w^mu^m\bigr) = \sum_ic(k,m,i)a_{(i)}$. So \lean{vSqZ W a k m} is Zhu's $a[m]$ \emph{provided $k = \mathrm{wt}\,a$};
the library passes the weight as a separate argument and imposes $a\in V_k$ where needed. $L_{[n]} = \tilde\omega[n+1]$ with
$\tilde\omega = \omega - \frac c{24}\mathbf 1$, as in Zhu.
\end{dictionary}
\begin{physics}
$a[m]$ are the modes of the field $a$ on the cylinder $z\in\C/2\pi i\Z$ (equivalently on the torus before imposing periodicity
in $\tau$); $L_{[0]}$ is the cylinder Hamiltonian. The shift $\tilde\omega = \omega - \frac c{24}$ is the Schwarzian
transformation of the stress tensor under $w = e^z - 1$, which produces the Casimir energy $-c/24$ (next theorem).
\end{physics}

\begin{theorem}[The Schwarzian of $e^z-1$]\label{l14:thm:schwarzian}
In $\C[[z]]$, with $\{f;z\} = f'''/f' - \frac32(f''/f')^2$: $\{e^z-1;z\} = -\frac12$, and $\{z;z\} = 0$.
\end{theorem}
\noindent File \leanpath{VOA/\allowbreak Modularity/\allowbreak Schwarzian.lean}.
\begin{leancode}
noncomputable def schwarzian (f : PowerSeries ℂ) : PowerSeries ℂ :=
  derivativeFun (derivativeFun (derivativeFun f)) * (derivativeFun f)⁻¹
    - (3 / 2 : ℂ) • ((derivativeFun (derivativeFun f) * (derivativeFun f)⁻¹) ^ 2)

noncomputable def expMinusOne : PowerSeries ℂ := PowerSeries.exp ℂ - 1

theorem schwarzian_expMinusOne :
    schwarzian expMinusOne = (-(1 / 2) : ℂ) • (1 : PowerSeries ℂ)
\end{leancode}
\begin{dictionary}
\lean{derivativeFun} is the formal derivative; \lean{⁻¹} on \lean{PowerSeries ℂ} is the inverse when the constant term is
nonzero (here $f' = e^z$, a unit), junk otherwise. So the definition is meaningful exactly for $f$ with $f'(0)\neq0$.
Faithful.
\end{dictionary}
\noindent\emph{Proof.} All derivatives of $e^z-1$ are $e^z$; $1 - \frac32 = -\frac12$.
\begin{physics}
The stress tensor transforms as $T(z) = f'(z)^2T(f(z)) + \frac{c}{12}\{f;z\}$. For the cylinder map the Schwarzian is the
constant $-\frac12$, so $T_{\text{cyl}} = \cdots - \frac c{24}$: the vacuum on a circle has energy $-c/24$. This is the origin
of the $q^{-c/24}$ in every character.
\end{physics}
\begin{status}
Unconditional; standard axioms. The docstring's ``Residual'' calls the operator $T_f$ of (C2) unbuilt; it is built in
\leanpath{ChangeCoord.lean} (documentation only). The two files also contradict each other on whether (C3) is a prerequisite of Zhu's theorem
(documentation only).
\end{status}

\begin{theorem}[$\exp\circ\log$ and the weight-matching identities]\label{l14:thm:explog}
In $\C[[w]]$: $\exp(\log(1+w)) = 1+w$; $\exp(n\log(1+w)) = \sum_r\binom nrw^r$ for $n\in\Z$. Consequently, for $k\ge1$,
$d\in\Z$, $r\in\N$:
\begin{align*}
 \text{(I1)}\quad&\sum_{m=0}^{r}\frac{d^m}{m!}c(k,m,r) = \binom{k-1+d}{r},\\
 \text{(I3)}\quad&\sum_{m\ge0}\frac{2d^{2m+1}}{(2m+1)!}c(k,2m+1,r) = \binom{k-1+d}{r} - \binom{k-1-d}{r},\\
 \text{(I2)}\quad&c(k,-1,r) - \binom{k-1}{r+1} - \tfrac12\binom{k-1}{r} = \sum_{m\ge0}g_{m+1}(0)\,c(k,2m+1,r).
\end{align*}
\end{theorem}
\noindent File \leanpath{VOA/\allowbreak Modularity/\allowbreak ExpLogBinomial.lean}.
\begin{leancode}
theorem exp_subst_log : (exp ℂ).subst (log ℂ) = 1 + X

theorem expLog_intCast_eq_binomSeries (n : ℤ) : expLog n = binomSeries n

theorem sum_pow_div_factorial_mul_sqBracketCoeff {k : ℕ} (hk : 1 ≤ k) (d : ℤ) (r : ℕ) :
    ∑ m ∈ Finset.range (r + 1), ((d : ℂ) ^ m / (m.factorial : ℂ)) * sqBracketCoeff k m r
      = cbz ((k : ℤ) - 1 + d) r

theorem finsum_odd_pow_mul_sqBracketCoeff {k : ℕ} (hk : 1 ≤ k) (d : ℤ) (r : ℕ) :
    ∑ᶠ m : ℕ, (2 * (d : ℂ) ^ (2 * m + 1) / ((2 * m + 1).factorial : ℂ))
        * sqBracketCoeff k (2 * m + 1) r
      = cbz ((k : ℤ) - 1 + d) r - cbz ((k : ℤ) - 1 - d) r

theorem sqBracketCoeffZ_neg_one_eq {k : ℕ} (hk : 1 ≤ k) (r : ℕ) :
    sqBracketCoeffZ k (-1) r - cbz ((k : ℤ) - 1) ((r : ℤ) + 1) - (1 / 2) * cbz ((k : ℤ) - 1) r
      = ∑ᶠ m : ℕ, zhuG (m + 1) 0 * sqBracketCoeff k (2 * m + 1) r
\end{leancode}
\begin{dictionary}
\lean{log ℂ} is Mathlib's $\log(1+X)$ series, \lean{exp ℂ} the exponential series, \lean{.subst} formal substitution
(meaningful as the substituted series has zero constant term). \lean{expLog c} $=\exp(c\log(1+X))$ and \lean{binomSeries n}
$=\sum_r\binom nrX^r$ (same file). \lean{cbz n r} is the generalised binomial $\binom nr$, $0$ for $r<0$ (Chapter~\ref{ch:01}).
\lean{sqBracketCoeff k m r} is the natural-index version $[w^r]((1+w)^{k-1}\log(1+w)^m)$, equal to $c(k,m,r)$. The
\lean{finsum}s are finite (terms vanish for $2m+1>r$). In (I2), $g_{m+1}(0) = -B_{2m+2}/(2m+2)!$, the cusp values of Zhu's
Eisenstein coefficients. Faithful.
\end{dictionary}
\noindent\emph{Proof.} Both $\exp(c\log(1+X))$ and $\sum\binom crX^r$ solve $(1+X)F' = cF$, $F(0) = 1$, whose solution is
unique coefficientwise (\lean{onePlusX_ode_unique}). (I1) is $[w^r]$ of $(1+w)^{k-1}e^{d\log(1+w)} = (1+w)^{k-1+d}$; (I3) is its
odd part; (I2) is $u = B(\log(1+w))$ with $B(z) = z/(e^z-1) = \sum B_nz^n/n!$, read through $\exp\circ\log = \mathrm{id}$.
\begin{physics}
These identities match the weights of the plane and cylinder mode expansions: they are the combinatorics behind Zhu's
recursion formula for torus one-point functions (Prop.~4.3.2), in which Bernoulli numbers enter as the cusp values of the
Eisenstein series.
\end{physics}
\begin{status}
Unconditional; standard axioms (all five).
\end{status}

\begin{theorem}[Torus Virasoro relations with the same central charge]\label{l14:thm:Lsq}
For all $m,n\in\Z$: $[L_{[m]},L_{[n]}] = (m-n)L_{[m+n]} + \delta_{m+n,0}\frac{m^3-m}{12}c$.
\end{theorem}
\noindent File \leanpath{VOA/\allowbreak Modularity/\allowbreak GelfandFuks.lean}.
\begin{leancode}
theorem Lsq_virasoro (W : VertexOperatorAlgebra 𝒱) (m n : ℤ) :
    ⁅Lsq W m, Lsq W n⁆
      = ((m : ℂ) - (n : ℂ)) • Lsq W (m + n)
        + (if m + n = 0 then (((m : ℂ) ^ 3 - (m : ℂ)) / 12 * W.centralCharge)
             • (1 : Module.End ℂ 𝒱) else 0)
\end{leancode}
\begin{dictionary}
\lean{⁅A, B⁆} is the commutator $AB - BA$ of endomorphisms. No hypotheses: true for every library VOA. Faithful (Zhu,
Theorem~4.2.1; Gannon (5.3.15)).
\end{dictionary}
\noindent\emph{Proof.} The commutator of $\omega[m+1]$ and $\omega[n+1]$ is computed from the plane Virasoro relations; the
central term is the moment $S(m,n) = \sum_ad[m,a]d[n,-a](a^3-a)$ of the change-of-basis coefficients, which equals
$(m^3-m)\delta_{m+n,0} - \frac{m-n}2\delta_{m+n,-2}$ (\lean{centralMomentClosed_holds}, proved by three formal Euler-operator
ODE computations). The anomalous $\delta_{m+n,-2}$ term is exactly cancelled by the $-\frac c{24}\mathbf 1[n+1]$ correction,
which is nonzero only at $n=-2$.
\begin{physics}
The Virasoro algebra on the cylinder has the same central charge as on the plane; the Schwarzian only shifts $L_0$ by $-c/24$.
The $\delta_{m+n,-2}$ cocycle is a coboundary, removed by redefining $\omega$ to $\tilde\omega$.
\end{physics}
\begin{status}
Unconditional; standard axioms.
\end{status}

\begin{definition}[The coordinate-change operator]\label{l14:def:changeCoord}
For coefficients $(a_k)_{k\ge1}$, $A_f = \sum_{k\ge1}a_kL_k$ is locally nilpotent on $V$, and $T_f = \exp(A_f)$ is a linear
automorphism of $V$.
\end{definition}
\noindent File \leanpath{VOA/\allowbreak Modularity/\allowbreak ChangeCoord.lean}.
\begin{leancode}
noncomputable def coordGen (W : VertexOperatorAlgebra 𝒱) (a : ℕ → ℂ) : 𝒱 →ₗ[ℂ] 𝒱 where
  toFun x := ∑ᶠ k : ℕ, a k • W.L ((k : ℤ) + 1) x

noncomputable def changeCoord (W : VertexOperatorAlgebra 𝒱) (a : ℕ → ℂ) : 𝒱 ≃ₗ[ℂ] 𝒱 :=
  expEquiv (coordGen_isLocNilp W a)
\end{leancode}
\begin{dictionary}
\lean{a k} plays $a_{k+1}$. Each $L_k$, $k\ge1$, lowers the $L_0$-grade by $k$, and the grading is bounded below, so $A_f$ is
locally nilpotent and the exponential is a finite sum (Chapter~\ref{ch:L04}'s \lean{expEquiv}). \lean{≃ₗ[ℂ]} is a linear
equivalence. Only the operator is built; the transformation law $T_fY(v,z)T_f^{-1} = \cdots$ (Huang's change of variables) is
\textbf{not} formalised (named residual (C3) in the docstring).
\end{dictionary}
\begin{physics}
$T_f$ implements a holomorphic reparametrisation $f(z) = z + a_1z^2+\cdots$ of the local coordinate at a puncture; on states
it acts through the positive Virasoro modes.
\end{physics}
\begin{status}
Definitions; standard axioms (\lean{changeCoord}). No theorem uses the transformation law.
\end{status}

%======================================================================
\section{The cusp: Zhu's relations at \texorpdfstring{$q=0$}{q=0}}\label{l14:sec:cusp}
%======================================================================

At $q=0$ the Eisenstein series become their constant terms, and Zhu's genus-one relations degenerate into the defining relations
of Zhu's algebra. This is the bridge from torus blocks to $A(V)$.

\begin{theorem}[A Bernoulli identity]\label{l14:thm:bernsq}
With $B(z) = z/(e^z-1)$: $B(z)^2e^z = B(z) - zB'(z)$, hence $[z^n](B^2e^z) = (1-n)B_n/n!$.
\end{theorem}
\noindent File \leanpath{VOA/\allowbreak Modularity/\allowbreak CuspDegeneration.lean}.
\begin{leancode}
theorem bernoulliSq_mul_exp :
    (bernoulliPowerSeries ℂ) ^ 2 * exp ℂ
      = bernoulliPowerSeries ℂ - X * derivative ℂ (bernoulliPowerSeries ℂ)
\end{leancode}
\noindent\emph{Proof.} Differentiate $B(e^z-1) = z$: $Be^z + (e^z-1)B' = 1$; multiply by $B$ and use $B(e^z-1) = z$ again.
\begin{dictionary}
\lean{bernoulliPowerSeries ℂ} is Mathlib's $\sum B_nz^n/n!$. Equivalent to the expansion $e^z/(e^z-1)^2 =
z^{-2} + \sum_{n\ge1}(2n-1)g_n(0)z^{2n-2}$, i.e.\ the Laurent expansion of $\wp$ on the cylinder.
\end{dictionary}
\begin{physics}
$e^z/(e^z-1)^2$ is the cylinder propagator (the two-point function of a free boson current); its Laurent coefficients are the
cusp values of the Eisenstein series.
\end{physics}
\begin{status}
Unconditional; standard axioms.
\end{status}

\begin{theorem}[DLM Lemma 9.2(c): the $q\to0$ limit of the {$[-2]$}-family]\label{l14:thm:negtwo}
For $k\ge1$ and $a,v\in V$:
\begin{equation*}
  a[-2]v + \sum_{n\ge2}(2n-1)g_n(0)\,a[2n-2]v \;=\; a\circ v + \tfrac1{12}a[0]v,
\end{equation*}
where all square-bracket modes are taken with weight parameter $k$ and $a\circ v = \sum_{i=0}^{k}\binom ki a_{(i-2)}v =
\Res_z\frac{(1+z)^k}{z^2}Y(a,z)v$ is Zhu's circle product.
\end{theorem}
\noindent File \leanpath{VOA/\allowbreak Modules/\allowbreak ZhuModule.lean} (a Chapter~\ref{ch:L12} file).
\begin{leancode}
noncomputable def Pres (W : VertexOperatorAlgebra 𝒱) (M : ℕ) (a b : 𝒱) (n : ℕ) : 𝒱 :=
  ∑ i ∈ Finset.range (M + n + 1), cbz ((M : ℤ) + n) (i : ℤ) • W.smul a ((i : ℤ) - 2 - n) b
\end{leancode}
\noindent File \leanpath{VOA/\allowbreak Modularity/\allowbreak CuspDegeneration.lean}.
\begin{leancode}
theorem negTwoFamily_cusp {k : ℕ} (hk : 1 ≤ k) (a v : 𝒱) :
    vSqZ W a k (-2) v
        + ∑ᶠ m : ℕ, (((2 * (m + 2) - 1 : ℕ) : ℂ) * zhuG (m + 2) 0)
            • vSqZ W a k (2 * ((m : ℤ) + 2) - 2) v
      = Pres W k a v 0 + (1 / 12 : ℂ) • vSqZ W a k 0 v
\end{leancode}
\begin{dictionary}
\lean{W.smul a j b} is $a_{(j)}b$; \lean{Pres W k a v 0} $=\sum_{i\le k}\binom ki a_{(i-2)}v$, Zhu's $a\circ v$ for
$a\in V_k$. The \lean{finsum} index $m$ is $n-2$. The $\frac1{12}$ is the missing $n=1$ term ($(2\cdot1-1)g_1(0) = -\frac1{12}$),
i.e.\ the $E_2$ anomaly. \textbf{Stronger than DLM:} no hypothesis $a\in V_k$; the identity holds for any $a$ once all modes use
the same weight parameter $k$ (it is a scalar identity of coefficients). Faithful.
\end{dictionary}
\noindent\emph{Proof.} All modes are finite combinations of $a_{(i)}v$ over a common window; coefficientwise, the identity is
$[w^{r}]$ of $(1+w)^{k-1}u^{-2}\cdot(B^2e^z)|_{z=\log(1+w)} = (1+w)^k$, i.e.\ Theorem~\ref{l14:thm:bernsq} pushed through
$z = \log(1+w)$ (\lean{sqBracketCoeffZ_neg_two_cusp}).
\begin{physics}
On a degenerate (infinitely long) torus, the torus Ward identity for an insertion of $a$ next to $v$ becomes the
operator-product relation defining $A(V)$: $a\circ v$ acts by zero on ground states. This is why Zhu's algebra controls the
characters: it is the $q\to0$ fibre of the sheaf of torus conformal blocks.
\end{physics}
\begin{status}
Unconditional; standard axioms.
\end{status}

\begin{definition}[Cusp relation spaces]\label{l14:def:cuspRel}
$R_\infty = O(V) + \mathrm{span}\{a[0]v : a\in V_k, k\ge1\}$; and $R'_\infty$ is the span of the two Zhu relation families with
the Eisenstein values replaced by their constant terms: all $a[0]v$, and all
$a[-2]v + \sum_{m<K}(2m+3)g_{m+2}(0)a[2m+2]v$ with $K$ beyond the support.
\end{definition}
\noindent File \leanpath{VOA/\allowbreak Modularity/\allowbreak CuspFunctionalCount.lean}.
\begin{leancode}
def cuspRelSpan : Submodule ℂ 𝒱 :=
  O W ⊔ Submodule.span ℂ
    {x : 𝒱 | ∃ (ka : ℕ) (a v : 𝒱), 1 ≤ ka ∧ a ∈ W.grading (ka : ℤ) ∧ x = vSqZ W a ka 0 v}
\end{leancode}
\noindent File \leanpath{VOA/\allowbreak Modularity/\allowbreak CuspLimit.lean}.
\begin{leancode}
def blockRelSpanAtCusp : Submodule ℂ 𝒱 :=
  Submodule.span ℂ
    ({x : 𝒱 | ∃ (ka : ℕ) (a v : 𝒱), 1 ≤ ka ∧ a ∈ W.grading (ka : ℤ) ∧ x = vSqZ W a ka 0 v} ∪
     {x : 𝒱 | ∃ (ka K : ℕ) (a v : 𝒱), 1 ≤ ka ∧ a ∈ W.grading (ka : ℤ) ∧
        (∀ m : ℕ, K ≤ m → vSqZ W a ka (2 * (m + 1) - 2 : ℤ) v = 0) ∧
        x = vSqZ W a ka (-2) v + ∑ m ∈ Finset.range K,
              (((2 * (m + 2) - 1 : ℕ) : ℂ) * zhuG (m + 2) 0)
                • vSqZ W a ka (2 * (m + 2) - 2 : ℤ) v})
\end{leancode}
\begin{dictionary}
\lean{O W} is Zhu's $O(V)$ (Chapter~\ref{ch:L12}, Section~\ref{l12:sec:zhu}); \lean{⊔} is the sum of subspaces. Only
homogeneous $a$ of positive weight $k$ enter (weight-$0$ elements are multiples of $\mathbf 1$ under
\lean{VacuumGradeZero}). The truncation $K$ in \lean{blockRelSpanAtCusp} is a witness that the tail vanishes, so the finite
sum is the full sum. These are the definitions used in DLM §9.
\end{dictionary}
\begin{physics}
These are the torus Ward identities on the infinitely long torus: the relations that ground-state traces must satisfy.
\end{physics}

\begin{theorem}[The cusp fibre is {$A(V)/[A(V),A(V)]$}]\label{l14:thm:cuspfibre}
Assume $V_0\subseteq\C\mathbf 1$. Then $R'_\infty = R_\infty$, and $V/R_\infty\cong A(V)/[A(V),A(V)]$; in particular
$\dim V/R'_\infty = \dim A(V)/[A(V),A(V)]$, which equals the number $n$ of Wedderburn factors when
$A(V)\cong\prod_{i\le n}\mathrm{Mat}_{d_i}(\C)$.
\end{theorem}
\noindent File \leanpath{VOA/\allowbreak Modularity/\allowbreak BlockZeroMode.lean} (a Chapter~\ref{ch:L15} file).
\begin{leancode}
def VacuumGradeZero : Prop := W.grading 0 ≤ Submodule.span ℂ ({W.vacuum} : Set 𝒱)
\end{leancode}
\noindent File \leanpath{VOA/\allowbreak Modularity/\allowbreak CuspLimit.lean}.
\begin{leancode}
theorem blockRelSpanAtCusp_eq_cuspRelSpan :
    blockRelSpanAtCusp W = cuspRelSpan W

theorem finrank_quot_blockRelSpanAtCusp_eq :
    letI := instRingA W
    letI := instAlgebraA W
    Module.finrank ℂ (𝒱 ⧸ blockRelSpanAtCusp W)
      = Module.finrank ℂ (A W ⧸ commSpan ℂ (A W))
\end{leancode}
\noindent File \leanpath{VOA/\allowbreak Modularity/\allowbreak CuspFunctionalCount.lean}.
\begin{leancode}
theorem finrank_quot_cuspRelSpan_eq_of_algEquiv_pi_matrix
    (hfd : letI := instRingA W; FiniteDimensional ℂ (A W))
    {n : ℕ} {d : Fin n → ℕ} (hd : ∀ i, NeZero (d i))
    (e : letI := instRingA W; letI := instAlgebraA W;
      A W ≃ₐ[ℂ] ∀ i : Fin n, Matrix (Fin (d i)) (Fin (d i)) ℂ) :
    finrank ℂ (𝒱 ⧸ cuspRelSpan W) = n
\end{leancode}
\noindent File \leanpath{VOA/\allowbreak Modularity/\allowbreak SemisimpleTrace.lean} (a Chapter~\ref{ch:L15} file).
\begin{leancode}
def commSpan : Submodule K A := Submodule.span K {x : A | ∃ a b : A, x = a * b - b * a}
\end{leancode}
\begin{dictionary}
\lean{letI := instRingA W} installs the ring structure on Zhu's algebra \lean{A W} (Theorem~\ref{l12:thm:zhuring});
\lean{commSpan ℂ (A W)} is $[A,A]$, the span of commutators; \lean{⧸} is the quotient space. \lean{Module.finrank} is $0$ on
infinite-dimensional spaces, so the equality of \lean{finrank}s is informative only when $A(V)$ is finite-dimensional (true
under $C_2$, Chapter~\ref{ch:L12}); the theorem is nonetheless correct as stated since the underlying linear equivalence
\lean{cuspRelSpanQuotEquiv} holds unconditionally. The Wedderburn form needs an algebra isomorphism
\lean{e} with a product of full matrix algebras. \textbf{Scope:} this is the $q=0$ fibre only; the analogous statement at
$z_0\in\lfourteenH$ is \emph{false} as an identity of relation spaces (\lean{blockRelSpan W z₀} moves with $z_0$, as the
docstring notes), and the dimension transfer is a residual of Chapter~\ref{ch:L15}.
\end{dictionary}
\begin{warning}
\textbf{Redundant hypothesis (moderate) --- removed from the cusp theorems.} The statements quoted above no longer carry
\lean{hzero}; the library dropped it after this finding. (Some other theorems of Chapter~\ref{ch:L15} still carry it.)
\lean{VacuumGradeZero W} holds for \emph{every} library VOA: it is the structure field
\lean{vacuum_grade : grading 0 = ℂ ∙ vacuum} weakened to $\le$, and the library itself proves this as
\lean{vacuumGradeZero_of_vacuum_grade} (in \leanpath{AKSl2Payoff.lean}, Chapter~\ref{ch:L16}). So every theorem of this
section carrying \lean{hzero} is in fact unconditional. Harmless, but a referee should not count it as an assumption.
\end{warning}
\noindent\emph{Proof.} $\subseteq$: Theorem~\ref{l14:thm:negtwo} turns each $[-2]$-relation into $a\circ v + \frac1{12}a[0]v$ with
$a\circ v\in O(V)$. $\supseteq$: $O(V)$ is spanned by $(L_{-1}+L_0)u\circ v$-type elements which are Zhu circle products
(\lean{zhuMul_L_base}); weight-$0$ $u$ contributes $\mathbf 1_{(-2)} = 0$. The quotient: in $A(V)$, $a[0]v$ maps to
$a*v - v*a$ (Zhu), so $V/R_\infty = A(V)/[A,A]$. For $\prod\mathrm{Mat}_{d_i}$, $[A,A]$ is the sum of trace-zero parts, of
codimension $n$.
\begin{physics}
At the cusp the torus conformal blocks are the symmetric (trace-like) linear functionals on $A(V)$; for a rational theory with
$A(V) = \bigoplus\mathrm{End}(M_{i,h})$ there is one for each sector, the trace over its ground states. This is the
infinite-cylinder version of ``number of torus blocks = number of primaries''.
\end{physics}
\begin{status}
No hypothesis on $V$; the Wedderburn form is conditional on the
decomposition \lean{e}. Standard axioms (all three). The companions \lean{exists_finset_span_sup_blockRelSpanAtCusp}
(\leanpath{CuspPoint.lean}: a spanning family of size $\dim A/[A,A]$ modulo $R'_\infty$) and
\lean{exists_finrank_cuspFunctionals_le} (functionals killing the relations; at most $n$ independent) are proved likewise.
\end{status}

%======================================================================
\section{Zhu's transformation law, the read-off of \texorpdfstring{$\rho(S)$}{rho(S)}, and characters}\label{l14:sec:rhoS}
%======================================================================

This section states, in the language built above, what ``the trace functions transform under $\lfourteenSL$'' means in the
library, and proves the parts that are pure modular-group theory. The deep inputs (that the trace functions span the torus
blocks) are Chapter~\ref{ch:L15}'s.

\begin{definition}[Trace functions; the law at $A$; trace spanning]\label{l14:def:realises}
For a family $(M_i)$ of irreducible modules with weights $h_i$, $v\in V$ and $\tau\in\lfourteenH$, the \emph{trace function}
is $S_i(v,\tau) = Z_{M_i}(v)$ evaluated at $\tau$ along $h_i - c/24$. With $j(A,\tau) = c\tau+d$, a matrix $\rho_A$
\emph{realises the law at $A$} if for every $n\in\N$, every $v$ of square-bracket weight $n$ ($L_{[0]}v = nv$), every $i$ and
$\tau$:
\begin{equation*}
  S_i(v,A\tau) = j(A,\tau)^n\sum_j(\rho_A)_{ij}S_j(v,\tau).
\end{equation*}
\emph{Trace spanning} means that every $A\in\lfourteenSL$ has such a $\rho_A$. The law \emph{on $\Gamma$} asks this only for
$A\in\Gamma$.
\end{definition}
\noindent File \leanpath{VOA/\allowbreak Modularity/\allowbreak ZhuTheorem.lean} (a Chapter~\ref{ch:L15} file).
\begin{leancode}
noncomputable def autFactor (A : Matrix.SpecialLinearGroup (Fin 2) ℤ) (τ : UpperHalfPlane) : ℂ :=
  ((A : Matrix (Fin 2) (Fin 2) ℤ) 1 0 : ℂ) * (τ : ℂ)
    + ((A : Matrix (Fin 2) (Fin 2) ℤ) 1 1 : ℂ)
\end{leancode}
\noindent File \leanpath{VOA/\allowbreak Modularity/\allowbreak BlockSpanning.lean} (a Chapter~\ref{ch:L15} file).
\begin{leancode}
noncomputable def traceFun (Mfam : ι → IrrModule W) (h : ι → ℂ) (i : ι) (v : 𝒱)
    (τ : UpperHalfPlane) : ℂ :=
  qEvalAt (torusOnePoint (Mfam i).isMod v) (h i - W.centralCharge / 24) τ

def TraceSpanning (Mfam : ι → IrrModule W) (h : ι → ℂ) : Prop :=
  ∀ A : SL(2, ℤ), ∃ ρA : ι → ι → ℂ, ∀ (n : ℕ) (v : 𝒱), v ∈ sqWeightSpace W (n : ℂ) →
    ∀ (i : ι) (τ : UpperHalfPlane),
      traceFun Mfam h i v (A • τ) = autFactor A τ ^ n * ∑ j, ρA i j * traceFun Mfam h j v τ

def RealisesLaw (Mfam : ι → IrrModule W) (h : ι → ℂ) (A : SL(2, ℤ)) (ρA : ι → ι → ℂ) : Prop :=
  ∀ (n : ℕ) (v : 𝒱), v ∈ sqWeightSpace W (n : ℂ) → ∀ (i : ι) (τ : UpperHalfPlane),
    traceFun Mfam h i v (A • τ) = autFactor A τ ^ n * ∑ j, ρA i j * traceFun Mfam h j v τ

def ZhuModularityAnalytic (D : ModularData W ι) : Prop :=
  ∀ (A : SL(2, ℤ)) (n : ℕ) (v : 𝒱), v ∈ sqWeightSpace W (n : ℂ) → ∀ (i : ι) (τ : UpperHalfPlane),
    traceFun D.Mfam D.h i v (A • τ)
      = autFactor A τ ^ n * ∑ j, D.rep A i j * traceFun D.Mfam D.h j v τ
\end{leancode}
\noindent File \leanpath{VOA/\allowbreak Modularity/\allowbreak SquareBracketGrading.lean} (a Chapter~\ref{ch:L15} file).
\begin{leancode}
noncomputable def sqWeightSpace (W : VertexOperatorAlgebra 𝒱) (k : ℂ) : Submodule ℂ 𝒱 :=
  LinearMap.ker (Lsq W 0 - k • (LinearMap.id : 𝒱 →ₗ[ℂ] 𝒱))
\end{leancode}
\begin{dictionary}
\lean{A • τ} is Mathlib's Möbius action; \lean{autFactor A τ} is $c\tau+d$. The law is Zhu's Theorem~5.3.2 / Gannon (5.3.18a)
with weights in the \emph{square-bracket} grading, which is correct: Zhu's one-point functions are modular of weight $n$ for
$v\in V_{[n]}$, not for $v\in V_n$. \lean{traceFun} is a \lean{qEvalAt}, i.e.\ junk $0$ where the series diverges; in all
theorems below the growth datum \lean{OnePointSubexp} is a hypothesis, so the values are honest. A family of trace functions
that are all identically $0$ would satisfy every law with any $\rho$; this is excluded where it matters by
\lean{TraceNonvanishing} / \lean{TraceIndependent}, which make $\rho_A$ unique (\lean{rho_unique}).
\end{dictionary}
\begin{physics}
$S_i(v,\tau)$ is the torus one-point function $\tr_{M_i}o(v)q^{L_0-c/24}$: the expectation value of the field $v$ inserted
on the torus, summed over states of the sector $M_i$ running around the time circle. The law says these amplitudes form a
finite-dimensional representation of the mapping class group of the torus, with $v$ contributing its conformal weight.
\end{physics}

\begin{theorem}[$\rho(S)$ from the $S$-slash spanning statement]\label{l14:thm:rhoS}
Assume that the $M_i$ have conformal weights $h_i$ and the growth datum, and that the $S$-slash of each trace block is a linear
combination of the trace blocks (\lean{SlashSpannedAtS}). Then trace spanning holds for every $A\in\lfourteenSL$. If moreover
the weights $h_i$ are pairwise incongruent mod $\Z$ (equivalently $T$-eigenvalues distinct) and the trace functions do not all
vanish, Zhu's law holds with a genuine representation $\rho$. (Of these hypotheses, the $S$-spanning statement is the deep
one, Zhu's Theorem~5.3.3; the growth datum is supplied under $C_2$ in Chapter~\ref{ch:L15}; nonvanishing follows from the
growth datum; distinctness of weights mod $\Z$ is a genuine restriction.)
\end{theorem}
\noindent File \leanpath{VOA/\allowbreak Modularity/\allowbreak RhoS.lean}.
\begin{leancode}
def SlashSpannedAtS (Mfam : ι → IrrModule W) (h : ι → ℂ) : Prop :=
  ∀ i : ι, ∃ c : ι → ℂ, ∀ (n : ℤ) (v : 𝒱), v ∈ sqWeightSpace W (n : ℂ) → ∀ z ∈ imGT 0,
    blockSlash W ModularGroup.S (traceBlock Mfam h i) v z
      = ∑ j, c j * traceBlock Mfam h j v z

theorem traceSpanning_of_slashSpannedAtS (hcw : ∀ i, HasConformalWeight (Mfam i).isMod (h i))
    (hsub : ∀ i, OnePointSubexp (Mfam i).isMod (h i - W.centralCharge / 24))
    (hS : SlashSpannedAtS Mfam h) : TraceSpanning Mfam h

theorem zhu_modularity_of_slashSpannedAtS (hcw : ∀ i, HasConformalWeight (Mfam i).isMod (h i))
    (hsub : ∀ i, OnePointSubexp (Mfam i).isMod (h i - W.centralCharge / 24))
    (hS : SlashSpannedAtS Mfam h)
    (hinj : Function.Injective (tEigen W h)) (hnv : TraceNonvanishing Mfam h) :
    ZhuModularityAnalytic
      (modularDataOf Mfam h (traceSpanning_of_slashSpannedAtS Mfam h hcw hsub hS)
        (traceIndependent_of_tEigen_injective Mfam h hcw hinj hnv) hcw)
\end{leancode}
\begin{dictionary}
\lean{traceBlock Mfam h i} is the trace function in the $\C$ picture (\lean{traceBlock_coe}: it agrees with \lean{traceFun} on
$\lfourteenH$); \lean{blockSlash W γ F v} $=$ \lean{slashC (sqWtOf W v) γ (F v)}, slashing each $v$ by its own square-bracket
weight (\leanpath{GradedBlock.lean}, Chapter~\ref{ch:L15}). \lean{tEigen W h i} $=e^{2\pi i(h_i-c/24)}$; \lean{TraceNonvanishing}
asks that each $S_i$ be nonzero somewhere. \lean{modularDataOf} assembles a \lean{ModularData} (Definition~\ref{l14:def:moddata})
from the matrices, and \lean{ZhuModularityAnalytic D} is the law of Definition~\ref{l14:def:realises} for every $A$ with
$\rho = D.\mathtt{rep}$. \textbf{Hypotheses.} \lean{hS} is the deep input (the spanning half of Zhu's Theorem~5.3.3, at $S$
only); \lean{hsub} the growth datum; \lean{hinj} an honest restriction (false for tensor products and many lattice VOAs, as
\leanpath{TEigenInjective.lean} records; Chapter~\ref{ch:L15} removes it via fibre independence). The conclusion is the
textbook law. Conditional, and the conditions are visible.
\end{dictionary}
\noindent\emph{Proof.} The law at $S$ follows from \lean{hS} by reading off coefficients (\lean{realisesLaw_of_slashSpanned}).
The law at $T$ is $T$-periodicity (Theorem~\ref{l14:thm:Tanalytic}). The set of $A$ admitting a law is closed under products
(the cocycle $j(AB,\tau) = j(A,B\tau)j(B,\tau)$, \lean{autFactor_mul}) and contains $1$; and $\lfourteenSL$ is generated by $S$, $T$
\emph{as a monoid}, because $S^{-1} = S^3$ and $T^{-1} = (ST)^5S$ (\lean{submonoid_closure_S_T_eq_top},
\leanpath{TraceSpanningS.lean}). Independence (from distinct $T$-eigenvalues) makes the $\rho_A$ unique, so $A\mapsto\rho_A$ is
multiplicative.
\begin{physics}
Modular invariance of a rational CFT is fixed by two pieces of data: the diagonal $T$ (spins) and the matrix $S$ (exchange of
space and time). Once one knows that the $S$-transformed amplitudes lie in the span of the old ones, the full
$\lfourteenSL$-representation follows from group theory. $\rho(S)$ is the CFT's modular $S$-matrix, input to the Verlinde
formula.
\end{physics}
\begin{status}
Conditional on the named hypotheses. \texttt{\#print axioms} of both: standard three; no \lean{VOA.pbw}. The library's
$S$-matrix is \lean{rhoS_of_blockSpanned} (same file), conditional likewise.
\end{status}

\begin{theorem}[$T$-eigenvalues are distinct iff the weights are distinct mod $\Z$]\label{l14:thm:teigen}
$e^{2\pi i(h_i-c/24)} = e^{2\pi i(h_j-c/24)}$ iff $h_i - h_j\in\Z$; hence \lean{tEigen W h} is injective iff the $h_i$ are pairwise
incongruent mod $\Z$.
\end{theorem}
\noindent File \leanpath{VOA/\allowbreak Modularity/\allowbreak TEigenInjective.lean}.
\begin{leancode}
theorem tEigen_injective_iff (h : ι → ℂ) :
    Function.Injective (tEigen W h) ↔ ∀ i j : ι, (∃ n : ℤ, h i - h j = (n : ℂ)) → i = j
\end{leancode}
\noindent\emph{Proof.} $e^x = e^y$ iff $x - y\in2\pi i\Z$ (Mathlib \lean{exp_eq_exp_iff_exists_int}).
\begin{physics}
Sectors with spins differing by an integer cannot be told apart by the Dehn twist; then independence of characters needs the
finer ground-state analysis of Theorem~\ref{l14:thm:qevallin}(iv).
\end{physics}
\begin{status}
Unconditional; standard axioms. Used by \lean{traceIndependent_of_weights_distinct_mod_int} (same file), which is conditional
on ordered spanning data and the distinctness.
\end{status}

\begin{definition}[Jacobi characters; Gannon §5.3.5]\label{l14:def:jacobi}
For $u,v\in V_1$ and a bilinear form $\beta$ on $V$:
$Z_M(\tau,u,v) = \sum_n\tr_{M_{h+n}}\bigl(e^{2\pi i(o(v)-\beta(u,v)/2)}e^{2\pi i\tau o(u)}\bigr)e^{2\pi i\tau(h+n-(c+12\beta(u,u))/24)}$,
operator exponentials taken blockwise on the finite-dimensional weight spaces; $\chi^J_M(\tau,v,w) = e^{2\pi iw}Z_M(\tau,0,v)$.
\end{definition}
\noindent File \leanpath{VOA/\allowbreak Modularity/\allowbreak JacobiCharacter.lean}.
\begin{leancode}
noncomputable def jacobiChar (N : OrdVModule W M) (h : ℂ) (β : 𝒱 →ₗ[ℂ] 𝒱 →ₗ[ℂ] ℂ) (v : 𝒱) (w : ℂ)
    (τ : UpperHalfPlane) : ℂ :=
  Complex.exp (twoPiI * w) * Zfun N h β 0 v τ

def JacobiTransformation (D : ModularData W ι) (β : 𝒱 →ₗ[ℂ] 𝒱 →ₗ[ℂ] ℂ) : Prop :=
  ∀ (A : SL(2, ℤ)) (v : 𝒱), v ∈ W.grading 1 → ∀ (w : ℂ) (i : ι) (τ : UpperHalfPlane),
    jacobiChar (D.Mfam i).isMod (D.h i) β ((autFactor A τ)⁻¹ • v)
        (w - ((A : Matrix (Fin 2) (Fin 2) ℤ) 1 0 : ℂ) * β v v / (2 * autFactor A τ)) (A • τ)
      = ∑ j, D.rep A i j * jacobiChar (D.Mfam j).isMod (D.h j) β v w τ

theorem jacobiTransformation_zero (D : ModularData W ι) (β : 𝒱 →ₗ[ℂ] 𝒱 →ₗ[ℂ] ℂ)
    (hJ : JacobiTransformation D β) (A : SL(2, ℤ)) (i : ι) (τ : UpperHalfPlane) :
    qEvalAt (gradedDim (D.Mfam i).isMod) (D.h i - W.centralCharge / 24) (A • τ)
      = ∑ j, D.rep A i j * qEvalAt (gradedDim (D.Mfam j).isMod) (D.h j - W.centralCharge / 24) τ
\end{leancode}
\begin{dictionary}
\lean{Zfun} (same file) is the displayed $Z_M$, a \lean{tsum} (junk if divergent); \lean{twoPiI} $=2\pi i$. \lean{β} is an
arbitrary bilinear form; the intended instance is an invariant form, but invariance is \emph{not} required by the statement.
\lean{JacobiTransformation} is the Jacobi-form law $\chi^J_i(A\tau,\frac v{c\tau+d},w-\frac{c\beta(v,v)}{2(c\tau+d)}) =
\sum\rho(A)_{ij}\chi^J_j(\tau,v,w)$ (Gannon (5.3.20); Miyamoto; Krauel--Mason), a named \lean{Prop}, \textbf{never proved}.
\lean{jacobiTransformation_zero} only checks that its $v=0$, $w=0$ instance is Zhu's law for characters (at $v=0$ the Jacobi
character is the character, \lean{jacobiChar_zero_zero}).
\end{dictionary}
\begin{physics}
$Z_M(\tau,u,v)$ is the partition function twisted by the flavour (current) charges of the conserved currents in $V_1$: a
torus with background gauge field. Its transformation as a Jacobi form encodes the anomaly (level) of the current algebra.
\end{physics}
\begin{status}
\lean{JacobiTransformation}: statement only (open). \lean{jacobiTransformation_zero}: proved, standard axioms, but it is a
consistency check, not a step toward the law.
\end{status}

%======================================================================
\section{Subgroups and the theta group}\label{l14:sec:theta}
%======================================================================

\begin{definition}[Slash relations on a subgroup; the theta group]\label{l14:def:theta}
For a family $F:V\to(\C\to\C)$ and $\gamma$: $G$ is the \emph{slash of $F$ by $\gamma$} if $G(v) = F(v)|_n\gamma$ on
$\lfourteenH$ for every integer $n$ and $v\in V_{[n]}$. $F$ is \emph{modular on $\Gamma$} if $F$ is its own slash by every
$\gamma\in\Gamma$. The \emph{theta group} is $\lfourteenGt = \langle S,T^2\rangle\le\lfourteenSL$.
\end{definition}
\noindent File \leanpath{VOA/\allowbreak Modularity/\allowbreak BlockSlash.lean} (a Chapter~\ref{ch:L15} file).
\begin{leancode}
def IsSlashOf (W : VertexOperatorAlgebra 𝒱) (γ : SL(2, ℤ)) (F G : 𝒱 → ℂ → ℂ) : Prop :=
  ∀ (n : ℤ) (v : 𝒱), v ∈ sqWeightSpace W (n : ℂ) → ∀ z ∈ imGT 0, G v z = slashC n γ (F v) z
\end{leancode}
\noindent File \leanpath{VOA/\allowbreak Modularity/\allowbreak SlashSubgroup.lean}.
\begin{leancode}
def ModularOn (W : VertexOperatorAlgebra 𝒱) (Γ : Subgroup SL(2, ℤ)) (F : 𝒱 → ℂ → ℂ) : Prop :=
  ∀ γ ∈ Γ, IsSlashOf W γ F F

def RealisesLawOn (Γ : Subgroup SL(2, ℤ)) : Prop :=
  ∀ A ∈ Γ, ∃ ρA : ι → ι → ℂ, RealisesLaw Mfam h A ρA

def thetaGroup : Subgroup SL(2, ℤ) :=
  Subgroup.closure ({ModularGroup.S, ModularGroup.T ^ 2} : Set SL(2, ℤ))

def ThetaGroupProper : Prop := ModularGroup.T ∉ thetaGroup
\end{leancode}
\begin{dictionary}
All in namespace \lean{VOA.SlashSub}. \lean{Subgroup.closure s} is the subgroup generated by $s$. The classical theta group is
$\{\gamma\equiv I\text{ or }S\bmod 2\}$, of index $3$; \lean{thetaGroup} is \emph{defined} by generators. The identification
with the congruence description, properness ($T\notin\lfourteenGt$) and index $3$ are proved in
\leanpath{ThetaGroup.lean} (Chapter~\ref{ch:L17}) and imported into \leanpath{GammaThetaStatement.lean} as
\lean{thetaGroup_eq}, \lean{thetaGroupProper}, \lean{thetaGroup_index}, \lean{gamma_two_le_thetaGroup}. So
\lean{ThetaGroupProper}, stated ``deliberately not proved'' in \leanpath{SlashSubgroup.lean}, is a theorem of the library (that
file's comment is stale, documentation only). \lean{IsSlashOf} only constrains $G$ on \emph{integer} square-bracket weights; for a bosonic VOA
every $V_{[n]}$ with $n\notin\N$ is $0$ (\lean{sqWeightSpace_eq_bot_of_not_natCast}), so nothing is lost there, but the layer
cannot express half-integral weights (the library proves this limitation explicitly: \lean{isSlashOf_of_eqOn_int},
\lean{blockSlash_apply_of_notInt}).
\end{dictionary}
\begin{physics}
$\lfourteenGt$ is the stabiliser of the Neveu--Schwarz spin structure on the torus (antiperiodic fermions in both directions):
$T$ changes the time-direction boundary condition of fermions (NS $\to$ $\widetilde{\mathrm{NS}}$) and is therefore not a
symmetry of the NS partition function, while $T^2$ and $S$ are. Fermionic theories, and bosonic theories with half-integral
weights, have characters modular only for $\lfourteenGt$.
\end{physics}

\begin{theorem}[$\lfourteenGt$ as a monoid, and bosonic $\lfourteenGt$-modularity]\label{l14:thm:thetamonoid}
(i) $\lfourteenGt$ is contained in the submonoid generated by $S$, $T^2$, $T^{-2}$. (ii) For any family with conformal weights, the
law at $S$ implies the law on $\lfourteenGt$. (iii) $[\lfourteenSL:\lfourteenGt] = 3$.
\end{theorem}
\noindent File \leanpath{VOA/\allowbreak Modularity/\allowbreak SlashSubgroup.lean}.
\begin{leancode}
theorem thetaGroup_le_submonoidClosure (A : SL(2, ℤ)) (hA : A ∈ thetaGroup) :
    A ∈ Submonoid.closure
      ({ModularGroup.S, ModularGroup.T ^ 2, (ModularGroup.T ^ 2)⁻¹} : Set SL(2, ℤ))

theorem realisesLawOn_thetaGroup_of_S (hcw : ∀ i, HasConformalWeight (Mfam i).isMod (h i))
    (hS : ∃ ρ : ι → ι → ℂ, RealisesLaw Mfam h ModularGroup.S ρ) :
    RealisesLawOn Mfam h thetaGroup
\end{leancode}
\noindent File \leanpath{VOA/\allowbreak Modularity/\allowbreak GammaThetaStatement.lean}.
\begin{leancode}
theorem thetaGroup_index : SlashSub.thetaGroup.index = 3 :=
  ThetaGroup.genSTsq_index_eq_three
\end{leancode}
\begin{dictionary}
(i) uses $S^{-1} = S^3$; $T^{-2}$ must be adjoined because $\lfourteenGt/\{\pm1\}\cong\Z/2*\Z$ is free on $\bar S$, $\bar T^2$ (as the
docstring says; not proved, and not needed). (ii) the laws at $T^{\pm2}$ come free from $T$-periodicity, since $\rho(T)$ is
diagonal and invertible; so for a \emph{bosonic} VOA $\lfourteenGt$-modularity costs exactly the $S$-law, i.e.\ restriction buys
nothing (correctly stated). (iii) a one-line import from Chapter~\ref{ch:L17}.
\end{dictionary}
\noindent\emph{Proof.} (i) Every element of the subgroup closure is a word in $S^{\pm1}$, $T^{\pm2}$. (ii) Realising matrices
form a submonoid (Theorem~\ref{l14:thm:rhoS}'s proof). (iii) Chapter~\ref{ch:L17}.
\begin{physics}
For a bosonic theory the NS-sector subtlety is absent, so modularity for $\lfourteenGt$ is weaker than, and implied by, modularity
for $\lfourteenSL$.
\end{physics}
\begin{status}
Unconditional ((i), (iii)); conditional on the $S$-law ((ii)). Standard axioms (all three).
\end{status}

\begin{theorem}[The mechanism: the $T$-phase of half-integral weights]\label{l14:thm:tphase}
For $m\in\Z$: $e^{2\pi i\,m/2} = (-1)^m$, which is $-1$ for odd $m$; and $e^{2\pi i\cdot2\cdot m/2} = 1$.
\end{theorem}
\noindent File \leanpath{VOA/\allowbreak Modularity/\allowbreak GammaThetaStatement.lean}.
\begin{leancode}
theorem tPhase_half (m : ℤ) :
    Complex.exp (2 * (Real.pi : ℂ) * Complex.I * ((m : ℂ) / 2)) = (-1 : ℂ) ^ m

theorem tSqPhase_half (m : ℤ) :
    Complex.exp (2 * (Real.pi : ℂ) * Complex.I * (2 * ((m : ℂ) / 2))) = 1
\end{leancode}
\begin{physics}
A state of weight $m/2$ picks up a sign under the Dehn twist $T$ when $m$ is odd, and nothing under $T^2$. This is the entire
reason the symmetry group drops to $\langle S,T^2\rangle$. The library's docstring stresses, correctly, that it is the
half-integral \emph{weight}, not the fermion parity, that matters: a $\Z$-graded super VOA has supertraces modular for all of
$\lfourteenSL$ (Dong--Zhao), while its plain traces are modular for $\Gamma_0(2)$-type subgroups (the other spin structures).
\end{physics}
\begin{status}
Unconditional arithmetic; standard axioms.
\end{status}

\begin{definition}[Half-integral automorphy factors; the super law; character systems]\label{l14:def:halfaut}
A \emph{half-integral automorphy factor} is a consistent choice of square root of $c\tau+d$, needed to give meaning to
$(c\tau+d)^{m/2}$; since a square root cannot be chosen multiplicatively, a multiplier $\mu$ records the failure of the cocycle
law (for the theta group it is a root of unity, the theta multiplier). Precisely, a half-integral automorphy factor on $\Gamma$
is a function $\sqrt{j}(A,\tau)$ with $\sqrt j(A,\tau)^2 = c\tau+d$, nowhere
zero, $\sqrt j(1,\tau) = 1$, and a $\tau$-independent multiplier $\mu(A,B)$ with $\sqrt j(AB,\tau) =
\mu(A,B)\sqrt j(A,B\tau)\sqrt j(B,\tau)$ for $A,B\in\Gamma$. A \emph{super character system} for a super VOA $W$ is a grading
$\mathrm{sqWt}$ of $V$ supported on $\frac12\Z$ with $\mathbf 1$ in weight $0$, and a family $\chi_i(v,\tau)$ with
$\chi_i(0,\tau) = 0$. The \emph{law on $\Gamma$} asks $\chi_i(v,A\tau) = \sqrt j(A,\tau)^m\sum_j\rho(A)_{ij}\chi_j(v,\tau)$ for
$v$ of weight $m/2$.
\end{definition}
\noindent File \leanpath{VOA/\allowbreak Modularity/\allowbreak SlashSubgroup.lean}.
\begin{leancode}
structure HalfAutFactor (Γ : Subgroup SL(2, ℤ)) where
  sq : SL(2, ℤ) → UpperHalfPlane → ℂ
  sq_sq : ∀ A ∈ Γ, ∀ τ : UpperHalfPlane, sq A τ ^ (2 : ℤ) = autFactor A τ
  sq_ne_zero : ∀ A ∈ Γ, ∀ τ : UpperHalfPlane, sq A τ ≠ 0
  mult : SL(2, ℤ) → SL(2, ℤ) → ℂ
  cocycle : ∀ A ∈ Γ, ∀ B ∈ Γ, ∀ τ : UpperHalfPlane,
    sq (A * B) τ = mult A B * (sq A (B • τ) * sq B τ)
  sq_one : ∀ τ : UpperHalfPlane, sq 1 τ = 1

def SuperRealisesLawOn (W : SuperVertexOperatorAlgebra 𝒱) {ι : Type*} [Fintype ι]
    (Γ : Subgroup SL(2, ℤ)) (J : HalfAutFactor Γ)
    (χ : ι → 𝒱 → UpperHalfPlane → ℂ) (ρ : SL(2, ℤ) → ι → ι → ℂ) : Prop :=
  ∀ A ∈ Γ, ∀ (m : ℤ) (v : 𝒱), v ∈ W.grading ((m : ℂ) / 2) → ∀ (i : ι) (τ : UpperHalfPlane),
    χ i v (A • τ) = J.sq A τ ^ m * ∑ j, ρ A i j * χ j v τ
\end{leancode}
\noindent File \leanpath{VOA/\allowbreak Modularity/\allowbreak GammaThetaStatement.lean}.
\begin{leancode}
structure SuperCharacterSystem (W : SuperVertexOperatorAlgebra 𝒱) (ι : Type v) [Fintype ι] where
  sqWt : ℂ → Submodule ℂ 𝒱
  sqWt_halfInt : ∀ k : ℂ, (∀ m : ℤ, ((m : ℂ) / 2) ≠ k) → sqWt k = ⊥
  vacuum_mem_sqWt : W.vacuum ∈ sqWt 0
  chi : ι → 𝒱 → UpperHalfPlane → ℂ
  chi_zero : ∀ (i : ι) (τ : UpperHalfPlane), chi i 0 τ = 0

def CompleteRationalFamily (_S : SuperCharacterSystem W ι) : Prop → Prop := fun P => P

def Conjecture_HalfWeight (S : SuperCharacterSystem W ι) (hFam : Prop) : Prop :=
  SuperC2Cofinite W → IsHalfIntegerGraded W → HasStrictHalfWeight W →
    CharacterIndependent S → CompleteRationalFamily S hFam → GammaThetaModularityHolds S
\end{leancode}
\begin{dictionary}
\lean{SuperVertexOperatorAlgebra} is Chapter~\ref{ch:L17}'s notion; \lean{W.grading} its $L_0$-grading (complex weights).
\lean{GammaThetaModularityHolds S} $=\exists J\,\exists\rho$, law of \lean{ModularLawOn} (the same shape as
\lean{SuperRealisesLawOn} but with \lean{S.sqWt} in place of \lean{W.grading}) on \lean{thetaGroup}. \lean{CharacterIndependent}
is linear independence of the \lean{chi i}; \lean{SuperC2Cofinite}, \lean{IsHalfIntegerGraded}, \lean{HasStrictHalfWeight} are the
evident predicates (super $C_2$; weights in $\frac12\Z$; some weight in $\frac12+\Z$ nonzero). \textbf{What is and is not tied
to $W$.} Nothing in \lean{SuperCharacterSystem} relates \lean{chi} to traces over $W$-modules, and \lean{sqWt} need not be
$W$'s square-bracket grading; these are explicitly abstracted because the library has no super module theory (the docstring
says so). \lean{SuperRealisesLawOn} uses $L_0$-weights rather than $L_{[0]}$-weights (also documented, and repaired by
\lean{ModularLawOn}). \lean{HalfAutFactor} is satisfiable: on $\langle T^2\rangle$ by \lean{translationHalfAutFactor}, and on all
of $\lfourteenGt$ by \lean{thetaHalfAutFactorExists} (\leanpath{ThetaHalfAutFactor.lean}, Chapter~\ref{ch:L17}, principal square
root); the docstring of \lean{ThetaHalfAutFactorExists} in \leanpath{GammaThetaStatement.lean} calling it ``OPEN (task T5)'' (an internal project task label) is
stale (documentation only).
\end{dictionary}
\begin{warning}
\textbf{Template statements (moderate).} \lean{CompleteRationalFamily S P} is \emph{definitionally} $P$, so \lean{Conjecture_HalfWeight
S hFam} takes an arbitrary proposition as its last hypothesis. With \lean{hFam := False} it is trivially true; with \lean{hFam :=
True} it asserts $\lfourteenGt$-modularity of an \emph{arbitrary} independent family \lean{chi} (unrelated to $W$), which is false
in general. So the ``conjecture'' has content only after the caller supplies, in \lean{hFam} and \lean{chi}, the module theory
the library lacks. Likewise \lean{SuperRealisesLawOn} (and \lean{IsGammaThetaModular}) is satisfied by the zero family for every
$\Gamma$, $J$, $\rho$. The library is candid that these are \lean{def}s, not theorems, and no theorem depends on them, but
the names suggest more. Checked in Lean against the library (the check is not part of this release):
\lean{completeRationalFamily_iff}, \lean{conjecture_halfWeight_false}, \lean{superRealisesLawOn_zero}; compiled for this chapter.
\end{warning}
\begin{physics}
The intended statement (Höhn; Dong--Zhao; van Ekeren) is that the NS-sector characters of a $C_2$-cofinite, $\frac12\Z$-graded
theory form a representation of the theta group, with the half-integral factor $(c\tau+d)^{m/2}$ coming from the
$\eta$-function multiplier (the theta multiplier $\theta(A\tau)/\theta(\tau)$).
\end{physics}

\begin{proposition}[Check for the template warning]\label{l14:prop:certtheta}
For every super character system $S$: $\mathtt{CompleteRationalFamily}\,S\,P\iff P$; $\mathtt{Conjecture\_HalfWeight}\,S\,
\mathrm{False}$ holds; and the zero family satisfies \lean{SuperRealisesLawOn} for every $\Gamma$, $J$, $\rho$.
\end{proposition}
\noindent\emph{Proof.} Definitional unfolding; and $0 = J^m\cdot\sum\rho\cdot0$.
\begin{status}
Checked in Lean against the library, compiled with \texttt{lake env lean -DmaxSynthPendingDepth=3}: no errors. The check is not
part of this release.
\end{status}
\begin{physics}
None: this is a logical check that the template statements have no physical content until the missing module theory is supplied.
\end{physics}

\begin{theorem}[Consistency: bosonic VOAs satisfy the $\lfourteenGt$ statements]\label{l14:thm:thetaboson}
(i) If $D$ is a modular datum satisfying Zhu's law, then the character system of $D$ (genuine square-bracket grading, genuine
trace functions) satisfies the half-integral law on $\lfourteenGt$ for every half-integral factor $J$. (ii) For a $C_2$-cofinite VOA
with semisimple $A(V)$, a finite irreducible family of the right cardinality with simple, pairwise non-isomorphic bottoms,
and the datum \lean{Pr}, the characters are $\lfourteenGt$-modular.
\end{theorem}
\noindent File \leanpath{VOA/\allowbreak Modularity/\allowbreak GammaThetaStatement.lean}.
\begin{leancode}
theorem modularLawOn_thetaGroup_of_zhuModularityAnalytic (D : ModularData W ι)
    (hD : ZhuModularityAnalytic D) (J : SlashSub.HalfAutFactor SlashSub.thetaGroup) :
    IsGammaThetaModular (ofModularData D) J (fun A i j => D.rep A i j)

theorem characterModularOn_thetaGroup_of_C2Cofinite [Finite (Φ W)]
    (Mfam : ι → IrrModule W) (hw : ι → ℂ)
    (hC2 : C2Cofinite W)
    (hss : letI := instRingA W; IsSemisimpleRing (A W))
    (Pr : BottomProjection (parabolicRepOfAMod W (regularAMod W)))
    (hcard : Fintype.card ι = Nat.card (Φ W))
    (hcw : ∀ i, HasConformalWeight (Mfam i).isMod (hw i))
    (hsep : ∀ j k : ι, j ≠ k → hw j = hw k →
      IsEmpty (AModEquiv W (lowestAMod (hcw j).2) (lowestAMod (hcw k).2))) :
    ∃ ρ : SL(2, ℤ) → ι → ι → ℂ,
      CharacterModularOn SlashSub.thetaGroup
        (fun i τ => traceFun Mfam hw i W.vacuum τ) ρ
\end{leancode}
\begin{dictionary}
(i) is a pure restriction: on even doubled weight $m=2n$, $J^{2n} = (c\tau+d)^n$ (\lean{HalfAutFactor.zpow_even}); on odd $m$
the bosonic weight space is $0$. (ii) chains (i)'s character form off Chapter~\ref{ch:L15}'s headline
\lean{zhuModularity_pinned_of_C2Cofinite}; \lean{Φ W} is the set of iso-classes of irreducibles, \lean{Pr : BottomProjection …}
is the datum discussed in Chapter~\ref{ch:L13} (supplied for every $V$ by \lean{voaPBWBottomProjection}, from the theorem \lean{VOA.pbw}; MISTAKES entry 61, now resolved),
the former redundant hypotheses \lean{hzero} and \lean{hsimp} have been removed from the current statement. \lean{CharacterModularOn Γ χ ρ} is $\chi_i(A\tau) =
\sum_j\rho(A)_{ij}\chi_j(\tau)$ for $A\in\Gamma$. As the library says, these buy nothing mathematically over $\lfourteenSL$-modularity;
they are type-checks of the statements. The hypotheses \lean{[Finite (Φ W)]}, \lean{hcard}, \lean{hcw}, \lean{hsep} and
\lean{Pr} of (ii) are all dischargeable: finiteness is the theorem \lean{finite_Φ_of_C2Cofinite_pbw} from $C_2$ alone, and
the canonical family of Chapter~\ref{ch:L15} (Theorem~\ref{l15:thm:canonical}) satisfies the family hypotheses. The
hypothesis \lean{hss} is semisimplicity of $A(V)$, not rationality; for a rational $V$ (\lean{AdmissiblyRational}) it is
supplied by the theorem \lean{isSemisimpleRing_}\allowbreak\lean{A_of_}\allowbreak\lean{admissiblyRational}
(Chapter~\ref{ch:L15}, Theorem~\ref{l15:thm:zhuss}).
\end{dictionary}
\begin{physics}
A bosonic theory trivially passes the NS-sector test.
\end{physics}
\begin{status}
(i) conditional on \lean{hD}; (ii) conditional on $C_2$, semisimplicity, the family hypotheses and \lean{Pr}. Both
\texttt{\#print axioms}-clean (standard three), \textbf{not} using \lean{VOA.pbw}; the datum \lean{Pr} is supplied for every $V$ by
\lean{voaPBWBottomProjection} (from the theorem \lean{VOA.pbw}; Chapter~\ref{ch:L15}'s \lean{_voaPBW} variants).
\end{status}

%======================================================================
\section{The analytic \texorpdfstring{$j$}{j}-invariant and the abstract Moonshine corollary}\label{l14:sec:moonshine}
%======================================================================

\begin{definition}[Analytic $j$]\label{l14:def:jinv}
$j(\tau) = E_4(\tau)^3/\Delta(\tau)$ on $\lfourteenH$. It is $\lfourteenSL$-invariant (\lean{jInvariant_smul}). The ``Hauptmodul
step'' \lean{eq_jInvariant_add_const} states: \textbf{if one is given} a weight-$0$ modular form $F$ (holomorphic, also at the cusp)
with $F = \chi - j$, then $\chi = j + c$. \textbf{The hypothesis is the Hauptmodul content}: it is assumed, not proved.
\end{definition}
\noindent File \leanpath{VOA/\allowbreak Modularity/\allowbreak JInvariantAnalytic.lean}.
\begin{leancode}
noncomputable def jInvariant (z : ℍ) : ℂ :=
  (ModularForm.E₄ z) ^ 3 / ModularForm.discriminant z

theorem jInvariant_smul (γ : SL(2, ℤ)) (z : ℍ) : jInvariant (γ • z) = jInvariant z

theorem eq_jInvariant_add_const (F : ModularForm 𝒮ℒ 0) (χ : ℍ → ℂ)
    (hF : ∀ z, F z = χ z - jInvariant z) : ∃ c : ℂ, ∀ z, χ z = jInvariant z + c
\end{leancode}
\begin{dictionary}
Division is honest because $\Delta$ has no zeros on $\lfourteenH$ (Mathlib \lean{discriminant_ne_zero}; \lean{jInvariant_eq_iff}).
Mathlib has $E_4$ and $\Delta$ but not $j$ (true at this Mathlib version), so this is new, though elementary.
\lean{discriminant_slash} first upgrades Mathlib's $S$- and $T$-invariance of $\Delta$ to all of $\lfourteenSL$ (the invariance
set is a subgroup containing the generators). \lean{eq_jInvariant_add_const}: \textbf{the hypothesis carries the whole
content.} It \emph{assumes} a holomorphic weight-$0$ form $F$, holomorphic at the cusp, with $F = \chi - j$; then $F$ is
constant (Mathlib \lean{levelOne_weight_zero_const}). Proving that $\chi - j$ is such a form (holomorphy of $\chi$ and
cancellation of the cusp poles) is exactly the Hauptmodul step, and it is \textbf{not} done anywhere in the library. The
textbook Hauptmodul theorem (Chapter~\ref{ch:05}, Theorem~\ref{ch05:thm:haupt}) is not formalised.
\end{dictionary}
\begin{physics}
$j$ is the unique (up to constant) modular-invariant function with a simple pole $q^{-1}$: the partition function of a
holomorphic $c=24$ CFT is $j + \text{const}$, the constant being $\dim V_1 - 744$.
\end{physics}
\begin{status}
\lean{jInvariant_smul}, \lean{eq_jInvariant_add_const}: unconditional, standard axioms. The second is nearly a restatement of
Mathlib's lemma.
\end{status}

\begin{definition}[Formal character of a VOA; the Moonshine hypotheses and identity]\label{l14:def:moonchar}
$\mathrm{ch}(V) = \sum_n(\dim V_n)q^n\in\C[[q]]$. The \emph{Moonshine hypotheses} are $c=24$, $V_1 = 0$, $V$ holomorphic and
$C_2$-cofinite; the \emph{Moonshine character identity} is $\mathrm{ch}(V) = qJ$.
\end{definition}
\noindent File \leanpath{VOA/\allowbreak Modularity/\allowbreak MoonshineCorollary.lean}.
\begin{leancode}
noncomputable def charSeries (W : VertexOperatorAlgebra 𝒱) : PowerSeries ℂ :=
  PowerSeries.mk (fun n => (Module.finrank ℂ (W.grading n) : ℂ))

structure MoonshineHypotheses (W : VertexOperatorAlgebra 𝒱) : Prop where
  centralCharge_eq : W.centralCharge = 24
  grading_one_eq_bot : W.grading 1 = ⊥
  holomorphic : Holomorphic W
  c2 : C2Cofinite W

def MoonshineCharacter (W : VertexOperatorAlgebra 𝒱) : Prop := charSeries W = JShift
\end{leancode}
\begin{dictionary}
\lean{JShift} $=qJ$ (Chapter~\ref{ch:05}, Definition~\ref{ch05:def:jshift}); \lean{Holomorphic} is Chapter~\ref{ch:L12}'s
Definition~\ref{l12:def:gannon}. \lean{MoonshineCharacter} is a \lean{Prop}; \textbf{the library never proves it for any $V$},
and \lean{MoonshineHypotheses} is not used by any theorem. The library's actual $\chi_{\Vnat} = J$ (Chapter~\ref{ch:05}) is stated
for the concrete $\Vnat$ and does not mention these definitions.
\end{dictionary}
\begin{physics}
$\mathrm{ch}(V)$ is $q\cdot Z(\tau)$, the torus partition function with the Casimir factor $q^{-1}$ removed (Chapter~\ref{ch:05}).
\end{physics}
\begin{status}
Definitions. The small lemmas \lean{moonshineCharacter_coeff_zero_iff} ($\Rightarrow\dim V_0 = 1$) and
\lean{coeff_one_charSeries_of_grading_one_eq_bot} are proved, standard axioms.
\end{status}

\begin{theorem}[Zhu's theorem forces invariance of a $c=24$ weight-$0$ character]\label{l14:thm:moonsl2}
Let $V$ have $c = 24$ and an irreducible module $M$ of conformal weight $0$ (for instance $M = V$). If trace spanning and trace
independence hold for the one-element family $\{M\}$ and the character of $M$ is not identically zero, then $\rho$ is trivial and
the character is $\lfourteenSL$-invariant. Under $C_2$-cofiniteness, the growth datum and $\dim(\text{analytic $S$-slash
ambient})\le1$, the same conclusion holds without the spanning/independence hypotheses. (For a one-element family,
independence is automatic once the character is nonzero; spanning is the real hypothesis of the first form, and the rank bound
replaces it in the second.) Finally, \textbf{if one is given} a weight-$0$ modular form equal to $\chi_M - j$, then $\chi_M = j +
\text{const}$; \textbf{this last hypothesis is the Hauptmodul content and is assumed, not proved}, anywhere in the library.
\end{theorem}
\noindent File \leanpath{VOA/\allowbreak Modularity/\allowbreak MoonshineCorollary.lean}.
\begin{leancode}
theorem char_SL2Z_invariant_of_spanning (V : IrrModule W)
    (hsp : TraceSpanning (singleFam V) (fun _ => (0 : ℂ)))
    (hind : TraceIndependent (singleFam V) (fun _ => (0 : ℂ)))
    (hcw : HasConformalWeight V.isMod 0) (hc : W.centralCharge = 24)
    (hne : ∃ τ : UpperHalfPlane, traceFun (singleFam V) (fun _ => (0 : ℂ)) () W.vacuum τ ≠ 0)
    (A : SL(2, ℤ)) (τ : UpperHalfPlane) :
    traceFun (singleFam V) (fun _ => (0 : ℂ)) () W.vacuum (A • τ)
      = traceFun (singleFam V) (fun _ => (0 : ℂ)) () W.vacuum τ
\end{leancode}
\noindent File \leanpath{VOA/\allowbreak Modularity/\allowbreak MoonshineFromBlocks.lean}.
\begin{leancode}
theorem char_SL2Z_invariant_of_analyticRank (V : IrrModule W) (hC2 : C2Cofinite W) {z₀ : ℂ}
    (hz₀ : z₀ ∈ imGT 0) (hcw : HasConformalWeight V.isMod 0)
    (hsub : ∀ i : Unit, OnePointSubexp (singleFam V i).isMod
      ((fun _ : Unit => (0 : ℂ)) i - W.centralCharge / 24))
    (hle : Module.finrank ℂ (analyticSlashAmbient (singleFam V) (fun _ => (0 : ℂ))) ≤ 1)
    (hc : W.centralCharge = 24)
    (A : SL(2, ℤ)) (τ : UpperHalfPlane) :
    traceFun (singleFam V) (fun _ => (0 : ℂ)) () W.vacuum (A • τ)
      = traceFun (singleFam V) (fun _ => (0 : ℂ)) () W.vacuum τ
\end{leancode}
\noindent File \leanpath{VOA/\allowbreak Modularity/\allowbreak MoonshineHauptmodul.lean}.
\begin{leancode}
theorem char_eq_jInvariant_add_const (V : IrrModule W) (F : ModularForm 𝒮ℒ 0)
    (hF : ∀ z : UpperHalfPlane,
      F z = traceFun (singleFam V) (fun _ => (0 : ℂ)) () W.vacuum z - jInvariant z) :
    ∃ c : ℂ, ∀ τ : UpperHalfPlane,
      traceFun (singleFam V) (fun _ => (0 : ℂ)) () W.vacuum τ = jInvariant τ + c
\end{leancode}
\begin{dictionary}
\lean{singleFam V : Unit → IrrModule W} is the one-element family; \lean{()} its index. \lean{traceFun … () W.vacuum} is the
character $\chi_M(\tau) = \tr_Mq^{L_0-1}$ (base point $0 - 24/24$, $o(\mathbf 1) = \mathrm{id}$). \textbf{Holomorphy of $V$ is not
assumed}: the theorem is about one module of weight $0$, and uses $c = 24$ only to make $\rho(T) = e^{2\pi i(0-1)} = 1$. Then
$\rho(-1) = 1$ (law at $-1$ on $\mathbf 1$, using \lean{hne}), $\rho(S)^2 = \rho(-1) = 1$ and $\rho(S)^3 = \rho((ST)^3)\rho(T)^{-3} =
1$, so $\rho(S) = 1$ and $\rho$ is trivial (\lean{modularRepOf_eq_one_of_holomorphic}). In
\lean{char_SL2Z_invariant_of_analyticRank}, \lean{analyticSlashAmbient} is Chapter~\ref{ch:L15}'s finite-dimensional space of
analytic blocks containing the trace block and its $S$-slash; its finite-dimensionality is a theorem there, and the bound
$\le1$ is the honest hypothesis (``$\#\Phi = 1$''). \lean{char_eq_jInvariant_add_const} inherits the vacuous-looking hypothesis
of \lean{eq_jInvariant_add_const}: it assumes the conclusion's analytic content.
\end{dictionary}
\noindent\emph{Proof.} As in the dictionary, plus Theorem~\ref{l14:thm:rhoS}'s machinery for the second statement (independence
is free for one module once the character is nonzero, and nonvanishing is free from the growth datum: on the imaginary axis the character is a convergent series with
nonnegative coefficients and positive leading coefficient, \lean{traceFun_vacuum_ne_zero}, Chapter~\ref{ch:L15}).
\begin{physics}
A $c=24$ theory with a single sector (holomorphic) has a modular-invariant partition function; with $V_1 = 0$ it is $J$. These
theorems get the first half from the torus Ward identities alone; the second half (Hauptmodul) is assumed.
\end{physics}
\begin{status}
Conditional on the visible hypotheses; standard axioms, no \lean{VOA.pbw} (all three).
\end{status}
\begin{warning}
\textbf{Not instantiated at $\Vnat$; formerly stale docstrings (documentation only).} No theorem of this section is applied to the library's $\Vnat$.
The library's $\chi_{\Vnat} = J$ (Chapter~\ref{ch:05}) is proved by computing both sides (theta series, $\dim M_{12} = 2$, the
level-two identity) and uses none of the modularity machinery of this chapter. When this chapter was first written, the docstrings of
\leanpath{MoonshineCorollary.lean}, \leanpath{MoonshineFromBlocks.lean}, \leanpath{MoonshineHauptmodul.lean} said that $\Vnat$ did
not exist in the repository (``Phase-3 project''), although the library already constructed $\Vnat$, and the
``Residual'' of \leanpath{MoonshineCorollary.lean} listed $[q^2]\mathtt{jShift} = 196884$ as open, although it is \lean{coeff_two_jShift}
(Chapter~\ref{ch:05}). At a pre-release build commit these docstrings have been corrected: they name \lean{Moonshine.moonshine} as the
construction of $\Vnat$ and record \lean{coeff_two_jShift} as proved.
\end{warning}

%======================================================================
\section{Recorded traps, dead ends and defects}\label{l14:sec:traps}
%======================================================================

Severities are graded on a four-step scale: serious (a statement that is wrong or vacuous), moderate (weaker or stronger than advertised, or a redundant
hypothesis), open gap (a conditional result whose hypothesis is not discharged, or a missing piece), and documentation
only (docstrings and notes).
\begin{enumerate}
\item \lean{levelOne_coe_mem_E46} is generation only; homogeneity and independence are elsewhere (vendored
  \lean{levelOne_eq_polynomial}); the two developments are disconnected (documentation only; Theorem~\ref{l14:thm:levelone}).
\item \lean{E46HasWeightedPolynomialReps} was documented as missing but follows from the vendored \lean{span_eisBasisFun}
  (Theorem~\ref{l14:thm:certreps}). \textbf{Closed}: now the library theorem \lean{e46HasWeightedPolynomialReps}.
\item \lean{E46SerreClosed} / \lean{E46SerreMonomialClosed} (Ramanujan identities, Serre closure) were open named predicates.
  \textbf{Closed}: \lean{e46SerreClosed}, \lean{e46SerreMonomialClosed} (Section~\ref{l14:sec:ramanujan}).
\item \lean{ZhuRecursionEisenstein} uses the superseded $E_{2j}$ normalisation and would be false; no consumer; docstring of
  \leanpath{EisensteinFormal.lean} stale (documentation only).
\item \lean{VacuumGradeZero} was a redundant hypothesis of the cusp theorems (it holds for every VOA,
  \lean{vacuumGradeZero_of_vacuum_grade}) (moderate; Section~\ref{l14:sec:cusp}); fixed in the library since.
\item \lean{CompleteRationalFamily} is the identity; \lean{Conjecture_HalfWeight}/\lean{Conjecture_ZGradedSuper} have a free
  proposition slot; \lean{SuperCharacterSystem}, \lean{SuperRealisesLawOn} are not tied to $W$ (moderate; checked in Lean,
  Proposition~\ref{l14:prop:certtheta}).
\item \lean{eq_jInvariant_add_const} and \lean{char_eq_jInvariant_add_const} assume the Hauptmodul content (moderate; open gap).
\item \lean{JacobiTransformation} (Gannon Thm 5.3.9), \lean{MoonshineCharacter}: statements only (open gap).
\item \lean{BuhlSpanning} is stronger than Buhl's theorem (all modules, finite bottom); its ordered monomials are weaker than
  GN/Buhl (no strict decrease) (moderate, documented by the library).
\item Docstrings: ``$p(n)\le2^n$ gives subexponential'' (false) in \leanpath{ConvergenceBridge.lean}; ``growth datum not built''
  (\leanpath{ConvergenceBridge.lean}, \leanpath{GrowthBound.lean}); ``not connected to a VOA'' (\leanpath{Majorant.lean});
  \lean{tendsto_indicial_quotient} misnamed (\leanpath{IndicialRecursion.lean}); (C2) unbuilt and (C3) prerequisite
  contradiction (\leanpath{Schwarzian.lean} vs \leanpath{ChangeCoord.lean}); \lean{ThetaGroupProper} ``not proved''
  (\leanpath{SlashSubgroup.lean}); \lean{ThetaHalfAutFactorExists} ``OPEN'' (\leanpath{GammaThetaStatement.lean}); $\Vnat$ ``does
  not exist'' and $[q^2]j$ ``open'' (\leanpath{MoonshineCorollary.lean} and siblings) (all documentation only).
\item Dead end, recorded by the library: the purely formal $\C[[q]]$ route cannot reach a point of $\lfourteenH$
  (Theorem~\ref{l14:thm:noalghom}). Dead end: a fixed positive lower bound on the indicial divisor does not close the majorant
  induction; $p_n\to\infty$ is needed (Theorem~\ref{l14:thm:majorant}).
\end{enumerate}

%======================================================================
\section{Referee's checklist}\label{l14:sec:checklist}
%======================================================================

\begingroup\small
\begin{longtable}{p{4.3cm}p{4.4cm}p{3.3cm}p{1.5cm}}
\toprule
Lean declaration & intended meaning & hypotheses; status & check \\
\midrule
\endhead
\lean{slashC_deriv_eq}, \lean{slashC_G2C} & slash vs.\ derivative; $G_2$ quasimodular & differentiability; unconditional & Thms.~\ref{l14:thm:slashderiv}, \ref{l14:thm:G2slash}\\
\lean{zhuG}, \lean{zhuGC} & Zhu's $\mathcal G_{2j} = -\frac{B_{2j}}{(2j)!}+\frac2{(2j-1)!}\sum\sigma q^n$ & definitions & Defs.~\ref{l14:def:zhuG}, \ref{l14:def:zhuGC}\\
\lean{zhuGC_eq_mul_E}, \lean{slashC_zhuGC} & $\mathcal G_{2j} = g_j(0)E_{2j}$, modular & $j\ge2$; unconditional & Thm.~\ref{l14:thm:zhuGCE}\\
\lean{zhuGC_one}, \lean{zhuG_one_eq_neg_E2_div_twelve} & $\mathcal G_2 = -E_2/12$ & unconditional & Thm.~\ref{l14:thm:zhuGCone}\\
\lean{tendsto_zhuGC_atImInfty} & cusp value $g_j(0)$ & $j\ge2$; unconditional & Thm.~\ref{l14:thm:cusplim}\\
\lean{levelOne_coe_mem_E46} & $M_k\subseteq\C\langle E_4,E_6\rangle$ & none; \textbf{generation only} & Thm.~\ref{l14:thm:levelone}\\
\lean{E46SerreClosed}, \lean{e46SerreClosed} & $\partial_kM_k\subseteq M_{k+2}$ & \textbf{proved} (\lean{serreForm}); unconditional & §\ref{l14:sec:ramanujan}\\
\lean{evInfty}, \lean{evAt}, \lean{ker_isMaximal} & cusp and interior points of $\lfourteenE$ & unconditional & Def.~\ref{l14:def:evInfty}\\
\lean{no_algHom_powerSeries_X_ne_zero} & no fibre of $\C[[q]]$ at $q_0\neq0$ (dead end) & unconditional & Thm.~\ref{l14:thm:noalghom}\\
\lean{Subexponential} and closure lemmas & radius of convergence $\ge1$ & unconditional & Def.~\ref{l14:def:subexp}, Thm.~\ref{l14:thm:subexpclosure}\\
\lean{qEvalAt}, \lean{qEvalC}, \lean{qSummableAt_of_subexponential} & evaluation of formal series & growth; unconditional & Def.~\ref{l14:def:qeval}, Thm.~\ref{l14:thm:qsummable}\\
\lean{differentiableOn_qEvalC}, \lean{deriv_qEvalC} & holomorphy; $\frac d{d\tau} = 2\pi iD$ & growth & Thm.~\ref{l14:thm:holo}\\
\lean{differentiableOn_char_adjoint} & character of a $C_2$ VOA holomorphic & $C_2$, $\mathbf 1\neq0$ & Thm.~\ref{l14:thm:holo}\\
\lean{qEvalC_qConv}, \lean{qEvalC_finsum}, \lean{qEvalAt_shift} & convolution $\to$ product; linearity; rebasing & support, summability & Thms.~\ref{l14:thm:qconv}, \ref{l14:thm:qevallin}\\
\lean{leading_coeff_eq_zero_of_qEvalAt_imAxis_eq_zero} & leading coefficient detected at $\tau=it$ & growth & Thm.~\ref{l14:thm:qevallin}\\
\lean{zhu_modularity_T_pow_analytic} & law on $\langle T\rangle$ & none (termwise) & Thm.~\ref{l14:thm:Tanalytic}\\
\lean{psSMul_mul}, \lean{instModulePowerSeriesQSeriesAt} & $\mathcal Q_h$ is a $\C[[q]]$-module & support on $h+\N$ & Def.~\ref{l14:def:QSeriesAt}\\
\lean{qDeriv_psSMul}, \lean{theta_mem_subexpPS} & Leibniz rule; $\theta$-stability of $\mathcal S$ & unconditional & Thm.~\ref{l14:thm:productrule}\\
\lean{qExpansion_zhuGC}, \lean{mk_zhuG_mem_GSeries}, \lean{isNoetherianRing_GSeries} & formal = analytic Eisenstein; Noetherian $\C[\mathcal G_2,\mathcal G_4,\mathcal G_6]$ & unconditional & Thm.~\ref{l14:thm:qexpzhuG}, Def.~\ref{l14:def:GSeries}\\
\lean{serreDerivC}, \lean{serreDerivC_slashC} & Serre derivative is equivariant & differentiability; unconditional & Def.~\ref{l14:def:serre}, Thm.~\ref{l14:thm:serreslash}\\
\lean{pcountLe_subexponential}, \lean{subexponentialCoeffs_of_pbw}, \lean{subexponentialCoeffs_of_buhl} & free-boson bound on degeneracies & spanning data / \lean{BuhlSpanning} & Thms.~\ref{l14:thm:pcount}, \ref{l14:thm:pbwgrowth}\\
\lean{subexponential_of_majorant_recursion} & majorant theorem & $p_n\to\infty$ & Thm.~\ref{l14:thm:majorant}\\
\lean{subexponential_of_indicial_recursion}, \lean{tendsto_indicialDivisor_atTop} & indicial recursion $\Rightarrow$ growth; divisor grows & monic, $s\ge1$ & Thm.~\ref{l14:thm:indicial}\\
\lean{onePointSubexp_of_formalODE}, \lean{onePointSubexp_of_GSeriesODE} & formal MLDE $\Rightarrow$ growth datum & ODE hypothesis & Thm.~\ref{l14:thm:formalODE}\\
\lean{subexponential_of_system} & first-order systems & subexponential coefficients & Thm.~\ref{l14:thm:system}\\
\lean{schwarzian_expMinusOne}, \lean{exp_subst_log}, (I1)--(I3) & $\{e^z-1;z\} = -\frac12$; weight identities & unconditional & Thms.~\ref{l14:thm:schwarzian}, \ref{l14:thm:explog}\\
\lean{Lsq_virasoro} & torus Virasoro, same $c$ & none & Thm.~\ref{l14:thm:Lsq}\\
\lean{changeCoord} & $T_f = \exp\sum a_kL_k$ & definition; law (C3) open & Def.~\ref{l14:def:changeCoord}\\
\lean{negTwoFamily_cusp} & DLM Lemma 9.2(c) & $k\ge1$; unconditional & Thm.~\ref{l14:thm:negtwo}\\
\lean{blockRelSpanAtCusp_eq_cuspRelSpan}, \lean{finrank_quot_cuspRelSpan_eq_of_algEquiv_pi_matrix} & cusp fibre $= A/[A,A]$; dimension $=$ \#Wedderburn factors & none (former \lean{VacuumGradeZero} removed); decomposition & Thm.~\ref{l14:thm:cuspfibre}\\
\lean{traceSpanning_of_slashSpannedAtS}, \lean{zhu_modularity_of_slashSpannedAtS} & $S$-spanning $\Rightarrow$ Zhu's law & weights, growth, $S$-spanning, distinct $T$-eigenvalues, nonvanishing & Thm.~\ref{l14:thm:rhoS}\\
\lean{tEigen_injective_iff} & distinct $T$-eigenvalues $\iff$ weights distinct mod $\Z$ & unconditional & Thm.~\ref{l14:thm:teigen}\\
\lean{JacobiTransformation} & Jacobi-form law & \textbf{statement only} & Def.~\ref{l14:def:jacobi}\\
\lean{thetaGroup}, \lean{thetaGroup_le_submonoidClosure}, \lean{realisesLawOn_thetaGroup_of_S}, \lean{thetaGroup_index} & $\lfourteenGt$; monoid generation; bosonic case & $S$-law (for the law) & Thm.~\ref{l14:thm:thetamonoid}\\
\lean{HalfAutFactor}, \lean{SuperCharacterSystem}, \lean{Conjecture_HalfWeight} & half-integral law; $\lfourteenGt$ conjecture & \textbf{templates}; free \lean{hFam} & Def.~\ref{l14:def:halfaut}\\
\lean{modularLawOn_thetaGroup_of_zhuModularityAnalytic}, \lean{characterModularOn_thetaGroup_of_C2Cofinite} & bosonic consistency & Zhu's law / $C_2$, semisimple, \lean{Pr} & Thm.~\ref{l14:thm:thetaboson}\\
\lean{jInvariant_smul}, \lean{eq_jInvariant_add_const} & $j$ invariant; Hauptmodul step & \textbf{assumes} $\chi - j\in M_0$ & Def.~\ref{l14:def:jinv}\\
\lean{char_SL2Z_invariant_of_spanning}, \lean{char_SL2Z_invariant_of_analyticRank} & $c=24$, $h=0$ character invariant & spanning data, or $C_2$, growth, rank $\le1$ & Thm.~\ref{l14:thm:moonsl2}\\
\lean{MoonshineCharacter} & $\mathrm{ch}\,V = qJ$ & \textbf{statement only}; not instantiated & Def.~\ref{l14:def:moonchar}\\
\bottomrule
\end{longtable}
\endgroup

All \texttt{\#print axioms} results quoted in this chapter were obtained at an earlier commit (not re-run for this edition) by running a scratch file
importing the relevant modules with \texttt{lake env lean -DmaxSynthPendingDepth=3} (never a bare \texttt{lake env lean}); 69
headline declarations were checked, and every one reports exactly
\begin{plaincode}
[propext, Classical.choice, Quot.sound]
\end{plaincode}
No headline of this chapter uses \lean{VOA.pbw}. No full \texttt{lake build} was run for this chapter. The
separate Lean checks of this chapter (not part of this release) compiled with no errors.

\paragraph{Release audit 2026-10-27: RELEASE\_AUDIT\_ZHU.md.} The release audit
(\texttt{Notes/\allowbreak Textbook/\allowbreak referee/\allowbreak RELEASE\_AUDIT\_ZHU.md}) found this chapter's two headlines
that it examined faithful as stated: \lean{levelOne_coe_mem_E46} (generation of the level-one modular forms by $E_4$, $E_6$
only, as Section~\ref{l14:sec:levelone} says) and \lean{subexponential_of_majorant_recursion} (an abstract lemma about
nonnegative real sequences). Checked again later, at a commit before this edition (not re-run for this edition), with \texttt{lake env lean -DmaxSynthPendingDepth=3}:
\lean{VOA.qSummableAt_of_subexponential}, standard three axioms. No other entry was re-run.

\paragraph{Second release revision 2026-10-28.} Only the description of the consumer theorem of Chapter~\ref{ch:L15} changed
(finding F2 of \texttt{RELEASE\_AUDIT\_NEW\_DECLS.md}: ``rational $\Rightarrow A(V)$ semisimple'' is proved in the library, so
Zhu's theorem holds for $C_2$-cofinite \lean{AdmissiblyRational} $V$). No statement of this chapter changed and no
\texttt{\#print axioms} of this chapter was re-run; the axiom checks of the new Chapter~\ref{ch:L15} theorems are recorded
there (standard three axioms, at an earlier commit, not re-run for this edition). No \texttt{lake build} was run.

%======================================================================
\section{Notes and references}\label{l14:sec:notes}
%======================================================================

\textbf{Modular forms.} J.-P.~Serre, \emph{A Course in Arithmetic}, Graduate Texts in Mathematics 7, Springer (1973), Ch.~VII
(the structure of $M_*$, Eisenstein series, $\Delta$, $j$). D.~Zagier, ``Elliptic modular forms and their applications'', in
\emph{The 1-2-3 of Modular Forms}, Universitext, Springer (2008), 1--103 (quasimodular $E_2$, eq.~(17); the Serre derivative and
Ramanujan's identities, §5). M.~Kaneko and D.~Zagier, ``Supersingular $j$-invariants, hypergeometric series, and Atkin's
orthogonal polynomials'', in \emph{Computational Perspectives on Number Theory}, AMS/IP Stud.\ Adv.\ Math.\ 7 (1998), 97--126
(the Kaneko--Zagier equation, a second-order MLDE). T.~Gannon, \emph{Moonshine beyond the Monster}, Cambridge University Press
(2006), §§2.3, 5.3 (the conventions transcribed by the library, eqs.~(5.3.15)--(5.3.21)).

\textbf{Modular linear differential equations.} S.~D.~Mathur, S.~Mukhi and A.~Sen, ``On the classification of rational conformal
field theories'', \emph{Phys.\ Lett.\ B} 213 (1988), 303--308; ``Reconstruction of conformal field theories from modular
geometry on the torus'', \emph{Nucl.\ Phys.\ B} 318 (1989), 483--540. G.~Mason, ``Vector-valued modular forms and linear
differential operators'', \emph{Int.\ J.\ Number Theory} 3 (2007), 377--390.

\textbf{Zhu's theorem and its analytic toolkit.} Y.~Zhu, ``Modular invariance of characters of vertex operator algebras'',
\emph{J.\ Amer.\ Math.\ Soc.} 9 (1996), 237--302 (square-bracket modes §4.2; recursion Prop.~4.3.2; ODE §4.4; Thm.~5.3.2). C.~Dong,
H.~Li and G.~Mason, ``Modular-invariance of trace functions in orbifold theory and generalized Moonshine'', \emph{Comm.\ Math.\
Phys.} 214 (2000), 1--56, arXiv:q-alg/9703016 (Lemma~9.2, the cusp degeneration; Thm.~5.4, the Serre derivative). G.~Buhl, ``A
spanning set for VOA modules'', \emph{J.\ Algebra} 254 (2002), 125--151, arXiv:math/0111296. M.~Gaberdiel and A.~Neitzke,
``Rationality, quasirationality and finite W-algebras'', \emph{Comm.\ Math.\ Phys.} 238 (2003), 305--331, arXiv:hep-th/0009235.
Y.-Z.~Huang, \emph{Two-dimensional Conformal Geometry and Vertex Operator Algebras}, Birkh\"auser (1997) (change of coordinates).

\textbf{Jacobi forms, theta group, super VOAs.} M.~Miyamoto, ``A modular invariance on the theta functions defined on vertex
operator algebras'', \emph{Duke Math.\ J.} 101 (2000), 221--236. M.~Krauel and G.~Mason, ``Vertex operator algebras and weak Jacobi
forms'', \emph{Int.\ J.\ Math.} 23 (2012), 1250024. G.~H\"ohn, \emph{Selbstduale Vertexoperatorsuperalgebren und das
Babymonster}, PhD thesis, Bonn (1995), arXiv:0706.0236. C.~Dong and Z.~Zhao, ``Modularity in orbifold theory for vertex operator
superalgebras'', \emph{Comm.\ Math.\ Phys.} 260 (2005), 227--256, arXiv:math/0411524. J.~van Ekeren, ``Modular invariance for twisted
modules over a vertex operator superalgebra'', \emph{Comm.\ Math.\ Phys.} 322 (2013), 333--371, arXiv:1111.0682.

\textbf{Physics.} J.~L.~Cardy, ``Operator content of two-dimensional conformally invariant theories'', \emph{Nucl.\ Phys.\ B} 270
(1986), 186--204 (the growth of degeneracies from modular invariance). E.~Verlinde, ``Fusion rules and modular transformations in
2D conformal field theory'', \emph{Nucl.\ Phys.\ B} 300 (1988), 360--376. P.~Di~Francesco, P.~Mathieu and D.~S\'en\'echal,
\emph{Conformal Field Theory}, Springer (1997), Ch.~10 (modular invariance on the torus).

\endgroup
